\RequirePackage[thinlines]{easybmat} %-- muss aufgrund eines Bugs vor etex und tikz geladen werden
\documentclass[a4paper, twoside, headsepline, index=totoc,toc=listof, fontsize=10pt, cleardoublepage=empty, headinclude, DIV=13, BCOR=13mm, titlepage,]{scrartcl}
\usepackage{scrtime} % KOMA, Uhrzeit ermoeglicht
\usepackage{scrpage2} % wie fancyhdr, nur optimiert auf KOMA-Skript, leicht andere Syntax
\usepackage{etoolbox}
\usepackage{letltxmacro}

% Compiler abhaengige IFs
\usepackage{ifxetex,ifluatex}
\newif\ifxetexorluatex
\ifxetex
  \xetexorluatextrue
\else
  \ifluatex
    \xetexorluatextrue
  \else
    \xetexorluatexfalse
  \fi
\fi

%--Farbdefinitionen
%-- muss vor tikz geladen werden
\usepackage[usenames, table, x11names]{xcolor}
\definecolor{dark_gray}{gray}{0.45}
\definecolor{light_gray}{gray}{0.6}

%--Zum Zeichnen
%-- muss vor polyglossia bzw. babel geladen werden
\usepackage{tikz}
\usepackage{tikz-cd}
\usetikzlibrary{external}
\tikzset{>=latex}
\usetikzlibrary{shapes,arrows,intersections}
\usetikzlibrary{calc,3d}
\usetikzlibrary{decorations.pathreplacing,decorations.markings}
\usetikzlibrary{angles}

%-- Konfiguration von tikzexternalize
\tikzexternalize[prefix=tikz/, up to date check=diff]

\ifxetexorluatex
\tikzset{external/system call={xelatex \tikzexternalcheckshellescape %-- verwende LuaLaTeX, wegen dynamischer Speicherallokation
    -halt-on-error -interaction=batchmode -jobname "\image" "\texsource"}}
\else
\tikzset{external/system call={pdflatex \tikzexternalcheckshellescape 
    -halt-on-error -interaction=batchmode -jobname "\image" "\texsource"}}
\fi

%-- tikzexternalize fuer tikzcd deaktivieren, da inkompatibel
\AtBeginEnvironment{tikzcd}{\tikzexternaldisable}
\AtEndEnvironment{tikzcd}{\tikzexternalenable}

%-- um Inkompatibilitaeten von quotes und polyglossia bzw. babel zu vermeiden
\tikzset{
  every picture/.append style={
    execute at begin picture={\shorthandoff{"}},
    execute at end picture={\shorthandon{"}}
  }
}
\usetikzlibrary{quotes}
\usepackage{pgfplots}
\usepgfplotslibrary{colormaps}
\pgfplotsset{compat=1.10}



%-- Mathe
\usepackage{mathtools} % beinhaltet amsmath
\mathtoolsset{showonlyrefs, centercolon}
\newtagform{brackets}[\textbf]{[}{]}
\usetagform{brackets}
\usepackage{wasysym}
\usepackage{amssymb} %zusätzliche Symbole
\usepackage{latexsym} %zusätzliche Symbole
\usepackage{stmaryrd} %für Blitz
\usepackage{nicefrac} %schräge Brüche
\usepackage{cancel} %Befehle zum Durchstreichen
\usepackage{mathdots}
% \usepackage{bm}
\DeclareSymbolFont{bbold}{U}{bbold}{m}{n}
\DeclareSymbolFontAlphabet{\mathbbold}{bbold}
\newcommand{\mathds}[1]{\mathbb{#1}} %Um Kompatibilitaet mit frueheren Benutzung von dsfont herzustellen
\newcommand{\ind}{\mathbbold{1}} % charakteristische-Funktion-Eins
\def\mathul#1#2{\color{#1}\underline{{\color{black}#2}}\color{black}} %farbiges Untersteichen im Mathe-Modus

%-- Alles was mit Schrift und XeteX zu tun hat
\ifxetexorluatex
    % XeLaTeX or LuaTeX
	\ifxetex
		\usepackage{mathspec} %beinhaltet fontspec 
	\else
		\usepackage[no-math]{fontspec}
	\fi
    \usepackage{polyglossia} %babel-ersatz
    \setmainlanguage[spelling=new,babelshorthands=false]{german}
	\newcommand\glqq{"}
	\newcommand\grqq{"}
	\defaultfontfeatures{Mapping=tex-text, WordSpace={1.4}} %
    \setmainfont[Ligatures=Common, BoldFont={* Bold}, ItalicFont={* Light Italic},SmallCapsFont={Linux Libertine O}, SmallCapsFeatures={Letters=SmallCaps}]{Source Sans Pro}
	\setsansfont[Scale=MatchLowercase,Ligatures=Common, BoldFont={* Medium}]{Ubuntu}
	\ifxetex
		\setallmonofonts[Scale=MatchLowercase, ItalicFont={*}]{Consolas} 
	\else
		\setmonofont[Scale=MatchLowercase, ItalicFont={*}]{Consolas}
	\fi
	\usepackage{xltxtra}
	\usepackage{fontawesome}
	\usepackage{microtype}
\else
    % default: pdfLaTeX
    \usepackage[ngerman]{babel}
    \usepackage[T1]{fontenc}
    \usepackage[utf8]{inputenc}
	\usepackage{textcomp} %verhindert ein paar Fehler bei den Fonts
    \usepackage[babel=true, tracking=true]{microtype}
	\usepackage{lmodern}
	\renewcommand{\familydefault}{\sfdefault} 
\fi


%--Mixed
\usepackage[neverdecrease]{paralist}
\usepackage[german=quotes]{csquotes}
\usepackage{makeidx}
\usepackage{booktabs}
\usepackage{wrapfig}
\usepackage{float}
\usepackage[margin=10pt, font=small, labelfont={sf, bf}, format=plain, indention=1em]{caption}
\captionsetup[wrapfigure]{name=Abb. }
\usepackage{stackrel}
\usepackage{ifthen}
\usepackage{multicol}
\usepackage{qrcode}
\flushbottom

%--Hyphenation
\hyphenation{di-men-si-o-na-le}


%--Unterstreichung
\usepackage[normalem]{ulem}
\setlength{\ULdepth}{1.8pt}

%--Indexverarbeitung
\newcommand{\bet}[1]{\textbf{#1}}
\newcommand{\Index}[1]{\textbf{#1}\index{#1}}
\makeindex
\setindexpreamble{{\noindent \itshape Die \emph{Seitenzahlen} sind mit Hyperlinks zu den entsprechenden Seiten versehen, also anklickbar} \ifxetexorluatex \faHandUp\fi \par \bigskip}
\renewcommand{\indexpagestyle}{scrheadings}


%--Marginnote & todonotes
\usepackage{marginnote}
\renewcommand*{\marginfont}{\color{dark_gray} \itshape \footnotesize }
\usepackage[textsize=small]{todonotes}

%--Todonotes mit tikz/externalize kompatibel machen
\LetLtxMacro{\oldtodo}{\todo}
\renewcommand{\todo}[2][]{\tikzexternaldisable\oldtodo[#1]{#2}\tikzexternalenable}
\LetLtxMacro{\oldmissingfigure}{\missingfigure}
\renewcommand{\missingfigure}[2][]{\tikzexternaldisable\oldmissingfigure[{#1}]{#2}\tikzexternalenable}

%--Konfiguration von Hyperref pdfstartview=FitH, 
\usepackage[hidelinks, pdfpagelabels,  bookmarksopen=true, bookmarksnumbered=true, linkcolor=black, urlcolor=SkyBlue2, plainpages=false,pagebackref, citecolor=black, hypertexnames=true, pdfauthor={Jannes Bantje}, pdfborderstyle={/S/U}, linkbordercolor=SkyBlue2, colorlinks=false]{hyperref}

% \let\orighref\href
\ifxetexorluatex
\newcommand{\hrefsym}[2]{\href{#1}{\texttt{#2\,\faExternalLink}}}
\newcommand{\hrefsymmail}[2]{\href{#1}{\texttt{\faEnvelope\,#2}}}
\renewcommand{\url}[1]{\hrefsym{#1}{\nolinkurl{#1}}}
\fi
\providecommand{\hrefsym}{\href}
% \usepackage{cleveref}
% \crefformat{footnote}{#2\footnotemark[#1]#3}

%--Römische Zahlen
\newcommand{\RM}[1]{\MakeUppercase{\romannumeral #1{}}}

%--Punkte (nach hyperref)
\usepackage{ellipsis}

%%--Abkürzungen etc.
\newcommand{\light}{\text{\Large $\lightning$}}

%-- Definitionen von weiteren Mathe-Befehlen
\DeclareMathOperator{\re}{Re} %Realteil
\DeclareMathOperator{\im}{Im} %Imaginaerteil
\DeclareMathOperator{\id}{id} %identische Abbildung
\DeclareMathOperator{\Sp}{Sp} %Spur
\DeclareMathOperator{\sgn}{sgn} %Signum
\DeclareMathOperator{\alt}{Alt} %Alternierende n-linearFormen
\DeclareMathOperator{\End}{End} %Endomorphismen
\DeclareMathOperator{\Vol}{Vol} %Volumen
\DeclareMathOperator{\dom}{dom} %Domain
\DeclareMathOperator{\Hom}{Hom} %Homomorphismen
\DeclareMathOperator{\bild}{Bild} %Bild
\DeclareMathOperator{\Kern}{Kern }
\DeclareMathOperator{\rg}{rg} %Rang
\DeclareMathOperator{\Rg}{Rg} %Rang
\DeclareMathOperator{\diam}{diam} %Durchmesser
\DeclareMathOperator{\dist}{dist} %Distanz
\DeclareMathOperator{\grad}{grad} %Gradient
\DeclareMathOperator{\dive}{div} %Gradient
\DeclareMathOperator{\rot}{rot} %Rotation
\DeclareMathOperator{\hess}{Hess} %Hesse-Matrix
\DeclareMathOperator{\koker}{Koker} %Kokern
\DeclareMathOperator{\aut}{Aut} %Automorphismen
\DeclareMathOperator{\ord}{ord} %Ordnung
\DeclareMathOperator{\ggT}{ggT} %ggT
\DeclareMathOperator{\kgV}{kgV} %kgV
\DeclareMathOperator{\Gr}{Gr} %Gerade
\DeclareMathOperator{\Kr}{Kr} %Kreis
\DeclareMathOperator{\Char}{char} %Charakteristik
\DeclareMathOperator{\Aut}{Aut} %Automorphismen
% \DeclareMathOperator{\D}{D} %formale Ableitung
\newcommand{\D}{\ensuremath{\mathrm{D}\mkern-1.5mu}}
\newcommand{\mathd}{\ensuremath{\mathrm{d}\mkern-1.0mu}}
\newcommand{\Tmap}{\ensuremath{\mathrm{T}\mkern-1.0mu}}
\newcommand{\opL}{\ensuremath{\mathrm{L}\mkern-0.6mu}}
\DeclareMathOperator{\cov}{cov} %Kovarianz
\DeclareMathOperator{\Gal}{Gal} %Galois
\DeclareMathOperator{\supp}{supp} %Träger
\DeclareMathOperator{\Sym}{Sym} %Symmetrische Gruppe
\DeclareMathOperator{\Zyl}{Zyl} %Zylinder
\DeclareMathOperator{\Mod}{Mod} %Moduln
\DeclareMathOperator{\EPK}{EPK} %Einpunktkompaktifizierung
\DeclareMathOperator{\conj}{conj} %Konjugation
\DeclareMathOperator{\ann}{ann} %Annulator
\DeclareMathOperator{\tor}{tor} %Torsionsmodul
\DeclareMathOperator{\Int}{Int} %Interior/Inneres
\DeclareMathOperator{\Br}{Br} %Brauerergruppe
\DeclareMathOperator{\GL}{GL} %allgemeine lineare Gruppe
\DeclareMathOperator{\SL}{SL} %spezielle lineare Gruppe
\DeclareMathOperator{\SU}{SU} %spezielle unitäre Gruppe
\DeclareMathOperator{\Tr}{Tr} %Spur/Trace
\DeclareMathOperator{\conv}{conv} %Konvexe Teilmenge
% \DeclareMathOperator{\tr}{tr}
\newcommand{\tr}{\mathrm{tr}}
\newcommand{\w}{\mathrm{w}}
\newcommand{\weakT}[1]{\ensuremath{\mathcal{T}_{#1}^{\mkern+1.0mu\text{\raisebox{0.4ex}{$\mathrm{w}$}}}}}
\newcommand{\weakTstar}[1]{\ensuremath{\mathcal{T}_{#1}^{\mkern+1.0mu\text{\raisebox{0.4ex}{$\mathrm{w}$}}^*}}}
\newcommand{\rand}[1]{\ensuremath{\partial^{\scriptscriptstyle #1}}}


% -- Kategorien
\DeclareMathOperator{\Ob}{Ob} %Objekte
\DeclareMathOperator{\Mor}{Mor} %Morphismen

% \DeclareMathOperator{\TOP}{TOP} %Topologische Räume
\newcommand{\TOP}{\textsc{Top}}
% \DeclareMathOperator{\HTOP}{HTOP} %Topologische Räume + Homotopieäquivalenz
\newcommand{\HTOP}{\textsc{HTop}}
% \DeclareMathOperator{\VR}{VR} %Vektorräume
\newcommand{\VR}{\textsc{VR}}
% \DeclareMathOperator{\MOD}{MOD} %Moduln
\newcommand{\MOD}{\textsc{Mod}}
% \DeclareMathOperator{\MENGEN}{MENGEN} %Mengen
\newcommand{\MENGEN}{\textsc{Mengen}}
% \DeclareMathOperator{\MAN}{MAN} %Mannigfaltigkeiten
\newcommand{\MAN}{\textsc{Man}}
% \DeclareMathOperator{\GRUPPEN}{GRUPPEN} %Gruppen
\newcommand{\GRUPPEN}{\textsc{Gruppen}}
% \DeclareMathOperator{\ABELGRUPPEN}{ABEL.GRUPPEN} %abelsche Gruppen
\newcommand{\ABELGRUPPEN}{\textsc{Abel.Gruppen}}
% \DeclareMathOperator{\KAT}{KAT} %Kategorien
\newcommand{\KAT}{\textsc{Kat}}
% \DeclareMathOperator{\FUN}{FUN} %Funktoren
\newcommand{\FUN}{\textsc{Fun}}


\newcommand{\op}{\mathrm{op}}

\newcommand{\Ker}{\ensuremath{\Kern \,}}
\newcommand{\zirkel}{\rotatebox[origin = c]{-90}{$\varangle$}}
\newcommand{\wirkung}{\!\curvearrowright\!}

\newcommand{\oE}{\OE}

%--Beweisende
\newcommand{\bewende}{\ifmmode \tag*{$\square$} \else \hfill $\square$ \fi}

%--Volumen
\newcommand{\vol}[1]{\ensuremath{\Vol_{#1}}}

%--Skalarprodukt
\DeclarePairedDelimiterX\sprod[2]{\langle}{\rangle}{#1\,\delimsize\vert\,#2}

%--Betrag, Gaußklammer
\DeclarePairedDelimiter{\abs}{\lvert}{\rvert}
\DeclarePairedDelimiter{\ceil}{\lceil}{\rceil}

%--Norm
\DeclarePairedDelimiter\doppelstrich{\Vert}{\Vert}
\newcommand{\norm}[2][\relax]{
\ifx#1\relax \ensuremath{\doppelstrich*{#2}}
\else \ensuremath{\doppelstrich*{#2}_{#1}}
\fi}


%--Umklammern
\DeclarePairedDelimiter\enbrace{(}{)}
\DeclarePairedDelimiter\benbrace{[}{]}
\DeclarePairedDelimiter\lenbrace{<}{>}

%--Mengen
\DeclarePairedDelimiterX\mengenA[1]{\lbrace}{\rbrace}{#1}
\DeclarePairedDelimiterX\mengenB[2]{\lbrace}{\rbrace}{#1\, \delimsize\vert \, #2}

\makeatletter
\newcommand{\set}[2][\relax]{
\ifx#1\relax \ensuremath{
\mengenA*{#2}}
\else \ensuremath{%
  \mengenB*{#1}{#2}}
\fi}
\makeatother


%--offen und abgeschlossen
\newcommand{\off}{\! \stackrel[\text{offen}]{}{\subset} \!}
\newcommand{\abg}{\! \stackrel[\text{abg.}]{}{\subset} \!}

%--Differential
\newcommand{\diff}[2]{\ensuremath{\frac{{\partial #1}}{{\partial #2}} }}
\newcommand{\diffd}[2]{\ensuremath{\frac{\mathd #1}{\mathd #2} }}

%--Diagonale Linien für Matrizen
\newcommand\tikzmark[1]{%
  \tikzexternaldisable\tikz[overlay,remember picture,baseline] \node [anchor=base] (#1) {};\tikzexternalenable}

\newcommand\MyLine[3][]{%
  \tikzexternaldisable\begin{tikzpicture}[overlay,remember picture]
    \draw[#1] (#2.south) -- (#3.north);
  \end{tikzpicture}\tikzexternalenable}
  
%--Jordankasten
\newcommand{\Jord}{ \tikzexternaldisable \begin{BMAT}{|ccc|}{|ccc|} %
				0 & & \\ %
				1\tikzmark{z} &  \diagdown & \\ %
				& \tikzmark{u}1 \MyLine[thick]{z}{u} &  0 %
			\end{BMAT} \tikzexternalenable}
			
%--Jordankasten mit Lambda
\newcommand{\JordLa}[1][\empty]{%
	\tikzexternaldisable\ifthenelse{\equal{#1}{\empty}}
		{\begin{BMAT}{|ccc|}{|ccc|} %
				\lambda & & \\ %
				1\tikzmark{z} &  \diagdown & \\ %
				& \tikzmark{u}1 \MyLine[thick]{z}{u} &  \lambda  %
			\end{BMAT}}
		{\begin{BMAT}{|ccc|}{|ccc|} %
				\lambda_{#1}  & & \\ %
				1\tikzmark{z} &  \diagdown & \\ %
				& \tikzmark{u}1 \MyLine[thick]{z}{u} &  \lambda_{#1}  %
			\end{BMAT}}\tikzexternalenable}

%--Inhaltsverzeichnis
\usepackage[tocindentauto]{tocstyle}
\usetocstyle{KOMAlike}			

\newcommand{\fach}{Topologie \RM{1}.}
% \newcommand{\shortFach}{Analysis, Topologie, Geometrie}
\newcommand{\semester}{WiSe 2014/2015}
\newcommand{\homepage}{https://wwwmath.uni-muenster.de/reine/u/topos/lehre/WS2014-2015/Topologie1/Topologie1.html}

\newcommand{\prof}{Prof.\,Dr.\,Arthur Bartels}
\providecommand{\verfasser}{Jannes Bantje}

%--Konfiguration von scrheadings
\setheadsepline{1pt}[\color{light_gray}]
\pagestyle{scrheadings}
%\fancyhf[H, F]{}
\clearscrheadfoot

\providecommand{\shortFach}{\fach}
\lehead{ \includegraphics[height=0.5 cm,keepaspectratio]{../!config/Bilder/Logo_WWU_Muenster_light_gray.pdf}}
\rehead{\rule{0cm}{0.5cm}\footnotesize \sffamily \color{light_gray} \verfasser -- Mitschrift \shortFach}
\lohead{\rule{0cm}{0.5cm} \footnotesize \sffamily \color{light_gray} Stand: \today \; \thistime[:]}
\rohead{\includegraphics[height=0.5 cm,keepaspectratio]{../!config/Bilder/fb10logo_gray.pdf}}


\ofoot[{ \color{dark_gray} \LARGE \sffamily \thepage}]{{ \color{dark_gray} \LARGE \sffamily \thepage}} %hier wir auch der plain Stil bearbeitet!
\automark{section}
\ifoot{ \color{dark_gray} \small \leftmark}

%--Metadaten
\author{\verfasser}
\titlehead{\includegraphics[height=1.5cm, keepaspectratio]{../!config/Bilder/Logo_WWU_Muenster.pdf}%
\hfill \includegraphics[height=1.3cm, keepaspectratio]{../!config/Bilder/fb10logo.pdf}}
\title{Skript \fach}
\subtitle{Mitschrift der Vorlesung  \enquote{\fach} von \prof}

\publishers{
Aktuelle Version verfügbar bei: \bigskip\\
\normalsize
\begin{minipage}[t]{0.45\textwidth}
	\centering{
	\qrcode[height=3.5cm, version=4]{https://github.com/JaMeZ-B/latex-wwu}
	\medskip\\
	\includegraphics[height=0.4cm, keepaspectratio]{../!config/Bilder/github_octo.pdf} 
	\includegraphics[height=0.4cm, keepaspectratio]{../!config/Bilder/GitHub_Logo.pdf} (inklusive Sourcecode)\\
	{\small\url{https://github.com/JaMeZ-B/latex-wwu}}}
\end{minipage}
\hfill
\begin{minipage}[t]{0.45\textwidth}
	\centering{
	\qrcode[height=3.5cm, version=4]{btsync://B6WH2DISQ5QVYIRYIEZSF4ZR2IDVKPN3I?n=latex_share}
	\medskip\\
	\includegraphics[height=0.4cm, keepaspectratio]{../!config/Bilder/bt_sync_logo.pdf} 
	{\large \textbf{Bittorrent} Sync} \\
	\texttt{B6WH2DISQ5QVYIRYIEZSF4ZR2IDVKPN3I}
	}
\end{minipage}
}


\begin{document}
\pagenumbering{Roman}
\maketitle
\begin{abstract}
\section*{Vorwort --- Mitarbeit am Skript}
Dieses Dokument ist eine Mitschrift aus der Vorlesung \enquote{\fach, \semester}, gelesen von \prof. Der Inhalt entspricht weitestgehend dem Tafelanschrieb. Für die
Korrektheit des Inhalts übernehme ich keinerlei Garantie! Für Bemerkungen und Korrekturen -- und seien es nur Rechtschreibfehler -- bin ich sehr dankbar. 
Korrekturen lassen sich prinzipiell auf drei Wegen einreichen: 
\begin{itemize}
	\item Persönliches Ansprechen in der Uni, Mails an \hrefsymmail{mailto:\mail}{\mail} (gerne auch mit annotieren PDFs) 
	\item \emph{Direktes} Mitarbeiten am Skript: Den Quellcode poste ich auf GitHub (siehe oben), also stehen vielfältige Möglichkeiten der Zusammenarbeit zur Verfügung:
	Zum Beispiel durch Kommentare am Code über die Website und die Kombination Fork + Pull Request. Wer sich verdient macht oder ein Skript zu einer Vorlesung, die 
	ich nicht besuche, beisteuern will, dem gewähre ich gerne auch Schreibzugriff.
	
	Beachten sollte man dabei, dass dazu ein Account bei \url{github.com} notwendig ist, der allerdings ohne Angabe von persönlichen Daten angelegt werden kann. 
	Wer bei GitHub (bzw. dem zugrunde liegenden Open-Source-Programm \enquote{\texttt{git}}) -- verständlicherweise -- Hilfe beim Einstieg braucht, dem helfe ich gerne 
	weiter. Es gibt aber auch zahlreiche empfehlenswerte Tutorials im Internet.\footnote{zB. \url{https://try.github.io/levels/1/challenges/1}, ist auf Englisch, aber dafür 
	interaktives LearningByDoing}
	\item \emph{Indirektes} Mitarbeiten: \TeX-Dateien per Mail verschicken. 
	
	Dies ist nur dann sinnvoll, wenn man einen ganzen Abschnitt ändern möchte (zB. einen alternativen Beweis geben), da ich die Änderungen dann per Hand einbauen muss! Ich freue mich aber auch über solche Beiträge!
\end{itemize}
\section*{Vorlesungshomepage}
{\centering 
\begin{minipage}[c][][c]{\textwidth}
	\centering \qrcode[height=3cm]{\homepage} \medskip\\
	\footnotesize \url{\homepage}
\end{minipage}
\par}

\end{abstract}

\setcounter{page}{1}
\tableofcontents
\cleardoubleoddemptypage

\pagenumbering{arabic}
\setcounter{page}{1}

\section{Kategorien, Funktoren und natürliche Transformationen} % (fold)
\label{sec:1}

\subsection[Definition: Kategorie]{Definition} % (fold)
\label{sub:11}
Eine \Index{Kategorie} $\mathcal{C}$ besteht aus:
\begin{enumerate}[i)]
	\item Einer Klasse $\Ob(\mathcal{C})$. Die Elemente von $\Ob(\mathcal{C})$ heißen die \Index{Objekte} von $\mathcal{C}$.
	\item Zu je zwei Objekten $A,B$ aus $\mathcal{C}$ einer Menge $\Mor_{\mathcal{C}}(A,B)$. Die Elemente von $\Mor_\mathcal{C}(A,B)$ heißen \Index{Morphismen} von
	$\mathcal{C}$.
	\item Zu je drei Objekten $A,B,C$ aus $\mathcal{C}$ einer Abbildung: 
	\[
		\Mor_\mathcal{C}(B,C) \times \Mor_\mathcal{C}(A,B) \to \Mor_\mathcal{C}(A,C) \enspace, \quad (f,g) \mapsto f \circ g
	\]
	genannt die \Index{Komposition} in $\mathcal{C}$. 
\end{enumerate}
Dabei müssen folgende Axiome  erfüllt sein:
\begin{enumerate}[(i)]
	\item Die Komposition ist \Index{assoziativ}: \\ Für Objekte $A,B,C,D$ von $\mathcal{C}$ und 
	$f \in \Mor_\mathcal{C}(C,D), g \in \Mor_\mathcal{C}(B,C), h \in \Mor_\mathcal{C}(A,B)$ gilt immer
	\[
		f \circ (g \circ h) = (f \circ g) \circ h
	\]
	\item Die Komposition ist \Index{unital}: Für jedes Objekt $A$ von $\mathcal{C}$ gibt es einen Morphismus $\id_A \in \Mor_\mathcal{C}(A,A)$ so, dass
	\[
		\forall f \in \Mor_\mathcal{C}(B,A) : \id_A \circ f = f  \qquad  \text{und} \qquad  \forall f \in \Mor_\mathcal{C}(A,B) : f \circ \id_A = f
	\]
\end{enumerate}
% subsection 11 (end)

\subsection[Beispiele für Kategorien]{Beispiele} % (fold)
\label{sub:12}
\begin{enumerate}[(1)]
	\item Die Kategorie der Mengen $\MENGEN$: Objekte sind Mengen und Morphismen sind Abbildungen.
	\item Die Kategorie der Gruppen $\GRUPPEN$: Objekte sind Gruppen und Morphismen sind Gruppenhomomorphismen.
	\item Die Kategorie der $K$-Vektorräume $K\text{-}\VR$: Objekte sind $K$-Vektorräume und die Morphismen sind $K$-lineare Abbildungen.
	\item Die Kategorie der $R$-Moduln $R\text{-}\MOD$: Objekte sind $R$-Moduln und Morphismen sind $R$-lineare Abbildungen.
	\item Die Kategorie der $C^\infty$-Mannigfaltigkeiten $C^\infty\text{-}\MAN$: Objekte sind $C^\infty$-Mannigfaltigkeiten und Morphismen sind $C^\infty$-Abbildungen.
	\item Die Kategorie der topologischen Räume $\TOP$: Objekte sind topologische Räume und Morphismen sind stetige Abbildungen.
	\item Die Kategorie der punktierten topologischen Räume $\TOP_\bullet$: Objekte sind punktiert-topologische Räume und Morphismen sind punktiert-stetige Abbildungen.
	\item Die Kategorie $\HTOP$: Objekte sind topologische Räume und Morphismen sind Homotopieklassen von stetigen Abbildungen.
	\item Die Kategorie $\HTOP_\bullet$: Objekte sind punktiert-topologische Räume und Morphismen sind Homotopieklassen von punktiert-stetigen Abbildungen.
	\item Sei $G$ eine Gruppe. Wir erhalten eine Kategorie $\mathcal{C}_G$ mit genau einem Objekt $*$ und $\Mor_{\mathcal{C}_G}(*,*) = G$. Die Komposition wird durch die
	Verknüpfung in der Gruppe festgelegt.
	\item Ist $\mathcal{C}$ eine Kategorie, so ist $\mathcal{C}^{\op}$ eine Kategorie, wobei $\Ob(\mathcal{C}^\op) = \Ob(\mathcal{C})$ und 
	\(
		\Mor_{\mathcal{C}^\op}(A,B) := \Mor_\mathcal{C}(B,A)
	\).
	Die Komposition ist gegeben durch $f \circ^\op g := g \circ f$.
\end{enumerate}
% subsection 12 (end)

\subsection[Bemerkungen zu Kategorien]{Bemerkung} % (fold)
\label{sub:13}
\begin{enumerate}[(i)]
	\item Eine Kategorie heißt \bet{klein}\index{Kategorie!klein}, wenn ihre Objekte eine Menge bilden.
	\item Statt $f \in \Mor_\mathcal{C}(A,B)$ schreiben wir oft $f : A \to B$ oder $A \xrightarrow{f} B$.
\end{enumerate}
% subsection 13 (end)

\subsection[Definition: Funktor]{Definition} % (fold)
\label{sub:14}
Seien $\mathcal{C}$ und $\mathcal{D}$ Kategorien: Ein \Index{Funktor} $F : \mathcal{C} \to \mathcal{D}$ ordnet jedem Objekt $C$ von $\mathcal{C}$ ein Objekt $F(C)$ von 
$\mathcal{D}$ und ordnet jedem Morphismus $f : C \to C'$ in $\mathcal{C}$ einen Morphismus $F(f) : F(C) \to F(C')$ in $\mathcal{D}$ zu. Dabei muss gelten:
\[
	F(f \circ g) = F(f) \circ F(g) \quad \text{und} \quad F(\id_C) = \id_{F(C)}
\] 
% subsection 14 (end)

\subsection[Beispiele für Funktoren]{Beispiele} % (fold)
\label{sub:15}
\begin{enumerate}[(1)]
	\item Es gibt offensichtliche \enquote{Vergiss}-Funktoren:
	\begin{align*}
		K\text{-}\VR \enspace\longrightarrow & \enspace\MENGEN \\
		R\text{-}\MOD \enspace\longrightarrow & \enspace\MENGEN \\
		\TOP \enspace\longrightarrow & \enspace\MENGEN \\
		\TOP_\bullet \enspace\longrightarrow & \enspace\MENGEN \\
		\TOP_\bullet \enspace\longrightarrow & \enspace \TOP
	\end{align*}
	\item Sei $\mathcal{C}$ eine Kategorie und $C$ ein Objekt von $\mathcal{C}$. Der durch $C$ \bet{dargestellte Funktor}\index{Funktor!dargestellter} 
	$F_C : \mathcal{C} \to \MENGEN$ ist definiert durch 
	\[
		\begin{array}{ll}
			F_C(A) = \Mor_\mathcal{C}(C,A) &\text{ für } A \in  \Ob(\mathcal{C}) \\
			F_C(f)  : \Mor_\mathcal{C}(C,A) \to \Mor_\mathcal{C}(C,A'), g \mapsto f \circ  g &\text{ für } f : A \to A' \text{ in } \mathcal{C}
		\end{array}
	\]
	\item Die Fundamentalgruppe definiert einen Funktor
	\[
		\pi_1 : \HTOP_\bullet \longrightarrow \GRUPPEN
	\]
\end{enumerate}
\minisec{Bemerkung}
\begin{itemize}
	\item Ist $F : \mathcal{C} \to \mathcal{D}$ ein Funktor so schreiben wir oft kürzer und ungenauer $f_* := F(f)$.
	\item Kleine Kategorien und Funktoren bilden die Kategorie $\KAT$.
\end{itemize}
% subsection 15 (end)

\subsection[Definition: Natürliche Transformation]{Definition} % (fold)
\label{sub:16}
Seien $F,G : \mathcal{C} \to \mathcal{D}$ zwei Funktoren. Eine \Index{natürliche Transformation} $\tau : F \to G$ ordnet jedem $C \in \Ob(\mathcal{C})$ einen Morphismus
$\tau_C : F(C) \to G(C)$ in $\mathcal{D}$ zu, sodass für jedes $f : C \to C'$ in $\mathcal{C}$
\[
	\begin{tikzcd}
		F(C) \rar["F(f)"] \dar["\tau_C"]& F(C') \dar["\tau_{C'}"] \\
		G(C) \rar["G(f)"] & G(C')
	\end{tikzcd}
\]
kommutiert.
\minisec{Bemerkung}
Natürliche Transformationen lassen sich komponieren. Für $\tau : F \to G$, $\eta : G \to H$ ist $\eta \circ \tau : F \to H$ gegeben durch
\marginnote{$F,G,H : \mathcal{C} \to \mathcal{D}$}
\[
	(\eta \circ \tau)_C :=  \eta_C \circ \tau_C : F(C) \to H(C)
\]
Genauer: Für eine feste kleine Kategorie $\mathcal{C}$ und $\mathcal{D}$ bilden die Funktoren $\mathcal{C} \to \mathcal{D}$ mit den natürlichen Transformationen eine 
Kategorie $\FUN(\mathcal{C}, \mathcal{D})$.
% subsection 16 (end)

\subsection[Ausblick auf Kategorien und Funktoren in der algebraischer Topologie]{Ausblick} % (fold)
\label{sub:17}
In der algebraischen Topologie werden topologische Fragen, wie zum Beispiel, wann $\mathds{R}^n$ und $\mathds{R}^m$ homöomorph sind, in algebraische Fragen übersetzt. Eine
Möglichkeit für eine solche Übersetzung sind Funktoren von einer Kategorie von topologischen Räumen, z.B. $\TOP, \TOP_\bullet, \HTOP_\bullet$, in eine algebraisch Kategorie,
z.B. $\GRUPPEN, K\text{-}\VR, \ABELGRUPPEN, R\text{-}\MOD$. Ein Beispiel für einen solchen Funktor ist die Fundamentalgruppe $\pi_1$. Ein Nachteil der Fundamentalgruppe ist,
dass diese oft schwierig zu berechnen ist. Wir werden in dieser Vorlesung weitere Funktoren und Methoden für ihr Berechnung kennenlernen.
% subsection 17 (end)

% section 1 (end)
\newpage

\section{Kombinatorische Beschreibungen topologischer Räume} % (fold)
\label{sec:2}
\subsection[Beispiel einer kombinatorischen Beschreibung eines Raumes]{Beispiel} % (fold)
\label{sub:21}
\Index{Graphen} sind einerseits topologische Räume und andererseits kombinatorische Objekte:
\begin{align*}
	\vcenter{\hbox{\begin{tikzpicture}[scale=1.2]
		\coordinate (A) at (0,1);
		\coordinate (B) at (1,1);
		\coordinate (C) at (2,1);
		\coordinate (1) at (0,0);
		\coordinate (2) at (1,0);
		\coordinate (3) at (2,0);
		\foreach \x in {1,2,3} {
			\foreach \y in {A,B,C}{
				\draw[thick] (\x) -- (\y);
			};
		}
		\foreach \x in {A,B,C} {
			\foreach \y in {1,2,3}{
				\draw[thick] (\x) -- (\y);
			};
			\draw[fill=DodgerBlue3, thick] (\x) circle[radius=0.05] node[above]{$\x$};
		}
		\foreach \x in {1,2,3} {
			\draw[fill=DodgerBlue3,thick] (\x) circle[radius=0.05] node[below]{$\x$};
		}
	\end{tikzpicture} }}\qquad 
	&\longleftrightarrow \quad \begin{array}{l}
		V =  \set{A,B,C,1,2,3} \\
		E =  \set{ \rule{0cm}{0.3cm} { \scriptstyle \set{A,1}, \set{A,2}, \set{A,3} , \set{B,1}, \set{B,2}, \set{B,3}, \set{C,1} , \set{C,2}, \set{C,3}}}  
	\end{array}\\
	\vcenter{\hbox{\begin{tikzpicture}[scale=1.2]
		\coordinate (1) at (0,1);
		\coordinate (2) at (0.5,1.5);
		\coordinate (3) at (1,1);
		\coordinate (4) at (1,0);
		\coordinate (5) at (0,0);
		\coordinate (6) at (0.5,0.5);
		\foreach \x in {2,3,5,6} {
			\draw[thick] (1) -- (\x);
		}
		\foreach \x/\y in {2/3,3/4,3/6,4/5,4/6,5/6}{
			\draw[thick] (\x) -- (\y);
		}
		\foreach \x/\y in {1/left,2/above,3/right,4/right,5/left,6/right}{
			\draw[fill=DodgerBlue3, thick] (\x) circle[radius=0.05] node[\y]{$\x$};
		}
	\end{tikzpicture} }}\qquad
	&\longleftrightarrow \quad \begin{array}{l}
		V = \set{1,2,3,4,5,6} \\
		E = \set{ \rule{0cm}{0.3cm} {\scriptstyle \set{1,2}, \set{1,3}, \set{1,5} , \set{1,6}, \set{2,3}, \set{3,4}, \set{3,6} , \set{4,5}, \set{4,6}, \set{5,6}} }  
	\end{array}
\end{align*}
% subsection 21 (end)

\subsection[Definition: Simplizialer Komplex]{Definition} % (fold)
\label{sub:22}
Ein \Index{simplizialer Komplex} $K=(V, \Sigma)$ besteht aus einer Menge $V$ und einer Menge $\Sigma$ von nichtleeren, endlichen Teilmengen von $V$, sodass gilt
\begin{enumerate}[(i)]
	\item $\set{v} \in \Sigma$ für alle $v \in V$
	\item $\sigma \in \Sigma$, $\emptyset \not= \tau \subseteq \sigma \Longrightarrow \tau \in\Sigma.$ \hfill {\footnotesize(Abgeschlossen bzgl. Teilmengenbildung)}
\end{enumerate}
Die Elemente von $V$ heißen die \Index{Ecken} oder \Index{Vertices} von $K$. Die Elemente von $\Sigma$ heißen die \Index{Simplizes} von $K$. Enthält $\sigma \in \Sigma$ 
genau $n+1$ Elemente, so heißt $\sigma$ ein \bet{$\mathbf{n}$-Simplex}\index{n-Simplex@$n$-Simplex}. Ist $\tau \subseteq \sigma$ mit $\sigma$ ein $n$-Simplex und $\tau$ ein $n-1$-Simplex, so heißt $\tau$ eine \Index{Seite} von $\sigma$. 

\noindent Ist $V$ geordnet, so heißt $K$ geordnet. Ist $V$ endlich, so heißt $K$ endlich.
% subsection 22 (end)

\subsection[Beispiel für simpliziale Komplexe]{Beispiel} % (fold)
\label{sub:23}
\begin{align*}
	\vcenter{\hbox{\begin{tikzpicture}[scale=1.]
		\coordinate (0) at (0,0);
		\coordinate (1) at (1,0.8);
		\coordinate (2) at (1,-0.7);
		\coordinate (3) at (2,0.2);
		\filldraw[DodgerBlue2!40] (0) -- (1) -- (2) -- cycle ;
		\foreach \x/\y in {0/1,0/2,1/2,1/3,2/3} {
			\draw[thick] (\x) -- (\y);
		};
		\foreach \x/\pos in {0/left,1/above,2/below,3/right} {
			\draw[fill=DodgerBlue3,thick] (\x) circle[radius=0.06] node[\pos]{$\x$};
		}
	\end{tikzpicture}}} \qquad 
	\begin{array}{l}
		V =  \set{0,1,2,3} \\
		E =  \set{ \rule{0cm}{0.3cm} {\emptyset, \set{0}, \set{1}, \set{2}, \set{3}, \set{0,1}, \set{0,2},\set{1,2},\set{1,3},\set{2,3},\set{0,1,2}}}  
	\end{array}
\end{align*}
Ist $\sigma$ eine endliche Menge, so heißt $\Delta^\sigma := \enbrace[\big]{\sigma, \mathcal{P}(\sigma) \setminus \set{\emptyset}}$ der $\sigma$-Simplex. Für 
$\sigma = \set{0, \ldots , n}$ schreiben wir $\Delta^n := \Delta^{\set{0, \ldots ,n}}$. Es ist
\[
	\Delta^0  = \vcenter{\hbox{\begin{tikzpicture}
		\draw[fill=DodgerBlue3,thick] (0,0) circle[radius=0.06];
	\end{tikzpicture}}} , \qquad 
	\Delta^1 = \vcenter{\hbox{\begin{tikzpicture}
		\coordinate (0) at (0,0);
		\coordinate (1) at (0.4,0.3);
		\draw[thick] (1) -- (0);
		\foreach \x in {0,1} {
			\draw[fill=DodgerBlue3,thick] (\x) circle[radius=0.06];
		}
	\end{tikzpicture}}} , \qquad 
	\Delta^2 = \vcenter{\hbox{\begin{tikzpicture}
		\coordinate (0) at (0,0);
		\coordinate (1) at (0.6,0);
		\coordinate (2) at (60:0.6);
		\draw[thick, fill=DodgerBlue2!40] (1) -- (0) -- (2) -- cycle;
		\foreach \x in {0,1,2} {
			\draw[fill=DodgerBlue3,thick] (\x) circle[radius=0.06];
		}
	\end{tikzpicture}}}, \qquad 
	\Delta^3 = 
	\vcenter{\hbox{\begin{tikzpicture}
			\coordinate (0) at (0,0);
			\coordinate (1) at (1,0);
			\coordinate (2) at (0.5,0,0.866);
			\coordinate (3) at (0.5,0.866,0.433);
			\foreach \x in {0,1,2,3}{
				\foreach \y in {0,1,2,3}{
					\foreach \z/\col in {0/DodgerBlue2,1/DodgerBlue1,2/DodgerBlue3,3/DodgerBlue2}{
					\draw[thick, fill=\col!40, fill opacity=0.4] (\x) -- (\y) -- (\z) -- cycle;
					};
				};
			};
			\foreach \x in {0,1,2,3} {
				\draw[fill=DodgerBlue3,thick] (\x) circle[radius=0.06];
			}
		\end{tikzpicture}}}
\]
% subsection 23 (end)

\subsection[Definition: Unterkomplex]{Definition} % (fold)
\label{sub:24}
Sei $K=(V, \Sigma)$ ein simplizialer Komplex. Ein \Index{Unterkomplex} von $K$ ist ein simplizialer Komplex $K_0 = (V_0, \Sigma_0)$ mit 
$V_0 \subseteq V$ und $\Sigma_0 \subseteq \Sigma$.
\minisec{Beispiel}
Sei $K=(V,\Sigma)$ ein endlicher simplizialer Komplex. Dann ist $K$ ein Unterkomplex von $\Delta^V$.
% subsection 24 (end)

\subsection[Definition: Simpliziale Abbildung]{Definition} % (fold)
\label{sub:25}
Eine \Index{simpliziale Abbildung} $f : K_1=(V_1, \Sigma_1) \to K_2(V_2, \Sigma_2)$ zwischen simplizialen Komplexen ist eine Abbildung $f : V_1 \to V_2$, sodass 
$f(\sigma_1) \in \Sigma_2$ für alle $\sigma_1 \in \Sigma_1$. 
% subsection 25 (end)

\subsection[Definition: Geometrische Realisierung]{Definition} % (fold)
\label{sub:26}
Sei $\sigma$ eine endliche Menge. Sei $\mathds{R}^\sigma = \prod_{v \in \sigma} \mathds{R}$ mit der Produkttopologie. Sei nun 
\[
	\abs*{\Delta^\sigma} := \set[x=(x_v)_{v \in \sigma} \in \mathds{R}^\sigma]{\sum_{v \in \sigma}x_v = 1 , x_v \in [0,1] \,\, \forall v \in \sigma}  
\]
Ist $\tau \subseteq \sigma$, so erhalten wir eine Abbildung $\iota_\tau^\sigma : \abs*{\Delta^\tau}  \to \abs*{\Delta^\sigma}$ indem wir 
$(x_v)_{v \in \tau} \in \abs*{\Delta^\tau}$ durch $x_v=0$ für $v \in \sigma \setminus \tau$ zu $(x_v)_{v \in \sigma}$ auffüllen. Ist $\eta \subseteq \tau \subseteq \sigma$ 
so gilt $\iota^\sigma_\eta = \iota_\tau^\sigma \circ  \iota_\eta^\tau$.

\noindent Sei $K=(V, \Sigma)$ ein simplizialer Komplex. Die \Index{geometrische Realisierung} $\abs*{K}$ von $K$ ist definiert als 
\[
	\abs*{K} := \coprod_{\sigma \in \Sigma} \set{\sigma} \times \abs*{\Delta^\sigma} \Big/ \sim
\] 
wobei $\sim$ die durch $(\tau, x) \sim \enbrace[\big]{\sigma, \iota_\tau^\sigma(x)}$ für $\tau \subseteq \sigma$ erzeugte Äquivalenzrelation ist. Versehen mit der 
Quotiententopologie ist $\abs*{K}$ ein topologischer Raum.
% subsection 26 (end)

\subsection[Beispiel: Rand des $n$-Simplizes und der Torus]{Beispiel} % (fold)
\label{sub:27}
\begin{itemize}
	\item Sei $\partial \Delta^n := \enbrace[\Big]{\set{0, \ldots ,n}, \mathcal{P}\enbrace[\big]{\set{0,\ldots ,n}} \setminus \set{\emptyset, \set{0, \ldots , n}}} = $ 
	\enquote{$\Delta^n \setminus \set{0,\ldots ,n}$}. $\partial \Delta^n$ heißt der \Index{Rand des $n$-Simplizes}. Es gilt
	\begin{align*}
		\abs*{\partial \Delta^1} &= \enspace\vcenter{\hbox{\begin{tikzpicture}
		\draw[fill] (0,0) circle[radius=0.04];
		\draw[fill] (0.15,0) circle[radius=0.04];
	\end{tikzpicture}}} \enspace\cong S^0\\ 
		\abs*{\partial \Delta^2} &=  \vcenter{\hbox{\begin{tikzpicture}
		\coordinate (0) at (0,0);
		\coordinate (1) at (0.6,0);
		\coordinate (2) at (60:0.6);
		\draw[thick] (1) -- (0) -- (2) -- cycle;
	\end{tikzpicture}}} \cong S^1 \\
		\abs*{\partial \Delta^3} &= \vcenter{\hbox{
		\begin{tikzpicture}
			\coordinate (0) at (0,0);
			\coordinate (1) at (1,0);
			\coordinate (2) at (0.5,0,0.866);
			\coordinate (3) at (0.5,0.866,0.433);
			\foreach \x in {0,1,2,3}{
				\foreach \y in {0,1,2,3}{
					\foreach \z/\col in {0/DodgerBlue2,1/DodgerBlue1,2/DodgerBlue3,3/DodgerBlue2}{
					\draw[thick, fill=\col!40, fill opacity=0.4] (\x) -- (\y) -- (\z) -- cycle;
					};
				};
			};
		\end{tikzpicture}}} \cong S^2 
	\end{align*}
	Allgemein gilt $\abs*{\partial \Delta^n} \cong S^{n-1}$.
	\item $T^2 = S^1 \times S^1 \cong 
	\vcenter{\hbox{\begin{tikzpicture}[scale=0.3]
		\draw (-3.5,0) .. controls (-3.5,2) and (-1.5,2.5) .. (0,2.5);
		\draw[xscale=-1] (-3.5,0) .. controls (-3.5,2) and (-1.5,2.5) .. (0,2.5);
		\draw[rotate=180] (-3.5,0) .. controls (-3.5,2) and (-1.5,2.5) .. (0,2.5);
		\draw[yscale=-1] (-3.5,0) .. controls (-3.5,2) and (-1.5,2.5) .. (0,2.5);

		\draw (-2,.2) .. controls (-1.5,-0.3) and (-1,-0.5) .. (0,-.5) .. controls (1,-0.5) and (1.5,-0.3) .. (2,0.2);

		\draw (-1.75,0) .. controls (-1.5,0.3) and (-1,0.5) .. (0,.5) .. controls (1,0.5) and (1.5,0.3) .. (1.75,0);
	\end{tikzpicture}}} \cong 
	\abs*{\vcenter{\hbox{\begin{tikzpicture}[>=to, scale=0.8]
		\filldraw[color=DodgerBlue2!20,] (0,0) rectangle (3,3);
		\foreach \x in {0,1,2}{
			\foreach \y in {0,1,2}{
				\draw (\x,\y) -- ++ (1,1);
			};
		};
		\draw (0,0) grid (3,3);
		\draw[very thick] (0,0) -- (0,3) -- (3,3) -- (3,0) -- cycle;
		\draw[->>, thick] (0,1) -- (0,1.5);
		\draw[->, thick] (1,3) -- (1.5,3);
		\draw[->, thick] (1,0) -- (1.5,0);
		\draw[->>, thick] (3,1) -- (3,1.5);
		\clip (-0.2,-0.2) rectangle (3.2,3.2);
	\end{tikzpicture}}}} $
\end{itemize}
% subsection 27 (end)

\subsection[Definition: Polyeder]{Definition} % (fold)
\label{sub:28}
Ein topologischer Raum $X$ heißt ein \Index{Polyeder}, falls er homöomorph zur Realisierung eines simplizialen Komplexes ist. 
% subsection 28 (end)

\subsection[Definition: Euler-Charakteristik]{Definition} % (fold)
\label{sub:29}
Sei $K=(V,\Sigma)$ ein endlicher simplizialer Komplex. Sei $a_n := \# \set[\sigma \in \Sigma]{\sigma \text{ ist ein $n$-Simplex von $K$}} $. Dann heißt 
$\chi(K) := \sum_{n \in \mathds{N}} (-1)^n a_n $ die \Index{Euler-Charakteristik} von $K$. 
\minisec{Beispiel}
\begin{align*}
	\chi \enbrace*{\vcenter{\hbox{\begin{tikzpicture}
		\coordinate (0) at (0,0);
		\coordinate (1) at (0.6,0);
		\coordinate (2) at (60:0.6);
		\draw[thick] (1) -- (0) -- (2) -- cycle;
	\end{tikzpicture}}} } &= 3 -3 = 0 \\
	\chi \enbrace[\bigg]{\vcenter{\hbox{\begin{tikzpicture}
		\coordinate (0) at (0,0);
		\coordinate (1) at (0.6,0);
		\coordinate (2) at (0.6,0.6);
		\coordinate (3) at (0,0.6);
		\draw[thick] (0) -- (1) -- (2) -- (3) -- cycle;
	\end{tikzpicture}}} } &= 4-4 =0 \\
	\chi \enbrace[\big]{\benbrace*{n\text{-Eck}}} &= n-n = 0
\end{align*}

\begin{align*}
	\chi \enbrace*{\partial \Delta^3 = \vcenter{\hbox{\begin{tikzpicture}
			\coordinate (0) at (0,0);
			\coordinate (1) at (1,0);
			\coordinate (2) at (0.5,0,0.866);
			\coordinate (3) at (0.5,0.866,0.433);
			\draw[sharp corners, dashed] (0) -- (1);
			\draw[sharp corners] (1) -- (2) -- (0) -- (3) -- (2) -- (1) -- (3) -- (0);
		\end{tikzpicture}}} } &= 4 - 6 +4 = 2 \\
	\chi\enbrace[\Bigg]{ \vcenter{\hbox{\begin{tikzpicture}
		\coordinate (0) at (0,0,-0.6);
		\coordinate (1) at (0.6,0,-0.6);
		\coordinate (2) at (0.6,0.6,-0.6);
		\coordinate (3) at (0,0.6,-0.6);
		\draw[thick, dashed] (3) -- (0) -- (1);
		\draw[thick] (1) -- (2) -- (3);
		\draw[thick, dashed] (0,0) -- (0);
		\draw[thick] (0.6,0) -- (0.6,0,-0.6);
		\draw[thick] (0.6,0.6) -- (0.6,0.6,-0.6);
		\draw[thick] (0,0.6) -- (0,0.6,-0.6) -- ++ (0.2,0);
		\draw[thick] (0,0) rectangle (0.6,0.6);
	\end{tikzpicture}}} } &= 8 -12 +6 = 2 \\
	\chi \enbrace*{ \vcenter{\hbox{\begin{tikzpicture}
		\coordinate (A1) at (0,0);
		\coordinate (A2) at (0.6,0.2);
		\coordinate (A3) at (1,0);
		\coordinate (A4) at (0.4,-0.2);
		\coordinate (B1) at (0.5,0.5);
		\coordinate (B2) at (0.5,-0.5);
		\begin{scope}[thick,dashed,]
		\draw (A1) -- (A2) -- (A3);
		\draw (B1) -- (A2) -- (B2);
		\end{scope}
		\begin{scope}[thick]
		\draw[] (A1) -- (A4) -- (B1);
		\draw[] (A1) -- (A4) -- (B2);
		\draw[] (A3) -- (A4) -- (B1);
		\draw[] (A3) -- (A4) -- (B2);
		\draw (B1) -- (A1) -- (B2) -- (A3) --cycle;
		\end{scope}
	\end{tikzpicture}}} } &= 6- 12  +8 = 2 
\end{align*}
% subsection 29 (end)

\subsection[Satz über die Euler-Charakteristik]{Satz} % (fold)
\label{sub:210}
Seinen $K$ und $K'$ endliche simpliziale Komplexe. Gilt $\abs*{K} \cong \abs*{K'}$, so gilt 
\[
	\chi(K)= \chi(K').
\]
\minisec{Bemerkung}
Die Euler-Charakteristik ist also eine topologische Invariante von simplizialen Komplexen. Sie hängt nur von der topologischen Struktur von $\abs*{K}$ und nicht
von der kombinatorischen Struktur von $K$ ab.
% subsection 210 (end)
% section 2 (end)
\newpage

\section{Simpliziale Homologie} % (fold)
\label{sec:3}

\subsection[Definition: Freier $R$-Modul mit Basis $S$]{Definition} % (fold)
\label{sub:31}
Sei $S$ eine Menge.Sei $R$ ein Ring.\marginnote{Ring mit Eins} Der \bet{freie $R$-Modul mit Basis $S$}\index{freier Modul}, $R[S]$, besteht aus allen endlichen formalen $R$-Linearkombinationen
\[
	\sum_{s \in S} r_s \cdot s \quad , r_s \in R, r_s \not=0 \text{ für endliche viele }s \in S
\]
Die $R$-Modulstruktur auf $R[S]$ ist definiert durch:
\begin{align*}
	r \cdot \enbrace*{\sum_{s\in S} r_s \cdot s} := \sum_{s \in S} (r \cdot r_s) \cdot s \qquad \qquad 
	\enbrace*{\sum_{s \in S} r_s \cdot s} + \enbrace*{\sum_{s \in S} r_s' s} := \sum_{s \in S} (r_s + r_s') \cdot s
\end{align*}
\minisec{Bemerkung}
Mittels $s= \sum_{s' \in S} \delta_{s,s'} \cdot s'$ fassen wir $S$ als Teilmenge von $R[S]$ auf.
% subsection 31 (end)

\subsection[Lemma: Universale Eigenschaft von freien Moduln]{Lemma} % (fold)
\label{sub:32}
Sei $S$ eine Menge, $R$ ein Ring. Es gilt\\
\begin{minipage}{0.7\textwidth}
	\begin{enumerate}[1)]
		\item Ist $M$ ein $R$-Modul und $i : S \to M$ eine Abbildung, so gibt es genau eine $R$-lineare Abbildung $\varphi : R[S] \to M$, die $i$ fortsetzt.
		\item Sei $f : R[S] \to M$ $R$-linear und $p : N \twoheadrightarrow M$ $R$-linear und surjektiv. Dann gibt es $\hat{f} : R[S] \to N$ mit $p \circ \hat{f} =f$.
	\end{enumerate}
\end{minipage}\hfill
\begin{minipage}{0.25\textwidth}
	\[
		\begin{tikzcd}
			\mbox{ } & N \dar["p", two heads] \\
			R[S] \rar["f"] \urar["\hat{f}", dashed] & M
		\end{tikzcd}
	\]
\end{minipage}
\minisec{Beweis}
\begin{enumerate}[1)]
	\item Eine solche Abbildung ist gegeben durch
	\(
		\varphi \enbrace*{\sum_{s \in S} r_s \cdot s} = \sum_{s \in S} r_s \cdot i(s)
	\).
	Ist $\varphi'$ eine zweite, so gilt 
	\begin{align*}
		\hat{\varphi} \enbrace*{\sum_{s \in S} r_s \cdot s}  = \sum_{s \in S} r_s \cdot \hat{\varphi}(s) = \sum_{s  \in S} r_s \cdot i(s)
	\end{align*}
	\item Wähle für jedes $s \in S$ ein Urbild $j(s)\in N$ für $f(s) \in M$ unter $p : N \twoheadrightarrow M$. Nun wende 1) auf $j : S \to N$ an, um $\hat{f} : R[S] \to N$
	mit $\hat{f}(s) = j(s)$ zu erhalten. Nun ist $ p \circ \hat{f}(s) = f(s)$ für alle $s \in S$. Mit der Eindeutigkeit aus 1) folgt $p \circ \hat{f} = f$. \bewende
\end{enumerate}
% subsection 32 (end)

\subsection[Beispiel für freie und nicht-freie Moduln]{Beispiel} % (fold)
\label{sub:33}
\begin{enumerate}[(i)]
	\item Sei $K$ ein Körper und $V,M,N$ seien $K$-Vektorräume, $f : V \to M$, $p : N \to M$ seien $K$-linear, $p$ surjektiv. Dann gibt es $\hat{f} : V \to N$ mit
	$p \circ \hat{f} = f$. Wir können Lemma \ref{sub:32} benutzen, da $V$ eine Basis $B$ hat, also $V \cong K[B]$.
	\item Sei $R=\mathds{Z}$, sei $V= \nicefrac{\mathds{Z}}{2 \mathds{Z}}$ und $M = \nicefrac{\mathds{Z}}{2 \mathds{Z}}$, sowie $f= \id : V \to M$, $N= \mathds{Z}$ und 
	$p : \mathds{Z}n \to \nicefrac{\mathds{Z}}{2 \mathds{Z}}$ die Projektion mit $p(n)= n + 2 \mathds{Z}$. Dann gibt es keine $\mathds{Z}$-lineare Abbildung 
	$\hat{f} : \nicefrac{\mathds{Z}}{2 \mathds{Z}} \to \mathds{Z}$ mit $p \circ \hat{f} = f$.
	Ist $\hat{f} (1+ 2 \mathds{Z}) = n \in \mathds{Z}$, so folgt 
	\[
		2 n = 2 \cdot \hat{f}(1+ 2 \mathds{Z}) = \hat{f} \enbrace[\big]{2 \cdot (1+ 2 \mathds{Z})} = \hat{f} (0+ 2 \mathds{Z}) =0
	\]
	also $n=0$.
\end{enumerate}
% subsection 33 (end)

\subsection[Definition: $n$-ter Kettenmodul]{Definition} % (fold)
\label{sub:34}
Sei $K=(V, \Sigma)$ ein simplizialer Komplex. Sei $\Sigma_n := \set[\sigma \in \Sigma]{\sigma \text{ ist ein $n$-Simplex}}$. Sei $R$ ein Ring. Der 
\bet{$\mathbf{n}$-te Kettenmodul}\index{n-ter Kettenmodul@$n$-ter Kettenmodul!simplizial} von $K$ über $R$ ist definiert als 
\[
	C_n(K;R) := R[\Sigma_n]
\]
Ist $R=\mathds{Z}$, so schreiben wir auch kurz $C_n(K) := C_n(K;\mathds{Z})$.
% subsection 34 (end)

\subsection[Definition: Orientierung eines $(n-1)$-Simplizes als Teilmenge eines $n$-Simplizes]{Definition} % (fold)
\label{sub:35}
Sei $K=(V,\Sigma)$ ein geordneter simplizialer Komplex. Sei $\sigma = \set{v_0, \ldots ,v_n}$ ein $n$-Simplex von $K$. Bezüglich der Ordnung von $K$ sei dabei 
$v_0 < v_1 < \ldots < v_n$. Sei $\tau$ ein $(n-1)$-Simplex von $K$. Definiere
\[
	\varepsilon_\sigma^\tau := \begin{cases}
		(-1)^i, &\text{ falls }\tau = \set{v_0, \ldots , v_{i-1}, v_{i+1}, \ldots , v_n} \\
		0 , &\text{ sonst}
	\end{cases}
\]
\minisec{Bemerkung}
$\varepsilon^\tau_\sigma \not= 0 \iff \tau$ ist eine Seite von $\sigma$.
\minisec{Beispiel}
Für den folgenden geordneten simplizialen Komplex gilt
\[
	\vcenter{\hbox{\begin{tikzpicture}
		\coordinate (0) at (0,0);
		\coordinate (1) at (0.5,0.866);
		\coordinate (2) at (1,0);
		\draw[very thick, 
			postaction={decorate},
			decoration={markings,
				mark=at position 1/6 with {\arrow{to}},
				mark=at position 0.5 with {\arrow{to}},
				mark=at position 5/6 with {\arrow{to reversed}}}] 
		(0) -- (1) -- (2) -- cycle;
		\foreach \x/\pos in{0/left,1/above,2/right}{
			\node[\pos] at (\x) {$\x$};
		}
	\end{tikzpicture}}} \qquad 
	\varepsilon_{\set{0,1,2}}^{\set{0,1}} = 1 , \quad \varepsilon_{\set{0,1,2}}^{\set{0,2}}=-1 , \quad \varepsilon_{\set{0,1,2}}^{\set{1,2}} =1    
\]

% subsection 35 (end)

\subsection[Definition: $n$-te Randabbildung]{Definition} % (fold)
\label{sub:36}
Sei $K=(V,\Sigma)$ ein geordneter simplizialer Komplex. Die \bet{$\mathbf{n}$-te Randabbildung}\index{n-te Randabbildung@$n$-te Randabbildung!simplizial} 
\[
	\partial_n : \underset{=R[\Sigma_n]}{C_n(K;R)} \longrightarrow \underset{=R[\Sigma_{n-1}]}{C_{n-1}(K;R)}
\]
ist definiert durch
\[
	\partial_n(\sigma) = \sum_{\mathclap{\tau \in \Sigma_{n-1}}} \varepsilon_\sigma^\tau \cdot \tau \qquad \text{ für } \sigma \in \Sigma_n
\]
% subsection 36 (end)

\subsection[Proposition: $\partial_{n-1} \circ \partial_n = 0$]{Proposition} % (fold)
\label{sub:37}
\[
	\partial_{n-1} \circ \partial_n = 0
\]
Damit folgt $\im \partial_n \subseteq \ker \partial_{n-1}$.
\minisec{Beweis (mit Lemma \ref{sub:38})}
Sei $\sigma$ ein $n$-Simplex von $K$. Dann gilt
\begin{align*}
	\partial_{n-1} \enbrace[\big]{\partial_n (\sigma)} = \partial_{n-1} \enbrace*{ \sum_{\tau \in \Sigma_{n-1}} \varepsilon_\sigma^\tau \cdot \tau} = 
	\sum_{\mathclap{\tau \in \Sigma_{n-1}}} \varepsilon_{\sigma}^\tau \cdot \partial_{n-1} (\tau) 
	 &= \smashoperator[r]{\sum_{\tau \in \Sigma_{n-1}}} \varepsilon_\sigma^\tau \enbrace*{ \sum_{\eta \in \Sigma_{n-2}} \varepsilon_\tau^\eta \cdot \eta} \\
	 &= \sum_{\eta \in \Sigma_{n-2}} \enbrace*{ \sum_{\tau \in \Sigma_{n-1}, \eta \subseteq \tau \subseteq \sigma} \!\!\!\!\! \! \! \varepsilon_\tau^\eta \cdot \varepsilon_\sigma^\tau}
	 \eta  \\
	 &\stackrel{\mathclap{\text{Lemma \ref{sub:38}}}}{=} \quad \enspace 0 \bewende
\end{align*}
% subsection 37 (end)

\subsection[Lemma: Ein $(n-2)$-Simplex ist Seite von genau zwei $(n-1)$-Simplizes]{Lemma} % (fold)
\label{sub:38}
Sei $\sigma$ ein $n$-Simplex. Sei $\eta \subseteq \sigma$ ein $(n-2)$-Simplex. Dann gibt es genau zwei $(n-1)$-Simplizes $\tau, \tau'$ von $K$, die eine Seite von $\sigma$
sind und $\eta$ als Seite enthalten. Es gilt
\[
	\varepsilon_\tau^\eta \cdot \varepsilon^\tau_\sigma = - \varepsilon_{\tau'}^\eta \cdot \varepsilon^{\tau'}_\sigma
\]
\minisec{Beweis}
Sei $\sigma= \set{v_0, \ldots ,v_n}$ mit $v_0 < v_1 < \ldots < v_n$. Dann ist $\eta = \set{v_0, \ldots , v_n} \setminus \set{v_i, v_j}$ mit $i<j$. Dann sind 
$\tau = \set{v_0, \ldots , v_{i-1}, v_{i+1}, \ldots , v_n}$ und $\tau' = \set{v_0, \ldots , v_{j-1}, v_{j+1}, \ldots , v_n}$ die gesuchten $(n-1)$-Simplizes und es gilt
\begin{align*}
	\varepsilon_\sigma^\tau = (-1)^i \enspace, \enspace \varepsilon_\sigma^{\tau'} = (-1)^j \enspace, \enspace \varepsilon^\eta_\tau = (-1)^{j-1} \enspace, \enspace 
	\varepsilon_{\tau'}^\eta = (-1)^i \bewende
\end{align*}
% subsection 38 (end)

\subsection[Notation für das Weglassen eines Knotens]{Notation} % (fold)
\label{sub:39}
Für $\sigma \in \Sigma_n$, $ \sigma= \set{v_0, \ldots ,v_n} $ mit $v_0 < \ldots < v_n$ schreiben wir 
\[
	\delta^j \sigma := \set{v_0 , \ldots , v_{j-1}, v_{j+1}, \ldots , v_n}  \in \Sigma_{n-1}
\]
Dann ist $\partial_n(\sigma) = \sum_{j=0}^{n} (-1)^j \cdot \delta^j \sigma$.
% subsection 39 (end)

\subsection[Definition: $n$-ter Homologiemodul]{Definition} % (fold)
\label{sub:310}
Sei $K$ ein geordneter simplizialer Komplex. Der \bet{$\mathbf{n}$-te Homologiemodul}\index{n-ter Homologiemodul@$n$-ter Homologiemodul!simplizial} von 
$K$ über $R$ ist definiert als
\[
	H_n(K;R) := \frac{\ker \partial_n : C_n(K;R) \to C_{n-1}(K;R)}{\im \partial_{n+1} : C_{n+1}(K;R) \to C_n(K;R)} 
\]
Für $n=0$ interpretieren wir $\partial_0$ als die Nullabbildung. Daher gilt 
\[
	H_0(K;R) = \frac{C_0(K;R)}{\im \partial_1 : C_1(K;R) \to C_0(K;R)} 
\]
% subsection 311 (end)

\subsection[Beispiele für Homologiemoduln]{Beispiele} % (fold)
\label{sub:311}
\begin{enumerate}[(1)]
	\item Sei $K= \bullet$. Dann ist $C_0(K;R) \cong R$ und $C_i(K;R) = 0$ für $i>0$. Weiter ist $H_0(K;R) = C_0(K;R) \cong R$ und $H_i(K;R) =0$ für $i>0$.
	\item Sein nun $K=K_n$ das $n$-Eck, wobei $n \ge 3$.  Also
	\begin{align*}
		\Sigma_0 &= \set{\tau_1=\set{1}, \tau_2=\set{2}, \ldots , \tau_n = \set{n} \rule{0cm}{0.9em}} \\
 		\Sigma_1 &= \set{\sigma_1 = \set{1,2}, \ldots , \sigma_{n-1} = \set{n_1, n}, \sigma_n=\set{n,1} \rule{0cm}{0.9em}} 
	\end{align*}
	Dann gilt
	\begin{align*}
		C_0(K;R) &= R[\Sigma_0] \cong R^{n} \\
		C_1(K;R) &= R[\Sigma_1] \cong R^{n} \\
		C_i(K;R) &= 0 \text{ für $i>1$, insbesondere } H_i(K;R) = 0 \text{ für } i>1
	\end{align*}
	Es ist dann $\partial_2 = \partial_3 = \ldots  = 0$. Interessant ist $\partial_1 : C_1(K;R) \to C_0(K;R)$, denn es gilt
	\[
		\partial_1(\sigma_i) = \begin{cases}
			\tau_{i+1} - \tau_i, &\text{ falls }i=1, \ldots ,n-1\\
			\tau_n - \tau_1 , &\text{ falls } i=n
		\end{cases}
	\]
	Mit $\sigma_1' := \sigma_1, \sigma_2' := \sigma_2, \ldots , \sigma_{n-1}' := \sigma_{n-1}$ und $ \sigma_n' := - \sigma_n$ gilt dann
	\[
		\partial_1 (\sigma_i') = \tau_{i+1} - \tau_i
	\]
	mit der Konvention $\tau_{n+1} = \tau_1$. Also gilt
	\begin{align*}
		\partial_1 \enbrace*{\sum_{i=1}^{n} r_i \sigma_i' } = \sum_{i=1}^{n} r_i \enbrace*{\tau_{i+1}- \tau_i} = \sum_{i=1}^{n} (r_{i-1} - r_i) \tau_i  
	\end{align*}
	Es folgt
	\begin{align*}
		\ker \partial_1 &= \set[\sum_{i=1}^{n} r_i \cdot \sigma_i']{r_1=r_2= \ldots =r_n} \subseteq C_1(K;R) \\
		\im \partial_1 &= \set[\sum_{i=1}^{n} s_i \cdot \tau_i]{\sum_{i=1}^{n} s_i =0} \subseteq C_0(K;R)
	\end{align*}
	Ist $\sum_{i=1}^{n} s_i =0$, so folgt für $r_n=0, r_1 = -s_1, r_2=-s_1 -s_2, \ldots , r_{n-1} = -s_1 - s_2 - \ldots - s_{n-1}$
	\[
		\partial_1 \enbrace*{\sum_{i=1}^{n} r_i \sigma_i'} = \sum_{i=1}^{n} s_i \tau_i 
	\]
	Es folgt ${\color{Firebrick1}[K] \in}\, H_1(K;R) = \nicefrac{\ker \partial_1}{\im \partial_2} = \ker \partial_1 \cong R \,{\color{Firebrick1}\ni 1}$ und 
	\[
		{ \color{Gold3} [\tau_i] =[p] \in }\, H_0(K;R) = \nicefrac{\ker \partial_0}{\im \partial_1} = \nicefrac{R[\Sigma_0]}{ \set[\sum_{i=1}^{n} s_i \tau_i]{\sum_{i=1}^{n} s_i =0} } \cong R \, { \color{Gold3} \ni 1 }
	\]
	(via $\sum_{i=1}^{n} s_i \tau_i \longmapsto \sum_{i=1}^{n} s_i$ und Homomorphiesatz)
\end{enumerate}
% subsection 311 (end)

\subsection[Definition: $R$-Kettenkomplex, $n$-te Homologie eines $R$-Kettenkomplexes]{Definition} % (fold)
\label{sub:312}
Sei $R$ ein Ring. Ein \bet{$\mathbf{R}$-Kettenkomplex}\index{R-Kettenkomplex@$R$-Kettenkomplex} $(C_*, d_*)$ ist eine Folge von $R$-Moduln $(C_n)_{n \in \mathds{N}}$ zusammen mit $R$-linearen Abbildungen 
$(d_n : C_n \to C_{n-1})_{n \ge 1}$, so dass $d_n \circ d_{n+1} = 0$. Die $d_n$ heißen die \Index{Randabbildungen} von $(C_*,d_*)$, die $C_n$ die \Index{Kettenmoduln}.
Die \bet{$\mathbf{n}$-te Homologie}\index{n-te Homologie@$n$-te Homologie} eines Kettenkomplexes $(C_*,d_*)$ ist definiert als
\[
	H_n(C_*, d_*) := \frac{\ker d_n : C_n \to C_{n-1}}{\im d_{n+1} : C_{n+1} \to C_n} 
\]
% subsection 312 (end)

\subsection[Bemerkungen zu Kettenkomplexen]{Bemerkung} % (fold)
\label{sub:313}
\begin{enumerate}[(i)]
	\item Oft werden auch Kettenkomplexe betrachtet, die $\mathds{Z}$ statt $\mathds{N}$ verwenden.
	\item Ein Kettenkomplex $(C_*,d_*)$ heißt endlich erzeugt, wenn alle $C_n$ endlich erzeugte $R$-Moduln sind und $C_n \not= 0$ nur für endlich viele $n$ ist.
\end{enumerate}
% subsection 313 (end)

\subsection[Definition: Simplizialer Kettenkomplex]{Definition} % (fold)
\label{sub:314}
Zu einem geordneten simplizialen Komplex $K$ heißt $\enbrace[\big]{C_*(K;R), \partial_*}$ der \bet{simpliziale Kettenkomplex}\index{simplizialer Kettenkomplex} von $K$ über 
$R$. 
\minisec{Bemerkung}
Ist $K$ endlich, so ist $\enbrace[\big]{C_*(K;R), \partial_*}$ endlich erzeugt.
% subsection 314 (end)
% section 3 (end)
\newpage

\section{Euler-Charakteristik von Kettenkomplexen} % (fold)
\label{sec:4}

\subsection[Wiederholung: Torsionsgruppe]{Wiederholung} % (fold)
\label{sub:41}
Sei $A$ eine abelsche Gruppe. Dann ist 
\[
	T A := \set[a \in A]{\exists n \ge 1 : n \cdot a = \underbrace{a + \ldots + a}_{\text{$n$-mal}} = 0} 
\]
die \Index{Torsionsgruppe} von $A$. Ist $A$ endlich erzeugt, so ist auch $TA$ endlich erzeugt und es gibt Primzahlpotenzen $p_1^{n_1}, \ldots , p_k^{n_k}$ mit
\marginnote{abelsche Gruppen sind $\mathds{Z}$-Moduln}
\[
	TA \cong \nicefrac{\mathds{Z}}{(p_1^{n_1})} \oplus \ldots \oplus \nicefrac{\mathds{Z}}{(p_k^{n_k})}
\]
Weiter gibt es dann $n$ mit $A \cong \mathds{Z}^n \oplus TA$. $\Rg A := n$ ist der \bet{Rang} von $A$\index{Rang einer abelschen Gruppe}.
% subsection 41 (end)

\subsection[Lemma: Rang einer endlich erzeugten abelschen Gruppe]{Lemma} % (fold)
\label{sub:42}
Sei $A$ eine endlich erzeugte abelsche Gruppe. Dann gilt $\Rg A = \max \set[m]{\exists C \le A, C \cong \mathds{Z}^m}$.
\minisec{Beweis}
Da $A \cong \mathds{Z}^n \oplus T A$ ist, genügt zu zeigen: Ist $\varphi : \mathds{Z}^m \to \mathds{Z}^n \oplus T A$ ein injektiver Gruppenhomomorphismus, so ist $m \le n$.
Schreibe $\varphi = \varphi_0 \oplus \varphi_1$, also $\varphi(a) = \varphi_0(a) + \varphi_1(a)$ mit $\varphi_0 : \mathds{Z}^m \to \mathds{Z}^n$, 
$\varphi_1 : \mathds{Z}^m \to T A$. 

Behauptung: $\varphi_0$ ist injektiv. Zu $v \in \mathds{Z}^m$ wähle $k\ge 1$ mit $k \cdot \varphi_1(v) = 0 \in T A$. Dann ist $\varphi_1(k v) = k \varphi_1(v)=0$.
Ist $v \not= 0$, so ist $k \cdot v \not= 0 \in \mathds{Z}^m$, also ist $\varphi(k v) \not= 0$ und damit $\varphi_0(k v) \not= 0$. Es folgt $\varphi_0(v) \not= 0$, da sonst
$\varphi_0(k \cdot v) = k \cdot \varphi_0(v)= 0$. Damit ist die Behauptung gezeigt.

Gruppenhomomorphismen $\varphi_0 : \mathds{Z}^m \to \mathds{Z}^n$ werden durch $n \times m$-Matrizen beschrieben: Es gibt $A= (a_{ij}) \in \mathds{Z}^{n \times m}$ mit
\[
	\varphi_0 \begin{pmatrix}
		z_1 \\ \vdots \\ z_m
	\end{pmatrix} = \begin{pmatrix}
		\sum_{j=1}^{m} a_{1j} z_j \\
		\vdots \\
		\sum_{j=1}^{m} a_{nj} z_j
	\end{pmatrix}
\]
Ist $m >n$, so hat $A$, aufgefasst als Matrix über $\mathds{Q}$, einen Kern. Es gibt also einen Vektor $w = \begin{psmallmatrix}a_1/b_1 \\ : \\ a_m/b_m\end{psmallmatrix}$
mit $A w = 0$, $w \not=0$. Dann ist $(b_1, \ldots , b_m) \cdot w \in \mathds{Z}^m$ und 
\[
	\varphi_0 \enbrace[\big]{(b_1, \ldots , b_m)\cdot w} = A \cdot \enbrace[\big]{(b_1, \ldots , b_m) \cdot w} = (b_1, \ldots , b_m) \cdot A \cdot w = 0 
\]
Also gilt $m \le n$. \bewende
% subsection 42 (end)

\subsection[Definition: Kurze exakte Sequenz]{Definition} % (fold)
\label{sub:43}
Seien $M_0, M_1, M_2$ $R$-Moduln und $f_0 : M_0 \to M_1$, $f_1 : M_1 \to M_2$ $R$-lineare Abbildungen. Dann heißt 
\[
	\begin{tikzcd}
		M_0 \rar["f_0", hook] & M_1 \rar["f_1", two heads]& M_2 
	\end{tikzcd} \tag*{$(\star)$}
\]
eine \Index{kurze exakte Sequenz}, wenn gilt:
\begin{multicols}{3}
	\begin{enumerate}[(i)]
		\item $f_0$ ist injektiv,
		\item $\im f_0 = \ker f_1$,
		\item $f_1$ ist surjektiv.
	\end{enumerate}
\end{multicols}
\minisec{Bemerkung}
Oft sagt man $M_0 \xrightarrow{f_0} M_1 \xrightarrow{f_1} M_2$ ist exakt in $M_1$, wenn $\im f_0 = \ker f_1$ ist. Dann ist $(\star)$ eine kurze exakte Folge, wenn
\[
	\begin{tikzcd}
		0 \rar & M_0 \rar["f_0"] & M_1 \rar["f_1"]& M_2 \rar & 0
	\end{tikzcd}
\]
exakt in $M_0, M_1$ und $M_2$ ist.
% subsection 43 (end)

\subsection[Beispiele für kurze exakte Folgen]{Beispiele} % (fold)
\label{sub:44}
\begin{enumerate}[1)]
	\item \[
		\begin{tikzcd}[row sep=0.3em]
			M_0 \rar["i",hook] & M_0 \oplus M_1 \rar["p",two heads] & M_1 \\
			v_0 \rar["i", mapsto] & (v_0,0) \enspace (v_0,v_1) \rar["p",mapsto] & v_1
		\end{tikzcd}
	\]
	ist eine kurze exakte Folge.
	\item $\begin{tikzcd}
		\mathds{Z} \rar[hook,"\cdot n"] & \mathds{Z} \rar[two heads] &  \nicefrac{\mathds{Z}}{n \mathds{Z}}
	\end{tikzcd}$
	ist eine kurze exakte Folge.
\end{enumerate}
% subsection 44 (end)

\subsection[Bemerkung: Spaltung einer kurzen exakten Sequenz]{Bemerkung} % (fold)
\label{sub:45}
Eine \Index{Spaltung} für eine kurze exakte Folge $M_0 \xrightarrow{f_0} M_1 \xrightarrow{f_1} M_2$ ist eine $R$-lineare Abbildung $s : M_2 \to M_1$ mit $f_1 \circ s = \id_{M_2}$. In diesem Fall erhalten wir einen Isomorphismus $M_0 \oplus M_2 \to M_1$, $(v_0, v_2) \mapsto f_0(v_0) + s(v_2)$.
\begin{description}
	\item[Injektivität:] Sei $(v_0,v_2) \in M_0 \oplus M_2$ mit $f_0(v_0) + s(v_2) =0$. Dann gilt
	\[
		0 = f_1 \enbrace[\big]{f_0(v_0) + s(v_2)} = v_2 \enspace \Longrightarrow \enspace f_0(v_0)=0 \enspace\Longrightarrow \enspace v_0 =0
	\]
	\item[Surjektivität:] Sei $v_1 \in M_1$. Betrachte $v_2 := f_1(v_1)$. Dann ist $v_1 -s(v_2) \in \ker f_1 = \im f_0$, also gibt es $v_0$ in $M_0$ mit
	$f_0(v_0) = v_1 - s(v_1)$. Damit ist $v_1= f_0(v_0)+ s(v_1)$.
\end{description}
% subsection 45 (end)

\subsection[Bemerkung: Existenz einer Spaltung, wenn $M_2$ frei ist]{Bemerkung} % (fold)
\label{sub:46}
\begin{enumerate}[1)]
	\item $\begin{tikzcd}
		\mathds{Z} \rar[hook,"\cdot n"] & \mathds{Z} \rar[two heads] &  \nicefrac{\mathds{Z}}{n \mathds{Z}}
	\end{tikzcd}$
	spaltet nicht.
	\item Ist $\begin{tikzcd}
		M_0 \rar["f_0"] & M_1 \rar["f_1"] & M_2
	\end{tikzcd}$
	eine kurze exakte Folge mit $M_2$ frei, also $M_2 \cong R[S]$, so spaltet $(\star)$.
	Siehe Lemma \ref{sub:32}. Insbesondere ist $M_1 \cong M_0 \oplus M_2$.
	
	Moduln mit dieser Eigenschaft bezeichnet man als \bet{projektiv}\index{projektiver Modul}. Freie Moduln sind also projektiv.
\end{enumerate}
% subsection 46 (end)

\subsection[Bemerkung: Der Rang ist additiv]{Bemerkung} % (fold)
\label{sub:47}
Der Rang für alle endlich erzeugten abelschen Gruppen ist additiv:
\[
	\Rg(A \oplus B) =\Rg A + \Rg B,
\]
da $A \cong \mathds{Z}^{\Rg A} \oplus T A$, $B \cong \mathds{Z}^{\Rg B} \oplus TB$, also 
\[
	A \oplus B \cong \mathds{Z}^{\Rg A + \Rg B} \oplus \underbrace{T A \oplus T B}_{=T(A \oplus B)}
\]
% subsection 47 (end)

\subsection[Proposition: Der Rang ist für kurze exakte Folgen endl. erz., abelscher Gruppen additiv]{Proposition} % (fold)
\label{sub:48}
Der Rang von endlich erzeugten abelschen Gruppen ist additiv für kurze exakte Folgen: Ist
\begin{equation*}
	\begin{tikzcd}
		A \rar["i",hook] & B \rar["p", two heads] & C
	\end{tikzcd} \label{eq:48}\tag{$\star$}
\end{equation*}
eine kurze exakte Folge von endlich 
erzeugten abelschen Gruppen, so gilt $\Rg B = \Rg A + \Rg C$.
% \minisec{Beispiel}
% mit $\mathds{Z}$
\minisec{Beweis}
Spaltet \eqref{eq:48}, so ist $B \cong A \oplus C$ und die Behauptung folgt aus der Bemerkung \ref{sub:47}. \smallskip \\
Wir können annehmen, dass $C = \mathds{Z}^n \oplus T C$ gilt. Wir erhalten folgendes Diagramm: 
\[
	\begin{tikzcd}
		A \rar["i",hook] \dar[equal] & p ^{-1}(\mathds{Z}^n) \rar["p", two heads] \dar[hook,"i''"]& \mathds{Z}^n \dar[hook,"i'"]\\
		A \rar["i",hook] & B \rar["p", two heads] \dar[two heads,"p''"]& C \dar[two heads,"p'"] & \text{\eqref{eq:48}}\\
		& \nicefrac{B}{p ^{-1}(\mathds{Z}^n)} \rar[two heads, hook,"p' \circ p"]& T C
	\end{tikzcd} 
\]
Ist $p' \circ p(b)=0$, so gibt es $v \in \mathds{Z}^n$ mit $i'(v)= p(b)$. Nun gibt es $v' \in p ^{-1}(\mathds{Z}^n)$ mit $p(v')=v$. Es folgt 
\[
	p \enbrace[\big]{i''(v')} = p(b) 
\]
also $i''(v') -b \in \ker p = i(A)$. Da $i(A) \subseteq i'' \enbrace*{p ^{-1} (\mathds{Z}^n)}$, folgt $b \in i'' \enbrace*{p ^{-1}(\mathds{Z}^n)}$.
In dem Diagramm sind die Spalten und Zeilen exakt. Da $\mathds{Z}^n$ frei ist, spaltet 
$\begin{tikzcd}
	A \rar[hook,"i"] & p^{-1}(\mathds{Z}^n) \rar["p", two heads] & \mathds{Z}^n
\end{tikzcd}$
und es gilt 
\[
	\Rg A + \Rg C = \Rg A + n = \Rg \enbrace*{p ^{-1}(\mathds{Z}^n)} 
\]
Es bleibt zu zeigen: $\Rg B = \Rg  \enbrace*{p ^{-1}(\mathds{Z}^n)}$. \medskip \\
Nebenbei: Ist $A' \subseteq B'$, so gilt $\Rg A' \le \Rg B'$, denn 
\[
	\Rg A' = \max \set[n']{\mathds{Z}^{n'} \cong C' \le A'} \le \max \set[n']{\mathds{Z}^{n'} \cong C' \le B'} = \Rg B'.
\]
Betrachte die kurze exakte Folge
$\begin{tikzcd}
	p^{-1}(\mathds{Z}^n) \rar & B \rar & \nicefrac{B}{p ^{-1}(\mathds{Z}^n)}
\end{tikzcd}$. 
Wegen $\nicefrac{B}{p ^{-1}(\mathds{Z}^n)} \cong T C$ gibt es ein
$k$ mit $k \cdot \enbrace*{b+ p ^{-1}(\mathds{Z}^n)} =0$ für alle $b \in B$. Also $k \cdot b \in p ^{-1}(\mathds{Z}^n)$. Es folgt 
$kB \subseteq  p ^{-1}(\mathds{Z}^n)$. Ist  $B \cong \mathds{Z}^{\Rg B}+ TB$, so ist $k B \cong (k \mathds{Z})^{\Rg B}+ T(kB)$ und $\Rg kB = n = \Rg B$. 
Mit der Nebenbemerkung folgt $\Rg p ^{-1}(\mathds{Z}^n) = \Rg B$ aus $kB \subseteq  p ^{-1}(\mathds{Z}^n) \subseteq B$. \bewende
\minisec{Beweis mit Tensorprodukten (Skizze)}
\begin{itemize}
	\item $\mathds{Q} \otimes_\mathds{Z} : \mathds{Z}\text{-}\MOD \to \mathds{Q}\text{-}\VR $ ist ein Funktor.
	\item $\Rg A= \dim_\mathds{Q} (\mathds{Q} \otimes_\mathds{Z} A)$
	\begin{enumerate}[(a)]
		\item $\mathds{Q} \otimes_\mathds{Z} (X \oplus Y) = \mathds{Q} \otimes_\mathds{Z} X \oplus  \mathds{Q} \otimes_\mathds{Z} Y$
		\item $\mathds{Q} \otimes_\mathds{Z} \mathds{Z}^n \cong \mathds{Q}^n$
		\item $\mathds{Q} \otimes_\mathds{Z} T \cong 0$ für $T$ ein Torsionsmodul.
	\end{enumerate}
	\item Ist $\begin{tikzcd}
		A \rar & B \rar & C
	\end{tikzcd}$ eine kurze exakte Folge, so ist auch 
	\[
		\begin{tikzcd}
				\mathds{Q} \otimes_\mathds{Z} A \rar & \mathds{Q} \otimes_\mathds{Z} B \rar & \mathds{Q} \otimes_\mathds{Z} C
			\end{tikzcd}
	\] 
	eine kurze exakte Folge.
	\item Jede kurze exakte Folge $\begin{tikzcd}
		V_0 \rar & V_1 \rar & V_2
	\end{tikzcd}$ von $\mathds{Q}$-Vektorräumen spaltet und daher gilt $V_1 \cong V_0 \oplus V_2$ und $\dim_\mathds{Q}(V \oplus W) = \dim_\mathds{Q}(V)+ \dim_\mathds{Q} (W)$.
\end{itemize}
% subsection 48 (end)

\subsection[Definition: Eulercharakteristik von $\mathds{Z}$-Kettenkomplexen]{Definition} % (fold)
\label{sub:49}
Sei $(C_*,d_*)$ ein endlich erzeugter $\mathds{Z}$-Kettenkomplex.
\[
	\chi(C_*) := \sum_{i=0}^{\infty} \Rg(C_i) = \sum_{i=0}^{n_0} \Rg(C_i) 
\]
heißt die \bet{Eulercharakteristik}\index{Euler-Charakteristik!von $\mathds{Z}$-Kettenkomplexen} von $(C_*,d_*)$. Dabei ist $n_0$ so gewählt, dass $C_n =0$ für alle 
$n >n_0$. 
% subsection 49 (end)

\subsection[Satz: Eulercharakteristik eines endlich erzeugten $\mathds{Z}$-Kettenkomplexes]{Satz} % (fold)
\label{sub:410}
Sei $(C_*,d_*)$ ein endlich erzeugter $\mathds{Z}$-Kettenkomplex. Dann gilt 
\[
	\chi(C_*) = \sum_{i=1}^{\infty} (-1)^i \Rg H_i(C_*,d_*).
\]
\minisec{Beweis}
Sei $B_{n+1} := \im d_{n+1} : C_{n+1} \to C_n$ und $Z_n := \ker d_n : C_n \to C_{n-1}$. Also $H_n(C_*,d_*) = \nicefrac{Z_n}{B_{n+1}}$.
Als Untermoduln von $C_n$ sind $Z_n$ und $B_n$ endlich erzeugt (LA2). Insbesondere ist auch $H_n(C_*,d_*)$ endlich erzeugt und der Rang somit definiert. Auch $B_{n+1}$
ist Untermodul des endlich erzeugten $\mathds{Z}$-Moduls $C_n$ und somit endlich erzeugt.
Wir erhalten kurze exakte Folgen:
\[
	\begin{tikzcd}[row sep=0.4em]
		B_{n+1} \rar[hook] & Z_n \rar[two heads] & H_n(C_*, d_*) \\
		Z_n \rar[hook] & C_n \rar[two heads] & B_n \quad 
	\end{tikzcd}
\]
Folglich gilt $\Rg Z_n = \Rg B_{n+1} + \Rg H_n(C_*,d_*)$ und $\Rg C_n = \Rg Z_n + \Rg B_n$. Also gilt
\begin{align*}
	\chi(C_*) = \sum_{i=0}^{\infty} (-1)^i \Rg C_i &= \sum_{i=0}^{\infty} (-1)^i \enbrace[\big]{\Rg Z_i + \Rg B_i} \\  
	&= \sum_{i=0}^{\infty} (-1)^i \enbrace[\Big]{\Rg H_i(C_*,d_*) + \Rg B_{i+1} + \Rg B_i} \\
	&= \sum_{i=0}^{\infty} (-1)^i \Rg H_i(C_*,d_*) + \sum_{i=0}^{\infty} (-1)^i \Rg B_{i+1} + \sum_{i=0}^{\infty}(-1)^i \Rg B_i \\
	&= \sum_{i=0}^{\infty} (-1)^i \Rg H_i(C_*,d_*)  \bewende  
\end{align*}
% subsection 410 (end)

\subsection[Korollar: Die Eulercharakteristik eines simplizialen Komplexes]{Korollar} % (fold)
\label{sub:411}
Sei $K$ ein endlicher, geordneter simplizialer Kettenkomplex. Dann gilt
\[
	\chi(K) = \sum_{i=1}^{\dim K} (-1)^i \Rg H_i(K;\mathds{Z})
\]
\minisec{Beweis}
Wende den Satz \ref{sub:410} auf den simplizialen Kettenkomplex von $K$ an. Da 
\[
	\Rg C_n(K;\mathds{Z}) = \Rg \mathds{Z}[\Sigma_n] = \# \text{ $n$-Simplizes in $K$}
\]
ist
\[
	\chi(K) = \chi \enbrace[\big]{C_*(K;\mathds{Z}), \partial_*} = \sum_{i=0}^{\infty} (-1)^i \Rg H_i \enbrace[\big]{C_*(K;\mathds{Z}), \partial_*}  =
	\sum_{i=0}^{\infty} (-1)^i \Rg H_i \enbrace*{K;\mathds{Z}} \bewende
\]
% subsection 411 (end)
% section 4 (end)
\newpage

\section{Singuläre Homologie} % (fold)
\label{sec:5}

\subsection[Definition: Singuläre Simplizes und $n$-ter singulärer Kettenmodul]{Definition} % (fold)
\label{sub:51}
Sei $X$ ein topologischer Raum. Sei $S_n(X)$ die Menge aller stetigen Abbildungen $\sigma : \abs*{\Delta^n} \to X$. Elemente von $S_n(X)$ heißen \Index{singuläre Simplizes} 
in $X$. Sei $R$ ein Ring. Der \bet{$\mathbf{n}$-te (singuläre) Kettenmodul}\index{n-ter Kettenmodul@$n$-ter Kettenmodul!singulär} von $X$ über $R$ ist
\[
	C_n(X;R) := R \benbrace*{S_n(X)} 
\] 

% subsection 51 (end)

\subsection[Definition: Einschränkung eines singulären Simplizes auf eine Seite]{Definition} % (fold)
\label{sub:52}
Für $\sigma \in S_n(X)$ und $j \in \set{0,\ldots ,n}$ sei $\delta_j \sigma \in S_{n-1}(X)$ die Einschränkung von $\sigma$ auf die $j$-te Seite von $\abs*{\Delta^n}$,
also $\delta_j(\sigma) = \sigma \circ \iota_{n,j}$, wobei $\iota_{n,j} : \abs*{\Delta^{n-1}} \to \abs*{\Delta^n}$ die Inklusion der $j$-ten Seite ist:

\[
	\iota_{n,j}(x_0, \ldots ,x_{n-1}) = \enbrace[\big]{x_0, \ldots , x_{j-1}, 0, x_j, \ldots , x_{n-1}}
\]
% subsection 52 (end)

\subsection[Bemerkung zur Inklusion der $j$-ten Seite]{Bemerkung} % (fold)
\label{sub:53}
Es gilt $\iota_{n,j} = \abs*{i_{n,j}}$ wobei $i_{n,j} : \Delta^{n-1} \to \Delta^n$ gegeben ist durch \marginnote{vergleiche Blatt 2, Aufgabe 2}
\[
	i_{n,j}(k) = \begin{cases}
	k, &\text{ falls }k < j\\
	k+1, &\text{ falls } k\ge j
\end{cases}
\]
% subsection 53 (end)

\subsection[Definition: $n$-te singuläre Randabbildung]{Definition} % (fold)
\label{sub:54}
Die \bet{$\mathbf{n}$-te singuläre Randabbildung}\index{n-te Randabbildung@$n$-te Randabbildung!singulär} $\partial_n : C_n(X;R) \to C_{n-1}(X;R)$ ist 
definiert durch 
\[
	\partial_n(\sigma) = \sum_{j=0}^{n} (-1)^j \cdot \delta_j(\sigma).
\]
% subsection 54 (end)

\subsection[Proposition: Für die Randabbildungen gilt $\partial_{n-1} \circ \partial_n = 0$]{Proposition} % (fold)
\label{sub:55}
\[
	\partial_{n-1} \circ \partial_n = 0
\]
\minisec{Beweis (mit Lemma \ref{sub:56})}
\begin{align*}
	\partial_{n-1} \partial_n(\sigma) = \partial_{n-1} \enbrace*{\sum_{j=0}^{n} (-1)^j \delta_j(\sigma)} &= 
	\sum_{k=0}^{n-1} \sum_{j=0}^{n} (-1)^{j+k} \delta_k \delta_j(\sigma)   \\
	&= \,\sum_{0 \le k < j \le n} \!\!\! (-1)^{k+j} \delta_k \delta_j(\sigma) +\!\!\!\smashoperator[r]{\sum_{0 \le j \le k \le n-1}} \enspace(-1)^{k+j} \delta_k\delta_j(\sigma) \\
	&\stackrel{\mathclap{\text{Lemma \ref{sub:56}}}}{=} \, \sum_{0 \le k < j \le n} \!\!\!(-1)^{k+j} \delta_k \delta_j(\sigma) + \!\!\!\smashoperator[r]{\sum_{0 \le j \le k \le n-1}} \enspace(-1)^{k+j} \delta_j \delta_{k+1}(\sigma) \\
	&= \,\sum_{0 \le k < j \le n} \!\!\! (-1)^{k+j} \delta_k \delta_j(\sigma) + \smashoperator[r]{\sum_{0 \le j < k \le n}} \enspace(-1)^{k-1+j} \delta_j \delta_k(\sigma) \\
	&= 0 \bewende
\end{align*}
% subsection 55 (end)

\subsection[Lemma: Hilfslemma für Proposition \ref{sub:55}]{Lemma} % (fold)
\label{sub:56}
Für $0 \le j \le k \le n-1$ und $\sigma \in S_n(X)$ ist $\delta_k \delta_j(\sigma) = \delta_j \delta_{k+1}(\sigma)$.
\minisec{Beweis}
Es ist
\begin{align*}
	\delta_k \enbrace[\big]{\delta_j(\sigma)} \enbrace*{x_0, \ldots , x_{n-2}} &= \delta_j(\sigma) \enbrace[\big]{x_0, \ldots , x_{k-1}, 0, x_k , \ldots ,x_{n-2}}  \\
	&= \enbrace[\big]{x_0, \ldots , x_{j-1},0,x_j, \ldots , x_{k-1},0,x_k, \ldots, x_{n-2}}  \\
	\intertext{und}
	\delta_j \delta_{k+1}(\sigma) \enbrace*{x_0, \ldots , x_{n-2}} &= \delta_{k+1}(\sigma) \enbrace[\big]{x_0, \ldots , x_{j-1},0,x_j,\ldots , x_{n-2}} \\
	&= \enbrace[\big]{x_0, \ldots , x_{j-1},0,x_j, \ldots , x_{k-1},0, x_k, \ldots , x_{n-2}} \bewende
\end{align*}
% subsection 56 (end)

\subsection[Definition: Singuläre Homologie von $X$]{Definition} % (fold)
\label{sub:57}
Sei $X$ ein topologischer Raum. Die Homologie des \bet{singulären Kettenkomplex}\index{singulärer Kettenkomplex} über $R$, $\enbrace[\big]{C_*(X;R), \partial_*}$, heißt die
\Index{singuläre Homologie von $X$} mit Koeffizienten in $R$:
\[
	H_n(X;R) := H_n \enbrace[\big]{C_*(X;R), \partial_*} 
\]  
Für $R=\mathds{Z}$ schreiben wir kürzer $C_*(X) := C_*(X;\mathds{Z})$ und $H_n(X) := H_n(X;\mathds{Z})$.
% subsection 57 (end)

\subsection[Definition: $n$-Ketten, $n$-Ränder und $n$-Zykel und Homologieklasse]{Definition} % (fold)
\label{sub:58}
\begin{itemize}
	\item Die Elemente von $C_n(X;R)$ heißen \bet{$\mathbf{n}$-Ketten}\index{n-Ketten@$n$-Ketten}.
	\item Die Elemente von $\im \partial_{n+1} \subseteq C_n(X;R)$ heißen \bet{$\mathbf{n}$-Ränder}\index{n-Ränder@$n$-Ränder}.
	\item Die Elemente von $\ker \partial_{n} \subseteq C_n(X;R)$ heißen \bet{$\mathbf{n}$-Zykel}\index{n-Zykel@$n$-Zykel}.
\end{itemize}
Jeder $n$-Zykel $\sigma \in C_n(X;R)$ bestimmt eine \Index{Homologieklasse} $[\sigma] := \sigma + \im \partial_{n+1} \in H_n(X;R)$. 
% subsection 58 (end)

\subsection[Beispiel: Die Homologie des Ein-Punkt-Raumes]{Beispiel} % (fold)
\label{sub:59}
Ist $X= \set{x_0}$ der Ein-Punkt-Raum, so ist 
\[
	H_* \enbrace[\big]{\set{x_0};R} \cong \begin{cases}
	R, &\text{ falls }*=0\\
	0 , &\text{ sonst } 
\end{cases} 
\]
\minisec{Beweis}
Es gibt für jedes $n$ genau eine Abbildung, nämlich die konstante, $\sigma_n : \abs*{\Delta^n} \to \set{x_0}$ und diese ist stetig. Also 
$C_n \enbrace*{\set{x_0};R}=R[\sigma_n]$. Für alle $j$ ist $\delta_j \sigma_n = \sigma_{n-1}$. Daher ist
\[
	\partial_n(\sigma_n) = \begin{cases}
		0, &\text{ falls }n \text{ ungerade}\\
		\sigma_{n-1} , &\text{ falls } n \text{ gerade}
	\end{cases}
\]
Der singuläre Kettenkomplex von $\set{x_0}$ hat also folgende Gestalt:
\[
	\begin{tikzcd}[row sep=0.4em]
		C_0 \enbrace[\big]{\set{x_0};R} & C_1 \enbrace[\big]{\set{x_0};R} \lar["0", swap] & 
		C_2 \enbrace[\big]{\set{x_0};R} \lar["\cong", swap] & C_3 \enbrace[\big]{\set{x_0};R} \lar["0", swap]& \cdots \lar["\cong",swap] \\
		R & R \lar["0"] & R \lar["\id"] & R \lar["0"] & R \lar["\id"]
	\end{tikzcd}
\]
Es folgt
\[
	H_n \enbrace[\big]{\set{x_0};R } \cong \begin{cases}
		R, &\text{ falls }n=0\\
		0 , &\text{ sonst}
	\end{cases} \bewende
\]
% subsection 59 (end)

\subsection[Proposition: Eigenschaften von $H_0(X)$ für $X \not= \emptyset$ und $X$ wegzusammenhängend]{Proposition} % (fold)
\label{sub:510}
\begin{enumerate}[1)]
	\item Ist $X \not= \emptyset$, so ist $H_0(X) \not=0$.
	\item Ist $X$ wegzusammenhängend, so gilt $H_0(X) \cong \mathds{Z}$
\end{enumerate}
\minisec{Beweis}
\begin{enumerate}[1)]
	\item Sei $\varepsilon : C_0(X) \to \mathds{Z}$ definiert durch $\varepsilon \enbrace*{\sum_{\sigma \in S_0(X)}r_\sigma \cdot \sigma } := \sum_\sigma r_\sigma$.
	Dann ist $\varepsilon \circ \partial_1 : C_1(X) \to \mathds{Z}$ trivial, denn für $\sigma \in S_1(X)$ ist 
	$\varepsilon \circ \partial_1(\sigma)= \varepsilon(\delta_0 \sigma - \delta_1 \sigma) = 0$.
	Daher induziert $\varepsilon$ eine Abbildung 
	$\overline{\varepsilon} : H_0(X) \to \mathds{Z}$. Da $X \not= \emptyset$ ist, gibt es einen singulären $0$-Simplex $\sigma : \abs*{\Delta^0} \to X$. Für $\sigma$ gilt 
	$\overline{\varepsilon}( [\sigma])=1$ und daher ist $\overline{\varepsilon}$ surjektiv und $H_0(X) \not= 0$.
	\item Wir zeigen: $\overline{\varepsilon} : H_0(X) \to \mathds{Z}$ ist ein Isomorphismus, falls $X$ wegzusammenhängend ist. Dazu zeigen wir 
	$\ker \varepsilon = \im \partial_1$. Dazu definieren wir $s : C_0(X) \to C_1(X)$ wie folgt: Sei $\sigma_0$ ein fest gewählter singulärer $0$-Simplex.
	Zu $\sigma \in S_0(X)$ gibt es, da $X$ wegzusammenhängend ist, $s(\sigma) \in S_1(X)$ mit $\delta_0 s(\sigma)= \sigma$, $\delta_1 s(\sigma)= \sigma_0$.
	Es gilt nun 
	\begin{align*}
		\partial_1 \circ s(\sigma) = \sigma - \sigma_0 = \enbrace*{\id_{C_0(X)} - i \circ \varepsilon }(\sigma)
	\end{align*}
	wobei $i : \mathds{Z} \to C_0(X)$, $n \mapsto n \cdot \sigma_0$. Ist nun $v \in \ker \varepsilon$, so folgt 
	\[
		\partial_1 \circ s(v)= v- i \circ \underbrace{\varepsilon(v)}_{=0} = v
	\]
	also $v \in \im \partial_1$. \bewende
	
	% Sei $\sigma_0 : \abs*{\Delta^0} \to X$ ein fest gewählter singulärer $0$-Simplex (also ein Punkt in X). Definiere $i : \mathds{Z} \to C_0(X)$ durch
% 	$i(n) := n \cdot \sigma$ und $\overline{i} : \mathds{Z} \to H_0(X)$ durch $\overline{i}(n) = \benbrace*{i(n)}$. Dann ist
% 	$\overline{\varepsilon}\circ\overline{i} = \id_\mathds{Z}$. Wir behaupten, dass auch $\overline{i} \circ \overline{\varepsilon}= \id_{H_0(X)}$. Dazu definieren
% 	wir $s : C_0(X) \to C_1(X)$ wie folgt:
%
% 	Zu jedem $\sigma \in C_0(X)$ wählen wir $s(\sigma) : \abs*{\Delta^1} \to X$ so, dass $\delta_0 s(\sigma) = \sigma$ und $\delta_1 s(\sigma)= \sigma_0$.
% 	Dies ist möglich, da $X$ wegzusammenhängend ist: $s(\sigma)$ sei ein Weg von $\im \sigma_0$ nach $\im \sigma$. Es gilt nun
% 	\[
% 		\partial \circ s (\sigma) = \sigma - \sigma_0 = \sigma- i \circ \varepsilon(\sigma)
% 	\]
% 	für jedes $\sigma \in S_0(X)$. Damit folgt $\partial \circ s = \id_{C_0(X)} - i \circ \varepsilon$. Also ist
% 	\[
% 		\overline{i} \circ \overline{\varepsilon} (\sigma) = \benbrace*{i \circ \varepsilon(\sigma)} = \benbrace[\big]{\sigma - \partial \circ s(\sigma)} = [\sigma]
% 	\]
% 	für jedes $\sigma \in C_0(X)$. Also $\overline{\imath} \circ \overline{\varepsilon} = \id_{H_0(X)}$. \bewende
\end{enumerate}
% subsection 510 (end)

\subsection[Bemerkung: Vorgehen, um später zu zeigen: $H_*(K;R) \cong H_*(\abs*{K};R )$]{Bemerkung} % (fold)
\label{sub:511}
Sei $K=(V,\Sigma)$ ein geordneter simplizialer Komplex. Sei $\sigma= \set{v_o, \ldots , v_n} \in \Sigma_n $ ein $n$-Simplex von $K$ mit $v_0 < v_1 < \ldots < v_n$. Wir
ordnen $\sigma$ den singulären $n$-Simplex $f_\sigma : \abs*{\Delta^n} \to \abs*{\Delta^\sigma} \subseteq \abs*{K}$ zu, wobei
\begin{align*}
	f_\sigma \enbrace*{\sum_{i=0}^{n} t_i e_i} = \sum_{i=0}^{n} t_i v_i. 
\end{align*}
Nun erhalten wir eine Abbildung 
\[
	C_*(K;R) \to C_*(\abs*{K},R) \quad , \qquad   \sum_{\sigma \in \Sigma_n} r_\sigma \cdot \sigma \longmapsto \sum_{\sigma \in \Sigma_n} r_\sigma \cdot f_\sigma.
\]
Diese Abbildung induziert(!) eine weitere Abbildung $H_*(K;R) \xrightarrow{\cong} H_*(\abs*{K};R )$ von der wir später zeigen werden, dass sie ein Isomorphismus ist.
% subsection 511 (end)
% section 5 (end)
\newpage

\section{Funktorialität} % (fold)
\label{sec:6}

\subsection[Definition: $R$-Kettenabbildung]{Definition} % (fold)
\label{sub:61}
Seien $(C_*,d_*)$ und $(C_*',d_*')$ zwei $R$-Kettenkomplexe. Eine \bet{$\mathbf{R}$-Kettenabbildung}\index{R-Kettenabbildung@$R$-Kettenabbildung}
$f_* : (C_*,d_*) \to (C_*',d_*')$ ist eine Folge von $R$-linearen Abbildungen $f_n : C_n \to C_n'$, sodass
\[
	d_n' \circ f_n = f_{n-1} \circ d_n
\]
für alle $n\ge 1$ ist. Also kommutiert folgendes Diagramm\todo{RevChap6}
\[
	\begin{tikzcd}
		C_0 \dar["f_0", Firebrick3] & C_1 \lar["d_1"] \dar["f_1", Firebrick3] & C_2 \lar["d_2"] \dar["f_2", Firebrick3]& C_3 \lar["d_3"] \dar["f_3", Firebrick3]& \ldots \lar \\
		C_0' & C_1' \lar["d_1'"]  & C_2' \lar["d_2'"]& C_3' \lar["d_3'"]& \ldots \lar
	\end{tikzcd}
\]

% subsection 61 (end)

\subsection[Bemerkung: Induzierte Abbildung einer Kettenabbildung]{Bemerkung} % (fold)
\label{sub:62}
Ist $f_* : (C_*,d_*) \to (C_*',d_*')$ eine Kettenabbildung, so erhalten wir eine induzierte Abbildung $H_n(f_*) : H_n(C_*,d_*) \to H_n(C_*',d_*')$ durch
\[
	H_n(f_*) \enbrace*{[v]} := \benbrace*{f_n(v)}  
\]
für $v \in \ker d_n$, denn:
\begin{itemize}
	\item $d_n' \enbrace[\big]{f_n(v)} = f_{n-1} \enbrace[\big]{d_n(v)} = 0 $, also $f_n(v) \in \ker d_n'$ für $v \in \ker d_n$.
	\item Ist $[v]=[w]$, mit $v,w \in\ker d_n$, so gibt es $x \in C_{n+1}$ mit $d_{n+1}(x)=v-w$. Dann ist 
	\[
		d_{n+1}'\enbrace[\big]{f_{n+1}(x)} = f_n \enbrace[\big]{d_{n+1}(x)} = f_n(v)- f_n(w) 
	\]
	also $\benbrace*{f_n(v)} = \benbrace*{f_n(w)} \in H_n(C_*',d_*')$.
\end{itemize}
% subsection 62 (end)

\subsection[Bemerkung: Homologie definiert einen Funktor]{Bemerkung} % (fold)
\label{sub:63}
\begin{enumerate}[a)]
	\item $R$-Kettenkomplexe mit Kettenabbildungen bilden die Kategorie $R\text{-}\textsc{Ketten}$.
	\item Homologie definiert nun einen Funktor: $H_n : R\text{-}\textsc{Ketten} \to R\text{-}\MOD$.
	
	Es gilt $H_n(\id_{(C_*,d_*)}) = \id_{H_n(C_*,d_*)}$ und $H_n(f \circ g) = H_n(f) \circ H_n(g)$, da
	\[
		H_n(f \circ g) \enbrace*{[v]} = \benbrace*{f_n \circ g_n(v)} = \benbrace*{f_n \enbrace[\big]{g_n(v)} } = H_n(f) \benbrace*{g_n(v)} = H_n(f) \enbrace[\big]{H_n(g)(v)}     
	\]
	\item Definiert man die Kategorie der \Index{graduierten} $R$-Moduln als die Kategorie, deren Objekte Folgen $(V_n)_n$ von $\mathds{R}$-Moduln sind und deren
	Morphismen Folgen von $R$-linearen Abbildungen $(f_n)_{n}$ sind, so kann man die $H_n$, $n \in \mathds{N}$ zu einem Funktor
	\[
		H_* : R\text{-}\textsc{Ketten} \longrightarrow \textsc{Gr}\text{-}R\text{-}\MOD \quad , \qquad (C_*,d_*) \longmapsto \enbrace[\big]{H_n(C_*,d_*)}_n 
	\]
	zusammensetzen.
\end{enumerate}
% subsection 63 (end)

\subsection[Definition: Induzierte Abbildung auf Kettenkomplexen]{Definition} % (fold)
\label{sub:64}
Sei $f : X \to Y$ stetig. Wir definieren $C_n(f;R) : C_n(X;R) \to C_n(Y;R)$ durch
\[
	C_n(f;R) \enbrace*{\smashoperator[r]{\sum_{\sigma \in S_n(X)}} r_\sigma \cdot  \sigma} := \smashoperator[r]{\sum_{\sigma \in S_n(X)}} r_\sigma \cdot \enbrace*{f \circ \sigma} 
\]
(Für $\sigma \in S_n(X)$ ist $f \circ \sigma \in S_n(Y)$.)
% subsection 64 (end)

\subsection[Proposition: Der singuläre Kettenkomplex bildet einen Funktor]{Proposition} % (fold)
\label{sub:65}
Mit dieser Definition von $C_*(f;R)$ wird der singuläre Kettenkomplex über $R$ zu einem Funktor
\[
	C_*(-;R) : \textsc{Top} \to R\text{-}\textsc{Ketten}
\]
\minisec{Beweis}
\begin{enumerate}[(i)]
	\item $C_*(f;R)$ ist eine $R$-Kettenabbildung:
	% \begin{align*}
	% 	C_{n-1}(f;R) \circ \partial_n \enbrace*{\sum_{\sigma \in S_n(X)} r_\sigma \cdot  \sigma} &= C_{n-1}(f;R) \enbrace*{\sum_{\sigma \in S_n(X)} \sum_{j=0}^{n} (-1)^j
	% 	\cdot r_\sigma (\delta_j \sigma)} \\
	% 	&= \sum_{\sigma \in S_n(X)} \sum_{j=0}^{n} (-1)^j \cdot r_\sigma \enbrace[\big]{f_{n-1} \circ (\delta_j \sigma)} \\
	% 	&= \sum_{\sigma \in S_n(X)} \sum_{j=0}^{n} (-1)^j \cdot r_\sigma \enbrace[\big]{f_{n-1} \circ (\sigma \circ \iota_{n,j})} \\
	% 	&= \sum_{\sigma \in S_n(X)} \sum_{j=0}^{n} (-1)^j \cdot r_\sigma \enbrace[\big]{(f_{n-1} \circ \sigma) \circ \iota_{n,j}} \\
	% 	&= \sum_{j=0}^{n} (-1)^j \enbrace[\Bigg]{\smashoperator[r]{\sum_{\sigma \in S_n(X)}} r_\sigma \underbrace{\enbrace*{f_{n-1} \circ \sigma} \circ \iota_{n,j}}_{\delta_j (f \circ \sigma)} } \\
	% 	&= \partial_n \circ C_n \enbrace*{f;R}(\sigma)
	% \end{align*}
	\begin{align*}
		C_{n-1}(F;R) \circ \partial_n(\sigma) = C_{n-1} (f;R) \enbrace*{\sum_{l=0}^{n}  (-1)^l \sigma \circ \iota_{n,l}} &=
		\sum_{l=0}^{n} (-1)^l (f \circ \sigma) \circ \iota_{n,l} \\
		&= \partial_n(f \circ \sigma) = \partial \enbrace*{C_n(f;R)(\sigma)} 
	\end{align*}
	\item Zu zeigen: $C_*(f \circ g;R) = C_*(f;R) \circ C_*(g;R)$:
	\begin{align*}
		C_n(f \circ g;R)(\sigma) = f \circ (g \circ \sigma) = C_n(f;R)(g \circ \sigma) = C_n(f;R) \enbrace[\big]{C_n(g;R)(\sigma)} 
	\end{align*}
	\item Zu zeigen: $C_*(\id_X;R) = \id_{C_*(X;R)}$:
	\[
		C_n(\id_X;R)(\sigma) = \id_X \circ \sigma = \sigma
	\]
\end{enumerate}
% subsection 65 (end)

\subsection[Korollar: Homologie ist ein Funktor $\TOP \to R\text{-}\textsc{Mod}$]{Korollar} % (fold)
\label{sub:66}
Mit $H_n(f;R) := H_n \enbrace*{C_*(f;R)}$ wird $H_n(-;R)$ zu einem Funktor: $H_n(-;R) : \TOP \to R\text{-}\textsc{Mod}$.
\minisec{Beweis}
$H_n(-;R)$ ist die Komposition der Funktoren $C_*(-;R) : \TOP \to R\text{-}\textsc{Ketten}$ und $H_* : R\text{-}\textsc{Ketten} \to R\text{-}\textsc{Mod}$. \bewende
% subsection 66 (end)

\subsection[Bemerkung: Notation und Formel für $H_n(f;R)$]{Bemerkung} % (fold)
\label{sub:67}
\begin{itemize}
	\item Oft schreiben wir $f_* = H_n(f;R)$.
	\item Es gilt für $\benbrace*{\sum_{\sigma \in S_n(X)} r_\sigma \cdot \sigma} \in H_n(X;R)$
	\[
		f_* \enbrace*{\benbrace*{\smashoperator[r]{\sum_{\sigma \in S_n(X)}} r_\sigma \cdot  \sigma}} =  \benbrace*{\smashoperator[r]{\sum_{\sigma \in S_n(X)}} r_\sigma
		\cdot f\circ \sigma} \in H_n(Y;R)
	\]
\end{itemize}
% subsection 67 (end)

\subsection[Definition: Summe oder Koprodukt von topologischen Räumen]{Definition} % (fold)
\label{sub:68}
Seien $X_i$, $i \in I$ topologische Räume. Mit 
\[
	X := \coprod_{i \in I} X_i 
\]
bezeichnen wir die \Index{Summe} (oder auch das \Index{Koprodukt}) der $X_i$. Als Menge
ist $X$ die disjunkte Vereinigung der $X_i$. $U \subseteq X$ ist offen genau dann, wenn $X_i \cap U \subseteq X_i$ offen ist für jedes $i \in I$. Für jedes $i_o$ erhalten
wir eine stetige Inklusion $j_{i_0} : X_{i_0} \to \coprod_{i} X_i$.
\minisec{Bemerkung}
Für jedes $i_o$ mit $X_{i_0} \subseteq \coprod_i X_i$ offen und abgeschlossen.
% subsection 68 (end)

\subsection[Definition: Summe von $R$-Moduln]{Definition} % (fold)
\label{sub:69}
Seien $V_i$, $i \in I$ $R$-Moduln. Mit 
\[
	V := \bigoplus_{i \in I} V_i
\]
bezeichnen wir die \Index{Summe} (oder auch das \Index{Koprodukt}) der $V_i$. Elemente von $V$ sind $I$-Folgen $(v_i)_{i \in I}$ mit $v_i \in V_i$ und $v_i=0$ für
alle bis auf endlich viele $i$. Die $R$-Modulstruktur ist erklärt durch:
\[
	(v_i)_{i \in I} + (w_i)_{i \in I} := (v_i +w_i)_{i \in I} \qquad ,\qquad r \cdot (v_i)_{i \in I} := (r \cdot v_i)_{i \in I}
\]
Für jedes $i_0 \in I$ erhalten wir eine $R$-lineare Abbildung $j_{i_0} : V_{i_0} \to V$ mit
\[
	\enbrace*{j_{i_0} (v)}_{i} = \begin{cases}
		v, &\text{ falls }i=i_0\\
		0 ,&\text{ sonst}
	\end{cases} 
\]
\minisec{Bemerkung}
Seien $V_i$, $i \in I$ $R$-Moduln. Sei $W$ ein weiterer $R$-Modul. Dann gibt es zu jeder Folge $f_i : V_i \to W$ von $R$-linearen Abbildungen eine
$R$-lineare Abbildung 
\(
	\oplus_{i \in I} f_i : \bigoplus_{i \in I} V_i \to W  
\)
mit
\[
	\enbrace*{\oplus_{i \in I} f_i} \enbrace[\big]{(v_i)_{i \in I}} = \sum_{i \in I} f_i(v_i)  
\]
Ist umgekehrt $f : \bigoplus_{i \in I} V_i \to W$ eine $R$-lineare Abbildung, so ist $f_i := f \circ j_{i}$ eine Folge von $R$-linearen Abbildungen mit 
$f = \bigoplus_{i \in I} f_i$
% subsection 69 (end)

\subsection[Satz: Homologie des Koproduktes topologischer Räume]{Satz} % (fold)
\label{sub:610}
Sei $X = \coprod_{i \in I} X_i$. Dann induzieren die $j_i : X_i \to X$ einen Isomorphismus
\[
	\bigoplus_{i \in I} H_n(X_i;R) \xrightarrow{\bigoplus_{i \in I} (j_i)_*}  H_n(X;R)
\]
\minisec{Beweis}
Da die $X_i \subseteq X$ offen und abgeschlossen sind und $\abs*{\Delta^n}$ zusammenhängend ist (sogar wegzusammenhängend), gibt es für jedes 
$\sigma : \abs*{\Delta^n} \to X $ ein eindeutiges $i$ mit $\im \sigma \subseteq  X_i$. Es gilt also $S_n(X) = \dot{\bigcup} S_n(X_i)$. Daher induzieren die $j_i$
einen Isomorphismus
\[
	\oplus_{i \in I} C_n(j_i;R) : \bigoplus_{i \in I}C_n(X_i;R) \xrightarrow{\enspace\cong \enspace} C_n(X;R) 
\]
Da diese Isomorphismen mit den Randabbildungen vertauschen, erhalten wir einen Isomorphismus von $R$-Kettenkomplexen
\[
	\varphi := \oplus_{i \in I} C_*(j_i;R) : \bigoplus C_*(X_i;R) \to C_*(X;R)
\]
Da $\bigoplus H_n \enbrace*{C_*(X_i;R),\partial_*} \cong H_n \enbrace*{\bigoplus_i C_*(X_i;R), \partial_*}$ induzieren dann auch die $(j_i)_*$ einen Isomorphismus
\[
	\bigoplus (j_i)_* : \bigoplus_{i \in I} H_n(X_i;R) \longrightarrow H_n(X;R) \bewende
\]
% subsection 610 (end)
% section 6 (end)
\newpage

\section{Homotopieinvarianz} % (fold)
\label{sec:7}

\subsection[Bemerkung: Die induzierte Abbildung eines Homöomorphismus ist ein Isomorphismus]{Bemerkung} % (fold)
\label{sub:71}
Sei $f : X \to Y$ eine Homöomorphismus. Dann ist 
\[
	f_* : H_n(X;R) \longrightarrow H_n(Y;R)
\]
ein Isomorphismus.
\minisec{Beweis}
Da Homologie ein Funktor ist, gilt
\begin{align*}
	f_* \circ (f^{-1})_* &= (f \circ f ^{-1})_* = (\id_Y)_* = \id_{H_n(Y;R)} \qquad \text{ und} \\
	(f ^{-1})_* \circ f_* &= (f ^{-1} \circ f)_* = (\id_X)_* = \id_{H_n(X;R)}
\end{align*}
Also ist $(f_*)^{-1}=(f ^{-1})_*$. Insbesondere ist $f_*$ ein Isomorphismus. \bewende
% subsection 71 (end)

\subsection[Definition: Homotopieäquivalenz und Homotopieinverse]{Definition} % (fold)
\label{sub:72}
Eine stetige Abbildung $f : X \to Y$ heißt eine \Index{Homotopieäquivalenz}, falls es eine stetige Abbildung $g : Y \to X$ gibt, so dass $g \circ f$ homotop zu $\id_X$ ist
und $f \circ g$ homotop zu $\id_Y$ ist. $g$ heißt dann eine \Index{Homotopieinverse} zu $f$.
\minisec{Bemerkung}
\begin{enumerate}[(i)]
	\item $f$ ist genau dann eine Homotopieäquivalenz, wenn $[f]$ in $\HTOP$ invertierbar ist. In $\HTOP$ ist dann $[f]^{-1} = [g]$.
	\item Die Homotopieinverse ist eindeutig bis auf Homotopie.
\end{enumerate}
% subsection 72 (end)

\subsection[Satz: Homotopieinvarianz der Homologie]{Satz} % (fold)
\label{sub:73}
Seien $f,g : X \to Y$ stetige Abbildungen. Sind $f$ und $g$ homotop, so gilt $H_n(f;R) = H_n(g;R)$ für alle $n \in \mathds{N}$.
\minisec{Beweis}
\emph{Siehe \ref{sub:78}, sowie \ref{sub:713} } \bewende

% subsection 73 (end)

\subsection[Korollar: Die induzierte Abbildung einer Homotopieäquivalenz ist ein Isomorphismus]{Korollar} % (fold)
\label{sub:74}
Ist $f : X \to Y$ eine Homotopieäquivalenz, so ist $f_* : H_n(X;R) \to H_n(Y;R)$ ein Isomorphismus.
\minisec{Beweis}
Sei $g : Y \to X$ ein Homotopieinverses zu $f$. Es folgt 
\begin{align*}
	f_* \circ g_* &\stackrel{\text{Funktor}}{=} (f \circ g)_* \stackrel{\text{Homotopieinvarianz}}{=} (\id_Y)_* \stackrel{\text{Funktor}}{=} \id_{H_n(Y;R)} \\
	g_* \circ f_* &\stackrel{\text{Funktor}}{=} (g \circ f)_* \stackrel{\text{Homotopieinvarianz}}{=} (\id_X)_* \stackrel{\text{Funktor}}{=} \id_{H_n(X;R)}
\end{align*}
% subsection 74 (end)

\subsection[Korollar: Homologie eines kontrahierbaren Raumes]{Korollar} % (fold)
\label{sub:75}
Ist $X$ \Index{kontrahierbar}, d.h. es gibt eine Homotopie $H : X \times [0,1] \to X$ mit $H_0 =\id_X$ und $H_1$ konstant, so gilt 
\[
	H_n(X;R) \cong \begin{cases}
		R, &\text{ falls }n=0\\
		0, &\text{ sonst}
	\end{cases}
\] 
\minisec{Beweis}
Sei $\set{x_0} =\im H_1$. Dann ist die Inklusion $\set{x_0} \to X$ eine Homotopieäquivalenz. Also gilt nach \ref{sub:74}
\[
	H_n(X;R) \cong H_n \enbrace[\big]{\set{x_0};R} \stackrel{\text{\ref{sub:59}}}{=} \begin{cases}
		R, &\text{ falls }n=0\\
		0, &\text{ sonst}
	\end{cases} \bewende
\]
\minisec{Beispiel}
\[
	H_n(\mathds{R}^k;R) \cong \begin{cases}
		R, &\text{ falls }n=0\\
		0, &\text{ sonst}
	\end{cases} \qquad \text{ und } \qquad 
	H_n(D^k;R) \cong \begin{cases}
		R, &\text{ falls }n=0\\
		0, &\text{ sonst}
	\end{cases}
\]
% subsection 75 (end)

\subsection[Definition: Kettenhomotopie]{Definition} % (fold)
\label{sub:76}
Seien $(C_*,d_*)$ und $(C_*',d_*')$ $R$-Kettenkomplexe. Seien $f_*,g_* : (C_*,d_*) \to (C_*',d_*')$ $R$-Kettenabbildungen. Eine \Index{Kettenhomotopie} von $f_*$ nach $g_*$
ist eine Folge von $R$-linearen Abbildungen $h_n :C_n \to C_{n+1}'$, $n \in \mathds{N}$, sodass für alle $n$ gilt:
\[
	d_{n+1}' \circ h_n + h_{n-1} \circ d_n = f_n - g_n
\] 
In diesem fall heißen $f_*$ und $g_*$ \Index{kettenhomotop}.
% subsection 76 (end)

\subsection[Prop.: induzierte Abbildungen von kettenhomotopen Kettenabbildungen sind gleich]{Proposition} % (fold)
\label{sub:77}
Seien $f_*,g_* : (C_*,d_*) \to (C_*',d_*')$ Kettenabbildungen. Sind $f_*$ und $g_*$ kettenhomotop, so gilt für alle $n$
\[
	H_n(f_*) = H_n(g_*)
\]
\minisec{Beweis}
Sei $x \in H_n(C_*,d_*)$. Also $x=[v]$ mit $v \in \ker d_n : C_n \to C_{n-1}$. Dann gilt, da $v \in \ker d_n$
\begin{align*}
	H_n(f_*)(x) = H_n(f_*) \enbrace*{[v]} = \benbrace*{f_n(v)}  &= \benbrace[\Big]{g_n(v) +\underbrace{ d_{n+1}' \circ h_n(v)}_{\in \im d_{n+1}'} + \underbrace{h_{n-1} \circ d_n(v)}_{=0}} \\
	&= \benbrace*{g_n(v)} = H_n(g_*) \enbrace*{[v]} = H_n(g_*)(x) \bewende
\end{align*}
% subsection 77 (end)

\subsection[Lemma: Reduktion von Satz \ref{sub:73} auf Beweis eines Spezialfalles]{Lemma} % (fold)
\label{sub:78}
Die Homotopieinvarianz von $H_n(-;R)$ folgt aus folgenden Spezialfall:
\begin{equation*}
	\text{Seien $i_0,i_1  :X \hookrightarrow X \times [0,1]$ Inklusionen mit $i_0(x)= (x,0)$, $i_1(x)=(x,1)$. $\Longrightarrow$ $(i_0)_* = (i_1)_*$} \label{eq:78} \tag{\#}
\end{equation*}
\minisec{Beweis}
Sei $H : X \times [0,1] \to Y$ eine Homotopie zwischen $f,g : X \to Y$, also $f = H \circ i_0$ und $g = H \circ i_1$. Dann folgt mittels Funktorialität 
\begin{align*}
	f_* = (H \circ i_0)_* = H_* \circ (i_0)_* = H_* \circ (i_1)_* = (H \circ i_1)_* = g_* \bewende
\end{align*}
% subsection 78 (end)

\subsection[Bemerkung zum Vorgehen beim Beweis von \textbf{[\#]}]{Bemerkung} % (fold)
\label{sub:79}
Zum Beweis von \eqref{eq:78} werden wir eine explizite Kettenhomotopie $h_*$ zwischen $C_*(i_0;R)$ und $C_*(i_1;R)$ konstruieren. Wir brauchen also 
$h_n : C_n(X;R) \to C_{n+1}(X \times [0,1];R)$ für alle $n \in \mathds{N}$ mit 
\begin{equation*}
	\partial_{n+1} \circ h_n + h_{n-1} \circ \partial_n = C_n(i_0;R)- C_n(i_1;R). \label{eq:79} \tag{\#\#}
\end{equation*}

\minisec{"Ansatz"}
Für $\sigma : \abs*{\Delta^n} \to X $ setze $h_n(\sigma) = \sigma \times \id_{[0,1]} : \abs*{\Delta^n} \times [0,1] \to X \times [0,1]$. Dann ist
"$h_{n-1}(\partial_n \sigma) = \partial_n \sigma \times \id_{[0,1]}$" und 
\[
	\partial_{n+1}(h_n \sigma) = \partial_{n+1}(\sigma \times \id_{[0,1]}) = (\sigma \times \id)\Big|_{\partial (\abs*{\Delta^n} \times [0,1] )} = \partial_n \sigma \times \id_{[0,1]} + i_0 \sigma + i_1 \sigma
\]
da $\partial \enbrace*{\abs*{\Delta^n} \times [0,1]} = \partial \abs*{\Delta^n} \times [0,1] \dot{\cup} \abs*{\Delta^n} \times \set{0,1}$.

Um daraus Sinn zu machen, zerlegen wir $\abs*{\Delta^n} \times [0,1]$ in eine Vereinigung von $(n+1)$-Simplizes. $\Delta^n \times [0,1]$ hat die folgende Form für
$n \in \set{0,1,2}$:
\begin{figure}[H]
	\centering{
	\scalebox{0.9}{\begin{tikzpicture}[scale=2.0]
		\coordinate (1) at (0,1);
		\coordinate (2) at (0,0);
		
		\coordinate (3) at (1.7,1);
		\path (3) ++ (1,0) coordinate (4) ++ (0,-1) coordinate (6) ++ (-1,0) coordinate (5);
		% \coordinate (4) at (3,1);
		% \coordinate (5) at (2,0);
		% \coordinate (6) at (3,0);
		
		\coordinate (7) at (4.5,1);
		\path (7) ++ (0.9,-0.15) coordinate (8);
		\path (7) ++ (0.5,0.18) coordinate (9);
		\path (7) ++ (0,-1) coordinate (10);
		\path (8) ++ (0,-1) coordinate (11);
		\path (9) ++ (0,-1) coordinate (12);
		
		\draw[thick] (2) -- (1) node[above, yshift=0.4cm]{$\Delta^0$} node[midway, left, xshift=-0.3cm]{$[0,1]$};
		
		\draw[thick] (3) node[left]{$(0,0)$}-- (4) node[midway, above, yshift=0.4cm]{$\Delta^1$} node[right]{$(1,0)$} -- (6) node[right]{$(1,1)$} 
		-- (5) node[left]{$(0,1)$} -- cycle;
		
		
		\draw[thick] (7) node[left]{$(0,0)$} node[above, yshift=0.4cm, xshift=0.9cm]{$\Delta^2$} -- (8) node[right]{$(2,0)$} -- (9) node[right]{$(1,0)$} -- cycle;
		
		\draw[thick] (7) -- (10) node[left]{$(0,1)$} -- (11) node[right]{$(2,1)$} -- (8);
		\draw[thick, dashed] (10) -- (12) node[above left]{$(1,1)$} -- (9);
		\draw[thick, dashed] (12) -- (11);
		
		\foreach \x in {1,...,12}{
			\filldraw (\x) circle(0.03);
		}
	\end{tikzpicture}}
	\caption{$\Delta^n \times [0,1]$ für $n=0,1,2$}}
\end{figure}
\noindent Die Zerlegung in $3$-Simplizes ist für $\Delta^2 \times [0,1]$ gegeben durch
\begin{figure}[H]
	\centering{
	\[
		\vcenter{\hbox{
		\begin{tikzpicture}[scale=2]
			\coordinate (7) at (0,1);
			\path (7) ++ (0.9,-0.15) coordinate (8);
			\path (7) ++ (0.5,0.18) coordinate (9);
			\path (7) ++ (0,-1) coordinate (10);
			\path (8) ++ (0,-1) coordinate (11);
			\path (9) ++ (0,-1) coordinate (12);
		
		
			\draw[thick] (7)  -- (8) -- (9) -- cycle;
			\draw[thick] (7) -- (10)  -- (11) -- (8);
			\draw[thick, dashed] (10) -- (12) -- (9);
			\draw[thick, dashed] (12) -- (11);
		
			\foreach \x in {7,...,12}{
				\filldraw (\x) circle(0.02);
			}
		\end{tikzpicture}
		}}
		\enspace =\enspace
		\vcenter{\hbox{
		\begin{tikzpicture}[scale=2]
			\coordinate (7) at (0,1);
			\path (7) ++ (0.9,-0.15) coordinate (8);
			\path (7) ++ (0.5,0.18) coordinate (9);
			\path (7) ++ (0,-1) coordinate (10);
			\path (8) ++ (0,-1) coordinate (11);
			\path (9) ++ (0,-1) coordinate (12);
		
		
			\draw[thick, gray] (7)  -- (8) -- (9) -- cycle;
			\draw[thick, gray] (7) -- (10)  -- (11) -- (8);
			\draw[thick, dashed, gray] (10) -- (12) -- (9);
			\draw[thick, dashed, gray] (12) -- (11);
			\draw[very thick, Firebrick3] (7) -- (10) -- (11) -- (12) -- (7) -- (11) -- (12) -- (10) -- cycle;
		
			\foreach \x in {7,...,12}{
				\filldraw (\x) circle(0.02);
			}
		\end{tikzpicture}
		}} 
		\enspace\cup \enspace
		\vcenter{\hbox{
		\begin{tikzpicture}[scale=2]
			\coordinate (7) at (0,1);
			\path (7) ++ (0.9,-0.15) coordinate (8);
			\path (7) ++ (0.5,0.18) coordinate (9);
			\path (7) ++ (0,-1) coordinate (10);
			\path (8) ++ (0,-1) coordinate (11);
			\path (9) ++ (0,-1) coordinate (12);
		
		
			\draw[thick, gray] (7)  -- (8) -- (9) -- cycle;
			\draw[thick, gray] (7) -- (10)  -- (11) -- (8);
			\draw[thick, dashed, gray] (10) -- (12) -- (9);
			\draw[thick, dashed, gray] (12) -- (11);
			
			\draw[very thick, Goldenrod3] (7) -- (9) -- (12) -- (7) -- (11) -- (12) -- (9) -- (11) -- cycle;
		
			\foreach \x in {7,...,12}{
				\filldraw (\x) circle(0.02);
			}
		\end{tikzpicture}
		}} 
		\enspace\cup \enspace
		\vcenter{\hbox{
		\begin{tikzpicture}[scale=2]
			\coordinate (7) at (0,1);
			\path (7) ++ (0.9,-0.15) coordinate (8);
			\path (7) ++ (0.5,0.18) coordinate (9);
			\path (7) ++ (0,-1) coordinate (10);
			\path (8) ++ (0,-1) coordinate (11);
			\path (9) ++ (0,-1) coordinate (12);
		
		
			\draw[thick, gray] (7)  -- (8) -- (9) -- cycle;
			\draw[thick, gray] (7) -- (10)  -- (11) -- (8);
			\draw[thick, dashed, gray] (10) -- (12) -- (9);
			\draw[thick, dashed, gray] (12) -- (11);
			
			\draw[very thick,RoyalBlue3] (7) -- (8) -- (9) -- (11) -- (7) -- (8) -- (11)-- (9) -- cycle;
		
			\foreach \x in {7,...,12}{
				\filldraw (\x) circle(0.02);
			}
		\end{tikzpicture}
		}}
	\]
	\caption{Zerlegung von $\Delta^2 \times [0,1]$ in $3$-Simplizes}}
\end{figure}
\noindent Diese verallgemeinern wir jetzt für beliebiges $n \in \mathds{N}$:
% subsection 79 (end)

\subsection[Definition: Affin lineare Abbildungen zur Konstruktion der Kettenhomotopie]{Definition} % (fold)
\label{sub:710}
Für $j=0, \ldots ,n$ seien $k_{n,j} : \abs{\Delta^{n+1}} \to \abs{\Delta^n} \times [0,1]$ und $\iota_{n,j} : \abs*{\Delta^{n-1}} \to \abs*{\Delta^n}$\marginnote{
$e_1, \ldots , e_n$ sind wieder die Einheitsvektoren von $\mathds{R}^n$} die eindeutigen affin linearen Abbildungn, für die gilt:
\begin{align*}
	k_{n,j}(e_0)&=(e_0,0) & \iota_{n,j}(e_0) &= e_0\\
	k_{n,j}(e_1)&= (e_1,0) \enspace\ldots & \iota_{n,j}(e_1)&=e_1  \enspace \ldots \\
	k_{n,j}(e_j)&= (e_j,0) & \iota_{n,j}(e_{j-1})&= e_{j-1}\\
	k_{n,j}(e_{j+1})&= (e_j,1) \enspace\ldots & \iota_{n,j}(e_j)&= e_{j+1} \enspace \ldots \\
	k_{n,j}(e_{n+1})&= (e_n,1)& \iota_{n,j}(e_{n-1}) &= e_n
\end{align*}
% subsection 710 (end)

\subsection[Lemma: Gleichungen für die Abbildungen aus \ref{sub:710}]{Lemma} % (fold)
\label{sub:711}
Es gelten folgende fünf Gleichungen für die Abbildungen $k_{n,j}$ und $\iota_{n,j}$:
\begin{enumerate}[(i)]
	\item Für $0 \le l < j \le n$ gilt $k_{n,j} \circ \iota_{n+1,l} = \enbrace*{\iota_{n,l} \times \id_{[0,1]}} \circ k_{n-1,j-1}$
	\item Für $1 \le j+1 < l \le n+1$ gilt $k_{n,j} \circ \iota_{n+1,l} = \enbrace*{\iota_{n,l-1} \times \id_{[0,1]}} \circ k_{n-1,j} $
	\item Für $1 \le j+1 = l \le n$ gilt $k_{n,j} \circ \iota_{n+1,l} = k_{n,j+1} \circ \iota_{n+1,l}$
	\item Für $l=0,j=0$ ist $k_{n,j} \circ \iota_{n+1,l} = i_1 :  \Delta^n \to \Delta^n \times [0,1]$
	\item Für $l=n,j=n+1$ ist $k_{n,j} \circ \iota_{n+1,l} = i_0 : \Delta^n \to \Delta^n \times [0,1]$
\end{enumerate}
\minisec{Beweis}
Wir zeigen hier nur exemplarisch die erste Gleichung:
Die linke Seite der Gleichung entspricht
	\[
		\begin{tikzcd}[row sep=2.4em, column sep=0.5em, font=\small]
			\abs*{\Delta^n} \dar["\iota_{n+1,l}" description] && e_0 \dar[mapsto] & e_{l-1} \dar[mapsto] & e_l \drar[mapsto]& \cdots & e_j \drar[mapsto]& \cdots 
			& e_n \drar[mapsto]& \\
			\abs*{\Delta^{n+1}} \dar["k_{n,j}" description] && e_0 \dar[mapsto] & e_{l-1} \dar[mapsto] & e_l \dar[mapsto] & e_{l+1} \dar[mapsto]& e_{j} \dar[mapsto] 
			& e_{j+1} \arrow[ddl,mapsto, bend left] & e_n & e_{n+1} \arrow[ddl,mapsto, bend left] \\
			\abs*{\Delta^{n}} \times [0,1] && (e_0,0) & (e_{l-1},0) & (e_l,0) & (e_{l+1},0)& (e_j,0) & & \\
			 &&  &  &  & & (e_j,1) & \cdots & (e_n,1)
		\end{tikzcd}
	\]
	und die rechte Seite entspricht
	\[
		\begin{tikzcd}[row sep=2.4em, column sep=0.2em, font=\small]
			\abs*{\Delta^n} \dar["k_{n-1,j-1}" description] && e_0 \dar[mapsto] & e_{l-1} \dar[mapsto] & e_l \dar[mapsto]& \cdots 
			& e_{j-1} \dar[mapsto]& e_j \arrow[ddl,bend left,mapsto]& e_{j+1} \arrow[ddl,bend left,mapsto] & \cdots & e_n \arrow[ddl,mapsto,bend left] \\
			\abs*{\Delta^{n-1}} \times [0,1] \arrow[dd,"\iota_{n,l}\times \id_{[0,1]}" description] && (e_0,0) \arrow[dd,mapsto] & (e_{l-1},0) \arrow[dd,mapsto] 
			& (e_l,0) \arrow[ddr, mapsto] & \cdots & (e_{j-1},0) \arrow[dd, controls={+(225:1.2) and +(135:1.2)}, mapsto] &   &  & \cdots \\
			 &&&&&& (e_{j-1},1) \arrow[dd, controls={+(-45:1) and +(45:1)}, end anchor={[xshift=2ex]north}, mapsto] & (e_{j},1) \arrow[dd, mapsto] & \cdots 
			 & (e_{n-1},1) \arrow[dd,mapsto]\\
			 \abs*{\Delta^n} \times [0,1] && (e_0,0)  & (e_{l-1},0) & (e_l,0) & (e_{l+1},0) &(e_j,0) & & & \\
			 &&&&&& (e_j,1) & (e_{j+1},1) & \cdots & (e_n,1)
		\end{tikzcd}
	\]
	Also folgt insgesamt die Gleichheit. Die anderen Gleichungen folgen genauso. \bewende
% subsection 712 (end)

\subsection[Beweis von \textbf{[\#\#]}/Konstruktion der Kettenhomotopie]{Beweis von \protect\eqref{eq:79}} % (fold)
\label{sub:712}
Sei $h_n : C_n(X;R) \to C_{n+1}(X \times [0,1];R)$ definiert durch
\[
	h_n(\sigma) := \sum_{j=0}^{n} (-1)^j \cdot \enbrace*{\sigma \times \id} \circ k_{n,j}
\]
$h_*$ ist die gesuchte Kettenhomotopie, da für $n \in \mathds{N}$ unter Anwendung der Gleichungen aus \ref{sub:712} gilt:
\begin{align*}
	\partial_{n+1} \enbrace[\big]{h_n(\sigma)} &= \partial_{n+1} \enbrace*{\sum_{j=0}^{n} (-1)^j \cdot (\sigma \times \id) \circ k_{n,j}} \\
	&=  \sum_{l=0}^{n+1} \sum_{j=0}^{n} (-1)^{j+l} \cdot \enbrace*{ \sigma \times \id} \circ k_{n,j} \circ \iota_{n+1,l} \tag*{(nach Definition \ref{sub:52})}\\
	&= \smashoperator[r]{\sum_{0 \le l < j \le n}} \, (-1)^{j+l} \enbrace[\big]{(\sigma \circ \iota_{n,l}) \times \id} \circ k_{n-1,j-1} 
	+ \enspace\smashoperator{\sum_{1 \le j+1 < l \le n+1}}\, (-1)^{j+l} \enbrace[\big]{(\sigma \circ \iota_{n,l-1}) \times \id} \circ k_{n-1,j} \\
	& \quad + \smashoperator[r]{\sum_{1 \le j+1=l \le n}} \,(-1)^{j+l} (\sigma \times \id) \circ k_{n,j+1} \circ \iota_{n+1,l} 
	+ \smashoperator[r]{\sum_{1 \le j=l \le n}} \, (-1)^{j+l} (\sigma\times\id) \circ k_{n,j}\circ \iota_{n+1,l} \\
	& \quad + \sum_{0=l=j}  (-1)^{j+l} (\sigma \times \id) \circ i_1 
	+ \smashoperator[r]{\sum_{n+1=j+1=l}}\, (-1)^{j+l} (\sigma \times \id) \circ i_0 \\
	&= \smashoperator[r]{\sum_{0 \le l \le j \le n-1}}\, (-1)^{j+l+1} \enbrace[\big]{(\sigma \circ \iota_{n,l}) \times \id} \circ k_{n-1,j} 
	+ \enspace\smashoperator{\sum_{1 \le j+1 \le l \le n}} \,(-1)^{j+l+1} \enbrace[\big]{(\sigma \circ \iota_{n,l}) \times \id} \circ k_{n-1,j} \\
	& \quad + \underbrace{\smashoperator[r]{\sum_{1 \le j=l \le n}}\, (-1)^{j+l-1} (\sigma \times \id) \circ k_{n,j} \circ \iota_{n+1,l} 
	+ \smashoperator[r]{\sum_{1 \le j=l \le n}} \,(-1)^{j+l} (\sigma\times \id) \circ k_{n,j} \circ \iota_{n+1,l}}_{=0} \\
	& \quad + \sum_{0=l=j}  (-1)^{j+l} (\sigma \times \id) \circ i_1 + \smashoperator[r]{\sum_{n+1=j+1=l}} (-1)^{j+l} (\sigma\times \id) \circ i_0 \\
	&= - \smashoperator[r]{\sum_{\substack{0 \le l\le n \\ 0 \le j \le n-1}}} \,(-1)^{j+l} \enbrace[\big]{ \enbrace*{\sigma \circ \iota_{n,l}} \times \id } \circ k_{n-1,j} 
	+ (\sigma\times \id) \circ i_1  - (\sigma\times \id) \circ i_0 \\
	&= - \sum_{j=0}^{n-1} (-1)^j \sum_{l=0}^{n} (-1)^{l} \enbrace[\big]{ \enbrace*{\sigma \circ \iota_{n,l}} \times \id } \circ k_{n-1,j}
	+ \underbrace{(\sigma\times \id) \circ i_1}_{= i_0 \circ \sigma}  - \underbrace{(\sigma\times \id) \circ i_0}_{= i_1 \circ \sigma} \\
	&= - h_{n-1}(\partial_n \sigma) + C_n(i_1;R)(\sigma) - C_n(i_0;R)(\sigma) \bewende  
\end{align*}
% subsection 713 (end)

\subsection[Bemerkung: Konstruktion einer Kettenhomotopie aus einer Homotopie]{Bemerkung} % (fold)
\label{sub:713}
Ist $H : X \times [0,1] \to Y$ eine Homotopie zwischen $f$ und $g$, so erhalten wir eine Kettenhomotopie $k$ zwischen $C_*(f;R)$ und $C_*(g;R)$ durch
\[
	k_n(\sigma) := C_{n+1}(H;R) \circ h_n = \sum_{j=0}^{n} (-1)^j H \circ (\sigma \times \id) \circ k_{n,j}
\]
% subsection 714 (end)
% section 7 (end)
\newpage

\section{Homologie von Paaren} % (fold)
\label{sec:8}

\subsection[Definition: Paar von topologischen Räumen]{Definition} % (fold)
\label{sub:81}
Sei $A$ ein Teilraum von $X$, dann heißt $(X,A)$ ein \Index{Paar von topologischen Räumen}. Eine Abbildung von Paaren $f : (X,A) \to (Y,B)$ ist eine stetige Abbildung
$f : X \to Y$ mit $f(A) \subseteq B$. Manchmal schreiben wir $(f,f|_{A}) : (X,A) \to (Y,B)$. 

\noindent Die Kategorie von Paaren von topologischen Räumen bezeichnen wir mit $\TOP^2$.
% subsection 81 (end)

\subsection[Definition: Singulärer Kettenkomplex und Homologiemodul von Paaren]{Definition} % (fold)
\label{sub:82}
Sei $(X,A)$ ein Paar. Dann definieren wir den \bet{singulären Kettenkomplex}\index{singulärer Kettenkomplex!von Paaren} von $(X,A)$ über $R$ durch
\begin{align*}
	C_n(X,A;R) &:= \nicefrac{C_n(X;R)}{C_n(A;R)} \\
	\partial_n \enbrace[\big]{\sigma+ C_n(A;R)} &:= \partial_n \sigma + C_{n-1}(A;R)
\end{align*}
Dies ist wohldefiniert, da $\partial_n \enbrace[\big]{C_n(A;R)} \subseteq C_{n-1}(A;R)$. Weiter heißt
\[
	H_n(X,A;R) := H_n \enbrace[\big]{C_*(X,A;R),\partial_*} 
\]
der \bet{$\mathbf{n}$-te singuläre Homologiemodul}\index{n-ter Homologiemodul@$n$-ter Homologiemodul!singulär, von Paaren} von $(X,A)$ mit Koeffizienten in $R$.
% subsection 82 (end)

\subsection[Definition: Der singuläre Kettenkomplex von Paaren definiert einen Funktor]{Bemerkung} % (fold)
\label{sub:83}
Der singuläre Kettenkomplex von Paaren definiert einen Funktor: $\TOP^2 \to R\text{-}\textsc{Ketten}$. Durch Komposition mit Homologie (als Funktor 
$R\text{-}\textsc{Ketten} \to R\text{-}\MOD$) erhalten wir einen Funktor
\[
	H_n(-,-;R) : \TOP^2 \longrightarrow R\text{-}\MOD
\]
% subsection 83 (end)

\subsection[Bemerkung: Topologische Räume als Paar auffassen]{Bemerkung} % (fold)
\label{sub:84}
Via $X \mapsto (X,\emptyset)$ können wir jeden Raum auch als Paar auffassen. Es gilt $H_n(X,\emptyset;R) = H_n(X;R)$.
% subsection 84 (end)

\subsection[Frage: Können wir $H_n(X,A;R)$ durch $H_n(X;R)$ und $H_n(A;R)$ ausdrücken?]{Frage} % (fold)
\label{sub:85}
Können wir $H_n(X,A;R)$ durch $H_n(X;R)$ und $H_n(A;R)$ ausdrücken? 
\begin{description}
	\item[Ansatz 1:] 
	\[
		H_n(X,A;R) \cong \nicefrac{H_n(X;R)}{H_n(A;R)}
	\]
	\ifxetexorluatex \faWarningSign \fi \, Problem: $H_n(A;R)$ ist kein Untermodul von $H_n(X;R)$
	\item[Ansatz 2:] Ist $H_n(X;R) \to H_n(X,A;R)$ surjektiv?
	
	Sei $x \in H_n(X,A;R)$. Dann gibt es $\sigma + C_n(A;R) \in C_n(X,A;R)$ mit $\partial_n \enbrace[\big]{\sigma + C_n(A;R)} = 0$ und $x= [\sigma]$. Es ist 
	$\sigma \in C_n(X;R)$, aber wir wissen nur $\partial_n \sigma \in C_{n-1}(A;R)$, nicht $\partial_n(\sigma)= 0$. \hfill \ifxetexorluatex {\large \faMeh} \fi
\end{description}
% subsection 85 (end)

\subsection[Beispiel: Homologie des Paares $(\abs*{\Delta^n}, \abs*{\partial \Delta^n})$]{Beispiel} % (fold)
\label{sub:86}
In $C_n \enbrace[\big]{\abs*{\Delta^n}, \abs*{\partial \Delta^n};R}$ gilt für $\sigma := \id : \abs*{\Delta^n} \to \abs*{\Delta^n}$, dass $\partial_n([\sigma])=0$ ist, da 
$\partial_n \sigma \in C_{n-1} \enbrace*{\abs*{\partial \Delta^n};R}$. Wir werden später sehen, dass $[\sigma]$ den Homologiemodul
\[
	H_n \enbrace[\big]{\abs*{\Delta^n}, \abs*{\partial \Delta^n};R  } \cong R 
\]
erzeugt.
% subsection 86 (end)

\subsection[Definition: Lange exakte Sequenz von $R$-Moduln]{Definition} % (fold)
\label{sub:87}
Eine Folge von $R$-linearen Abbildungen zwischen $R$-Moduln $M_i$
\[
	\begin{tikzcd}
		\cdots \rar & M_n \rar["f_n"]  & M_{n-1} \rar["f_{n-1}"] & M_{n-2} \rar & \cdots \rar & M_1 \rar & M_0
	\end{tikzcd}
\]
heißt \Index{lange exakte Folge}, wenn sie exakt an jeder Stelle $M_i$ ist, d.h. für alle $i$ gilt $\im f_i = \ker f_{i-1}$.
% subsection 87 (end)

\subsection[Satz: Existenz der Paarsequenz]{Satz} % (fold)
\label{sub:88}
Es gibt eine natürliche Transformation $\partial_n$ von $(X,A) \mapsto H_n(X,A;R)$ nach $(X,A) \mapsto H_{n-1}(A;R)$, sodass für jedes Paar $(X,A)$
\begin{equation*}
	\begin{tikzcd}[row sep=1em, column sep=2em]
		\cdots  \rar["\partial_{n+1}"] & H_n(A;R) \rar["i_*"] & H_n(X;R) \rar["j_*"] & H_n(X,A;R) \rar["\partial_n"] 
		& H_{n-1}(A;R) \rar & \cdots \\
		\cdots \rar["\partial_1"] & H_0(A;R) \rar & H_0(X;R) \rar & H_0(X,A;R) \rar & 0
	\end{tikzcd} \label{eq:88}\tag{*}
\end{equation*}
eine lange exakte Folge ist. Dabei sind $i : A \hookrightarrow X$ und $j : (X,\emptyset) \hookrightarrow (X,A)$ die Inklusionen. Das bedeutet: \smallskip\\
Für jedes Paar $(X,A)$ haben wir eine $R$-lineare Abbildung $\partial_n : H_n(X,A;R) \to H_{n-1}(A;R)$, sodass für jede Abbildung $(f,f|_{A}) : (X,A) \to (Y,B)$ von Paaren
folgendes Diagramm kommutiert
\[
	\begin{tikzcd}
		H_{n}(X,A;R) \rar["\partial_n"] \dar["{(f,f|_A)_*}"] & H_{n-1}(A;R) \dar["{(f|_A)_*}"] \\
		H_n(Y,B;R) \rar["\partial_n"] & H_{n-1}(B;R)
	\end{tikzcd}
\]
\minisec{Beweis}
\emph{siehe \ref{sub:813}.}
% subsection 88 (end)

\subsection[Korollar: Hologiemodul wenn $A$ kontraktibel ist]{Korollar} % (fold)
\label{sub:89}
\begin{enumerate}[(1)]
	\item Ist die Inklusion $i  :A \to X$ eine Homotopieäquivalenz, so ist $H_n(X,A;R)= 0$ für alle $n$.
	\item Sei $A$ kontraktibel. Dann gilt für alle $n \ge 1$. $H_n(X;R) \cong H_n(X,A;R)$
\end{enumerate}
\minisec{Beweis}
\begin{enumerate}[(1)]
	\item Betrachte den folgenden Ausschnitt aus der langen exakten Folge:
	\[
		\begin{tikzcd}
			H_n(A;R) \rar["i_*"] & H_n(X;R) \rar["j_*"] & H_n(X,A;R) \rar["\partial_n"] & H_{n-1}(A;R) \rar["i_*"] & H_n(X;R)
		\end{tikzcd}
	\]
	Nach \ref{sub:74} ist $i_* : H_n(A;R) \to H_n(X;R)$ ein Isomorphismus. Wegen der Exaktheit gilt
	$\ker j_* = \im i_* = H_n(X;R)$, also folgt $j_* = 0$. Da $\im \partial_n = \ker i_* = 0$, folgt $\partial_n=0$. Nun ist 
	\[
		H_n(X,A;R) = \ker \partial_n = \im j_* = 0
	\]
	\item Übung! \bewende
\end{enumerate}
% subsection 89 (end)

\subsection[Definition: Kurze exakte Folge von Kettenkomplexen]{Definition} % (fold)
\label{sub:810}
Seien 
\begin{equation*}
	\begin{tikzcd}
		(C_*,d_*) \rar["i_*"] & (C_*',d_*') \rar["p_*"] & (C_*'',d_*'') 
	\end{tikzcd} \label{eq:810}\tag{\#}
\end{equation*}
\hyperref[sub:61]{$R$-Kettenabbildungen}. Ist für jedes $n$ die Folge
\(
	\begin{tikzcd}
		C_n \rar["i_n"] & C_n' \rar["p_n"] & C_n''
	\end{tikzcd}
\)
kurz exakt, so heißt \eqref{eq:810} eine \Index{kurze exakte Folge von Kettenkomplexen}.
% subsection 810 (end)

\subsection[Beispiel einer kurzen exakten Folge von Kettenkomplexen]{Beispiel} % (fold)
\label{sub:811}
Für jedes Paar $(X,A)$ ist 
\[
	\begin{tikzcd}
		\enbrace*{C_*(A;R),\partial_*} \rar["i_*"] & \enbrace*{C_*(X;R),\partial_*} \rar["j_*"] & \enbrace*{C_*(X,A;R),\partial_*}   
	\end{tikzcd}
\]
eine kurze exakte Folge von $R$-Kettenkomplexen.
% subsection 811 (end)

\subsection{Schlangenlemma} % (fold)
\label{sub:812}
Sei 
\[
	\begin{tikzcd}
		(C_*,d_*) \rar["i_*"] & (C_*',d_*') \rar["p_*"] & (C_*'',d_*'')
	\end{tikzcd}
\]
eine kurze exakte Folge von $R$-Kettenkomplexen.
\begin{enumerate}[a)]
	\item Für jedes $n$ gibt es eine eindeutige wohldefinierte $R$-lineare Abbildung $\partial_n : H_n(C_*'',d_*'') \to H_{n-1}(C_*,d_*)$ mit:
	Für $v' \in C_n'$ mit $d_n'' \circ p_n(v')=0$ ist 
	\[
		\partial_n \enbrace[\Big]{\underbrace{\benbrace*{p_n(v')}}_{\mathclap{\in H_n(C_*'',d_*'')}} } = [v] \in H_{n-1}(C_*,d_*)
	\]
	wobei $v \in C_{n-1}$ bestimmt ist durch $i_{n-1}(v)= d_n'(v')$.
	\item 
	\[
		\begin{tikzcd}
			\cdots \rar["\partial_{n+1}"] & H_n(C_*,d_*) \rar["i_*"] & 
				H_n(C_*',d_*') \arrow[d, phantom, ""{coordinate, name=Z}] \rar["p_*"] 
					& H_n(C_*'',d_*'') \arrow[dll, "\partial_n" very near start, rounded corners, to path={
																	-- ([xshift=2ex]\tikztostart.east)
																	|- (Z) [near end]\tikztonodes
																	-| ([xshift=-2ex]\tikztotarget.west)
																	-- (\tikztotarget)
																	}] & \\
			& H_{n-1}(C_*,d_*) \rar["i_*"] 
				& H_{n-1}(C_*',d_*') \rar["p_*"] \arrow[d, phantom, ""{coordinate, name=R}]
					& \quad \cdots  \quad \arrow[dll, "\partial_1" very near start, rounded corners, to path={
																	-- ([xshift=2ex]\tikztostart.east)
																	|- (R) [near end]\tikztonodes
																	-| ([xshift=-2ex]\tikztotarget.west)
																	-- (\tikztotarget)
																	}]& \\
			& H_0(C_*,d_*) \rar["i_*"] & H_0(C_*',d_*') \rar["p_*"] & H_0(C_*'',d_*'') \rar & 0
		\end{tikzcd}
	\]
	ist eine lange exakte Folge. \todo{RevChap8}
\end{enumerate}
\minisec{Beweis}
\[
	\begin{tikzcd}[row sep=2em, column sep=3em]
		C_{n+1} \rar["i_{n+1}"] \dar["d_{n+1}"] & C_{n+1}' \rar["p_{n+1}"] \dar["d_{n+1}'"] & C_{n+1}'' \dar["d_{n+1}''"] \\
		C_n \rar["i_n"] \dar["d_{n}"]& C_n' \rar["p_n"] \dar["d_{n}'"]& C_n'' \dar["d_{n}''"]\\
		C_{n-1} \rar["i_{n-1}"] \dar["d_{n-1}"]& C_{n-1}'  \rar["p_{n-1}"] \dar["d_{n-1}'"]& C_{n-1}'' \dar["d_{n-1}''"]\\
		C_{n-2} \rar["i_{n-2}"]& C_{n-2}' \rar["p_{n-2}"] & C_{n-2}''
	\end{tikzcd}
\]
Sei $x'' \in H_n(C_*'',d_*'')$, also $x'' = [v'']$ mit $v'' \in \ker d_n''$. Da $p_n$ surjektiv ist, existiert $v' \in C_n'$ mit $p_n(v')=v''$. Es gilt
\[
	p_{n-1} d_n'(v') = d_n'' p_n(v') = d_n''(v'') = 0
\]
$\Rightarrow v' \in \ker p_{n-1}$. Weiter gilt $\ker p_{n-1}= \im i_{n-1}$, also $\exists v \in C_{n-1} : i_{n-1}(v)= d_n'(v')$. Dann gilt
\[
	i_{n-2} \enbrace[\big]{d_{n-1}(v)} = d'_{n-1} \enbrace[\big]{i_{n-1}(v)} = d_{n-1}' \enbrace[\big]{d_n'(v')} = 0   
\]
Da $i_{n-2}$ injektiv ist, folgt somit $d_{n-1}(v)=0$. Setze $\partial_n(x'') := [v]$. Zu zeigen: $[v] \in H_{n-1}(C_*,d_*)$ ist unabhängig von der Wahl von $v''$ und $v'$.
Seien $w'' \in \ker d_n''$, $w' \in C_n'$, $w \in C_{n-1}$ mit $x''=[w'']$, $p_n(w')=w''$ und $i_{n-1}(w)= d_n'(w')$. 

$[v'']=[w'']$, also folgt $\exists a'' \in C_{n+1}'' : d_{n+1}'' (a'')=v'' - w''$. $p_{n+1}$ ist surjektiv, also existiert $a' \in C_{n+1}'$ mit
$p_{n+1}(a')=a''$. Es gilt
\[
	p_n(v'-w'- d_{n+1}'(a')) = v'' - w'' - \underbrace{d_{n+1}(p_{n+1}(a'))}_{v''- w''} = 0
\]
Mit $\ker p_n = \im i_n$ folgt: $\exists a \in C_n$ mit $i_n(a)=v'- w' - d_{n+1}(a')$. Es bleibt zu zeigen: $d_n(-a)= v-w$.
\begin{align*}
	i_{n-1} \enbrace[\big]{-d_n(a)- (v-w)} & = -i_{n-1} \enbrace*{d_n(a)} - i_{n-1}(v-w) \\ &= - d_n' \enbrace*{i_n(a)} - \enbrace*{d_n'(v')- d_n'(w')}  \\  
	&= -d_n' \enbrace*{ d_{n+1}' (a') - (v'-w')} - \enbrace*{ d_n'(v') - d_n'(w')}  \\
	&= - d_n' d_{n+1}' (a') + d_n'(v'- w') -  \enbrace*{d_n'(v'-w')} = 0
\end{align*}
da $i_{n-1}$ injektiv ist, folgt $-d_n(a)- (v-w) = 0$, also $d_n(-a)=(v-w)$.
\begin{description}
	\item[\uline{$\im (i_n)_* = \ker (p_n)_*$:}] 
	\begin{description}
		\item["$\subseteq$":] $(p_n)_* \circ (i_n)_* = \enbrace*{p_n-i_n}_* = (0)_* =0 $
		\item["$\supseteq$":] Sei $x' \in \ker (p_n)_*$. Sei $x'= [v']$ mit $v' \in C_n'$, $d_n'(v')=0$. Da $\benbrace*{p_n(v')}= (p_n)_* [x'] = 0$ gibt es $a'' \in C_n''$
		mit $d_{n+1}''(a'') = p_n(v')$. Da $p_{n+2}$ surjektiv ist, existiert $a' \in C_{n+1}'$ mit $p_{n+1}(a')=a''$. Dann gilt
		\begin{align*}
			p_n \enbrace*{v' - d_{n+1}'(a')} = p_n(v') - p_n d_{n+1}'(a') = p_n(v') -  \underbrace{ d_{n+1}''\enbrace[\Big]{\underbrace{p_{n+1}(a')}_{=a''}}}_{p_n(v')} = 0  
		\end{align*}
		Da $\ker p_n = \im i_n$ gibt es $a \in C_n$ mit $i_n(a)= v'- d_{n+1}'(a')$. Nun ist 
		\[
			i_{n-1} \enbrace[\big]{d_n(a)} = d_n' \enbrace[\big]{i_n(a)} = d_n' \enbrace[\big]{v'- d_{n+1}(a')}  = d_n'(v') = 0 
		\]
		Da $i_{n-1}$ injektiv ist, folgt $d_n(a)=0$. Insbesondere $[a] \in H_n \enbrace*{(C_*,d_*)}$. Nun ist 
		\[
			(i_n)_* [a] = [i_n(a)] = \benbrace*{v'- d_{n+1}'(a')} = [v'] = x' 
		\]
		Also $x' \in \im (i_n)_*$.
	\end{description}
	\item[\uline{$\im \partial_{n+1} = \ker (i_n)_*$:}] \begin{description}
		\item["$\subseteq$":] $i_n \enbrace*{d_{n+1}(x'')}= 0 \enspace\checkmark$
		\item["$\supseteq$":] Sei $x \in \ker (i_n)_*$. Sei $v \in C_n$, $d_n(v)=0$, $[v]=x$. $\leadsto x= d_n [v'']$.
	\end{description} 
\end{description}
% subsection 812 (end)

\subsection[Beweis von Satz 8.8]{Beweis von Satz \ref{sub:88}} % (fold)
\label{sub:813}
Für jedes Paar $(X,A)$ ist die Folge  der singulären Kettenkomplexe
\[
	\begin{tikzcd}
		\enbrace*{C_*(A;R), \partial_*^A} \rar["C_*(i;R)"] &  \enbrace*{C_*(X;R), \partial_*^X} \rar["C_*(j;R)"] & \enbrace*{C_*(X,A;R), \partial_*^{(X,A)}}
	\end{tikzcd}
\]
kurz exakt. Das Schlangenlemma {\ref{sub:812}} produziert $\partial_n : H_n(X,A;R) \to H_{n-1}(A;R)$ und die lange exakte Sequenz \eqref{eq:88}.
Es bleibt zu zeigen, dass die Randabbildungen aus dem Schlangenlemma wie behauptet eine natürliche Transformation definieren: Sei $(f,F|_{A}) : (X,A) \to (Y,B)$ eine
Abbildung von Paaren. Zu zeigen ist, dass
\[
	\begin{tikzcd}
		H_n(X,A;R) \rar["\partial_n"]  \dar["{(f,f|_A)_*}"]& H_{n-1}(A;R) \dar["{(f|_{A})_*}"] \\
		H_n(Y,B;R) \rar["\partial_n"] & H_{n-1}(B;R)
	\end{tikzcd}
\]
kommutiert. Sei $x'' \in H_n(X,A;R)$. Sei $v' \in C_n(X;R)$ mit $\partial_n^X(v') \in C_n(A;R)$ und $\benbrace*{v' + C_n(A;R)}= x''$. Dann ist 
$\partial_n(x'')= \benbrace*{\partial_n^X(v')} \in H_n(A;R)$. Dann ist 
\[
	\enbrace*{f,f|_A}_*(x'') =  \enbrace*{f,f|_A}_* \benbrace*{v' + C_n(A;R)}  = \benbrace*{f(v')+ C_n(B;R)} \in H_n(Y,B;R) 
\]
Es ist 
\(
	\partial_n^Y \enbrace[\big]{f_*(v')} = \enbrace*{f|_A}_* \enbrace[\Big]{\underbrace{\partial_n^X(v')}_{\in C_n(A;R)}} \in C_n(B;R)   
\).
Also
\begin{align*}
	\partial_n \enbrace[\big]{\enbrace*{f,f|_A}_*(x'') } = \benbrace[\Big]{\partial_n^Y \enbrace[\big]{f_*(v')} } = \benbrace*{f_* \circ \partial_n^X(v')} = f_* \partial_n(X'') \bewende  
\end{align*}
% subsection 813 (end)

\subsection{Bemerkung} % (fold)
\label{sub:814}
Für eine Abbildung $(f,f|_A) : (X,A) \to (Y,B)$ erhalten wir ein kommutierendes Diagramm:
\missingfigure{langes Diagramm}

% subsection 814 (end)

\subsection{Fünfer-Lemma} % (fold)
\label{sub:815}
Seien die Zeilen in folgendem kommutativen Diagramm von $R$-Moduln exakt.
\[
	\begin{tikzcd}
		A_4 \rar["\alpha_4"] \dar["f_4" swap, "\cong"] & A_3 \rar["\alpha_3"] \dar["f_3" swap, "\cong"] & A_2 \rar["\alpha_2"]  \dar[Firebrick3, "f_2" swap] & A_1 \rar["\alpha_1"] \dar["f_1" swap, "\cong"] & A_0  \dar["f_0" swap, "\cong"]\\
		B_4 \rar["\beta_4"] & B_3 \rar["\beta_3"] & B_2 \rar["\beta_2"] & B_1 \rar["\beta_1"] & B_0 
	\end{tikzcd}
\]
Sind $f_4, f_3, f_1$ und $f_0$ Isomorphismen, so ist auch $f_2$ ein Isomorphismus.
\minisec{Beweis}
\begin{description}
	\item[Injektivität:] Sei $a_2 \in A_2$ und $f_2(a_2)=0$. $f_1 \enbrace[\big]{\alpha_2(a_2)} = \beta_2 \enbrace[\big]{f_2(a_2)} = \beta_2(0)=0$. Da $f_1$ ein 
	Isomorphismus ist, folgt $\alpha_2(a_2)=0$. $a_2 \in \ker \alpha_2 = \im \alpha_3$, also existiert $a_2 \in A_3$ mit $\alpha_3(a_3)= a_2$. Es gilt
	\[
		\beta_3 \enbrace[\big]{f_3(a_3)} = f_2 \enbrace[\big]{\alpha_3(a_3)} = f_2(a_2)=0  
	\]
	Also $f_3(a_3) \in \ker \beta_3 = \im \beta_4$. Da $f_4$ Isomorphismus ist, gibt es $a_4 \in A_4$ mit $\beta_4 \enbrace[\big]{f_4(a_4)} = f_3(a_3)$. Nun gilt
	\[
		f_3 \enbrace[\big]{\alpha_4(a_4)- a_3} = f_3 \enbrace[\big]{\alpha_4(a_4)} - f_3(a_3) = \beta_4 \enbrace[\big]{f_4(a_4)} - f_3(a_3) = f_3(a_3)- f_3(a_3) =0
	\]
	Da $f_3$ ein Isomorphismus ist, folgt nun $\alpha_4(a_4)=a_3$. Nun ist 
	\[
		a_2= \alpha_3(\alpha_3)= \alpha_3 \enbrace[\big]{\alpha_4(a_4)} = \underbrace{\alpha_3 \circ \alpha_4}_{=0}(a_4) =0
	\]
	\item[Surjektivität:] Sei $b_2 \in B_2$. Fall 1: $\beta_2(b_2)=0$. Da $\im \beta_3 = \ker \beta_2$ und $f_3$ Isomorphismus, folgt $\exists a_3 \in A_3$ mit
	$\beta_3 \enbrace[\big]{f_3(a_3)} = b_2$. Also ist $b_2= \beta_3 \enbrace[\big]{f_3(a_3)} = f_2  \enbrace[\big]{\alpha_3(a_3)} \in \im f_2$. 
	
	Ist $b_2$ beliebig, so genügt es zu zeigen: $\exists a_2 \in A_2$ mit $\beta_2 \enbrace[\big]{b_2 - f_2(a_2)} =0$. $f_1$ ist ein Isomorphismus, also $\exists a_1 \in A_1$
	mit $f_1(a_1) = \beta_2(b_2)$. Dann folgt
	\[
		f_0 \enbrace[\big]{\alpha_1(a_1)} = \beta_1 \enbrace[\big]{f_1(a_1)} = \beta_1 \enbrace[\big]{\beta_2(b_2)} = \underbrace{\beta_1 \circ \beta_2}_{=0} (b_2)=0
	\]
	Da $f_0$ ein Isomorphismus ist, folgt $\alpha_1(a_1)=0$. Da $\ker \alpha_1 = \im \alpha_2$, folgt $\alpha_2(a_2)=a_1$. Es gilt nun
	\begin{align*}
		\beta_2 \enbrace[\big]{b_2 - f_2(a_2)} = \beta_2(b_2) - \beta_2 \enbrace[\big]{f_2(a_2)} &= \beta_2(b_2) - f_1 \enbrace[\big]{\alpha_2(a_2)} \\ 
		&= \beta_2(b_2) - f_1(a_1)
		= \beta_2(b_2) - \beta_2(b_2) =0 \bewende
	\end{align*}
\end{description}
% subsection 815 (end)

\subsection{Lemma ("2 von 3")} % (fold)
\label{sub:816}
Sei $f : (X,A) \to (Y,B)$ eine Abbildung von Paaren. Seien von den drei Abbildungen
\begin{enumerate}[(i)]
	\item $(f|_A)_* :  H_n(A;R) \to H_n(B;R)$
	\item $f_* : H_n(X;R) \to H_n(Y;R)$
	\item $(f,f|_A)_* : H_n(X,A;R) \to H_n(Y,B;R)$
\end{enumerate}
zwei für jeweils alle $n$ Isomorphismen. Dann ist auch die dritte für alle $n$ ein Isomorphismus.
\minisec{Beweis}
Da die Randabbildung in der Paarfolge eine natürliche Transformation ist, erhalten wir ein kommutatives Leiterdiagramm:\marginnote{"$R$" aus Platzgründen weggelassen}
\[
	\begin{tikzcd}[column sep=1.7em, font=\small]
		H_{n+1}(A) \rar \dar["{(f|_A)_*}"] & H_{n+1}(X) \rar \dar["f_*"] & H_{n+1}(X,A) \rar["\partial_{n+1}"] \dar["{(f,f|_A)_*}"] & H_n(A) \rar \dar["(f|_A)_*"] & H_n(X) \rar \dar["f_*"] & H_n(X,A) \rar["\partial_n"] \dar["{(f,f|_A)_*}"] & H_{n-1}(A) \dar["{(f|_A)_*}"] \\
		H_{n+1}(B) \rar & H_{n+1}(Y) \rar & H_{n+1}(Y,B) \rar["\partial_{n+1}"] & H_n(B) \rar & H_n(Y) \rar & H_n(Y,B) \rar["\partial_n"] & H_{n-1}(B) \\
	\end{tikzcd}
\]
Da die Zeilen lang exakt sind, folgt die Behauptung aus dem Fünfer-Lemma. \bewende
% subsection 816 (end)
% section 8 (end)
\newpage

\section{Ausschneidung und die Mayer-Vietoris Folge} % (fold)
\label{sec:9}

\subsection[Satz: Ausschneidungsisomorphismus]{Satz} % (fold)
\label{sub:91}
Sei $(X,A)$ ein Paar und $L \subseteq A$, sodass der Abschluss $\overline{L}$ von im Inneren $\mathring{A}$ von $A$ liegt. Dann induziert die Inklusion
$\enbrace*{X \setminus L, A \setminus L} \xrightarrow{j} (X,A)$ einen Isomorphismus $j_* : H_n(X \setminus L, A \setminus L;R) \to H_n(X,A;R)$.
\todo{RevChap9}
% subsection 91 (end)

\subsection{Satz (Mayer-Vietoris Folge)} % (fold)
\label{sub:92}
Seien $A,B \subseteq X$ offen und $X=A \cup B$. Seien $i_A : A \hookrightarrow X$, $i_B : B \hookrightarrow X$, $j_A : A \cap B \hookrightarrow A$ und 
$j_B : A \cap B \hookrightarrow B$ die Inklusionen. Dann gibt es eine Randabbildung $\partial_n = \rand{X=A\cup B}_n : H_n(X;R) \to H_{n-1}(A \cap B;R)$, sodass
\marginnote{"$R$" aus Platzgründen weggelassen}
\[
	\begin{tikzcd}[column sep=3.5em]
		\cdots \rar["\partial_{n+1}"] & H_n(A \cap B) \rar["{\begin{psmallmatrix}(j_A)_* \\ (j_B)_* \end{psmallmatrix}}"] 
		&  \bigoplus\limits_{H_n(B)}^{H_n(A)} \rar["(i_A)_* - (i_B)_*"] 
		& H_n(X) \rar["\partial_n"] & H_{n-1}(A\cap B) \rar & \cdots
	\end{tikzcd}
\]
eine lange exakte Folge ist. Die Randabbildung ist dabei eine natürlicher Transformation: Sei $A',B' \subseteq X'$ offen mit $X'=A' \cup B'$. ist $f :X \to X'$ mit
$f(A) \subseteq A'$ und $f(B) \subseteq B'$, so kommutiert
\[
	\begin{tikzcd}[column sep=4em, row sep=3em]
		H_n(X;R) \rar["\rand{X=A \cup B}_n"] \dar["f_*"]  & H_{n-1}(A \cap B;R) \dar["\enbrace*{f|_{A \cap B}}_*"] \\
		H_n(X';R) \rar["\rand{X'=A' \cup B'}_n"] & H_n(A' \cap B';R)
	\end{tikzcd}
\]
\minisec{Beweis}
Sei $L := X \setminus B$. Dann ist $L$ abgeschlossen und $L \subseteq A=\mathring{A}$. Wir erhalten daher einen Ausschneideisomorphismus 
\[
	H_n(B,A \cap B;R) = H_n(X \setminus L, A \setminus L;R) \xrightarrow{\enspace \cong \enspace} H_n(X,A;R) 
\]
Betrachte das folgende kommutative Leiterdiagramm: \marginnote{"$R$" aus Platzgründen weggelassen}
\[
	\begin{tikzcd}[font=\small, column sep=1.8em, row sep=2.5em]
		\cdots \rar & H_{n+1}(B,A \cap B) \rar["\rand{B,A \cap B}_{n+1}"] \arrow{d}{(i_B,j_A)_*}[swap]{\cong} & H_n(A \cap B) \rar["(j_B)_*"] \dar["(j_A)_*"] 
		& H_n(B) \rar["(l_B)_*"] \dar["(i_B)_*"] & H_n(B,A \cap B) \rar["\rand{B,A\cap B}_n"] \arrow{d}{(i_B,j_A)_*}[swap]{\cong}
		& H_{n-1}(A \cap B) \rar \dar["(j_A)_*"] & \cdots \\
		\cdots \rar & H_{n+1}(X,A) \rar["\rand{X,A}_{n+1}"] & H_n(A) \rar["(i_A)_*"] & H_n(X) \rar["(l_X)_*"] & H_n(X,A) \rar["\partial_n^{X,A}"] & H_{n-1}(A) \rar & \cdots
	\end{tikzcd}
\]
dabei sind $l_B : (B,\emptyset) \to (B,A \cap B)$ und $l_X : (X, \emptyset) \to (X,A)$ die Inklusionen von Paaren.
Definiere nun $\partial_n : H_n(X;R) \to H_{n-1}(A\cap B;R)$ durch
\[
	\partial_n := \rand{B, A\cap B}_n \circ  (i_B,i_A)_*^{-1} \circ (l_X)_* 
\]
Die Natürlichkeit von $\partial_{n}^{B,A \cap B}$ impliziert die Natürlichkeit von $\partial_n$. \smallskip \\

\begin{itemize}
	\item Zu zeigen: $\im \begin{psmallmatrix} (j_A)_* \\ (j_B)_* \end{psmallmatrix} = \ker \enbrace[\big]{(i_A)_* - (i_B)_*} $: 
	\begin{description}
		\item["$\supseteq $":] Sei $(a,b) \in \ker \enbrace*{(i_A)_* - (i_B)_*}$. Also $(i_A)_*(a) = (i_B)_*(b)$. Es ist 
		\[
			(i_B,j_A)_* \enbrace[\big]{(l_B)_*(b)} = (l_X)_* \circ (j_B)_* (b) = \underbrace{(l_X)_* \circ (i_A)_*}_{=0}(a)= 0 
		\]
		$(i_B,j_A)_*$ ist ein Isomorphismus, also $(l_B)_*(b)=0$. $\ker (l_B)_* = \im (j_B)_*$. Damit folgt $\exists x \in H_n(A \cap B;R) : (j_B)_*(x)=b$. Dann gilt
		\[
			(i_A)_* \enbrace[\big]{a - (j_A)-*(x)} = (i_A)_*(a) - i_B \enbrace[\big]{(j_B)_* (x)} = (i_A)_*(a) - (i_B)_* (b) = 0  
		\]
		Mit $\ker (i_A)_*  =\im \rand{X,A}_n$ und $(i_B,j_A)_*$ ist Isomorphismus, folgt $\exists y \in H_{n+1}(B,A\cap B;R)$ mit 
		\[
			\rand{X,A}_n \circ  \enbrace*{i_B,j_A}_* (y) = a - (j_A)_*(x) 
		\]
		Es folgt für $z := x + \partial_n^{B,A \cap B}(y)$.
		\begin{align*}
			(j_B)_*(z) & =(j_B)_*(x) = b \\
			(j_A)_*(z) &= (j_A)_*(x) + (j_A)_* \enbrace*{\rand{B,A \cap B}(y)}_n = (j_A)_*(x) + \rand{X,A}_n \circ  \enbrace*{i_B, j_A}_*(y) = (j_A)_*(x) + a -
			(j_A)_*(x) = a  
		\end{align*}
		\item["$\subseteq$":] Sei $x \in H_n(A\cap B;R)$. Dann gilt
		\begin{align*}
			\enbrace*{(i_A)_* - (i_B)_*} \begin{pmatrix}
				(j_A)_* \\ (j_B)_*
			\end{pmatrix}  (x) = (i_A)_* \circ (j_A)_* (x) - (i_B)_* \circ (j_B)_*(x) = (i_A \circ j_A)_* (x) - (i_B \circ j_B)_* (x) = 0
		\end{align*}
		da $i_A \circ j_A = i_B \circ j_B$.
	\end{description}
	\item Wir zeigen nun $\im (i_A)_* - (i_B)_* = \ker \partial_n$:
	\begin{description}
		\item["$\supseteq$":] Sei $v \in \ker \partial_n$. \todo{Rest hinzufügen}
		\item["$\subseteq$":] Sei $\begin{psmallmatrix} a \\ b \end{psmallmatrix} \in \bigoplus\limits_{H_n(B)}^{H_n(A)}$. Dann gilt
		\[
			\partial_n \enbrace*{(i_A)_* - (i_B)_*}  \begin{psmallmatrix} a \\ b \end{psmallmatrix} = \partial_n \enbrace[\big]{(i_A)_*(a)- (i_B)_*(b)}  =0
		\]
	\end{description}
	\item Noch zu zeigen: $\im \partial_n = \ker \begin{psmallmatrix} (j_A)_* \\ (j_B)_* \end{psmallmatrix}$.
	
	\emph{Übung!}
\end{itemize}
% subsection 92 (end)

\subsection[Bemerkung zur Mayer-Vietoris-Folge]{Bemerkung} % (fold)
\label{sub:93}
Für die Mayer-Vietoris-Folge müssen $A$ und $B$ nicht notwendig offen sein. Es genügt, dass die Inklusion $(B,A \cap B) \hookrightarrow (X,A)$ einen Isomorphismus
$H_*(B,A \cap B;R) \to H_*(X,A;R)$ induziert.
% subsection 93 (end)

\subsection[Satz: Homologie der Sphäre $S^n$]{Satz} % (fold)
\label{sub:94}
\begin{equation*}
	H_k \enbrace*{S^n ;R} \cong \begin{cases}
		R \oplus R , &\text{ falls }k=n=0\\
		R, &\text{ falls } k=0, n \not= 0 \\
		R, &\text{ falls } k=n\not= 0 \\
		0, &\text{ sonst}
	\end{cases} \label{eq:94} \tag{\#}
\end{equation*}
% subsection 94 (end)

\subsection[Definition: Reduzierte Homologie]{Definition} % (fold)
\label{sub:95}
Sei $X$ ein nichtleerer topologischer Raum. Sei $p : X  \to \set{\mathrm{pt}}$ die Projektion auf den Ein-Punkt-Raum. die \Index{reduzierte singuläre Homologie} von $X$ ist definiert
durch 
\[
	\tilde{H}_k(X;R) := \ker \enbrace[\Big]{ p_* : H_k(X;R) \to H_k(\set{\mathrm{pt}};R )} 
\]
\minisec{Bemerkung}
\begin{itemize}
	\item Es gilt $\tilde{H}_k(X;R) =H_k(X;R)$ für $k>0$ und $H_0(X;R) = R \oplus \tilde{H}_0(X;R)$
	\item Es ist nicht schwer die Eigenschaften von singulärer Homologie auf die reduzierte singuläre Homologie zu übertragen. Insbesondere ist $\tilde{H}_*$ 
	homotopieinvariant und es gibt eine Mayer-Vietoris-Folge.
	\item Es ist $\tilde{H}_k \enbrace*{\set{\mathrm{pt}}}=0$.
\end{itemize}
% subsection 95 (end)

\subsection[Bemerkung: Die Aussage von Satz \ref{sub:94} in reduzierter Homologie]{Bemerkung} % (fold)
\label{sub:96}
In reduzierter singulärer Homologie wird \eqref{eq:94} zu
\begin{equation*}
	\tilde{H}_k(S^n;R) \cong \begin{cases}
		R, &\text{ falls }k=n\\
		0 , &\text{ sonst}
	\end{cases} \label{eq:96}\tag{\#\#}
\end{equation*}
% subsection 96 (end)

\subsection[Beweis von \textbf{[\#\#]}]{Beweis von \eqref{eq:96}} % (fold)
\label{sub:97}
Induktion nach $n$:
\begin{description}
	\item[Induktionsanfang: $n=0$:] Es gilt
	\begin{align*}
		H_k(S^0;R) \stackrel{\xleftarrow{\varphi}}{\cong} H_k \enbrace[\big]{\set{\mathrm{pt}};R } \oplus H_k \enbrace[\big]{\set{\mathrm{pt}};R } = \begin{cases}
			R \oplus R, &\text{ falls }k=0\\
			0 ,&\text{ sonst}
		\end{cases}
	\end{align*} 
	Sei $i_- : \set{\mathrm{pt}} \hookrightarrow S^0$, $i_+ : \set{\mathrm{pt}} \hookrightarrow S^0$. $\im i_{pm} = \set{\pm 1}$.
	Der Isomorphismus $\varphi$ ist gegeben durch $\varphi = (i_-)_* + (i_+)_*$. Für
	$(a,b) \in H_k \enbrace*{\set{-1};R}\oplus H_k \enbrace*{\set{+1};R}$ gilt dann
	\[
		p_* \enbrace[\big]{i_-(a), i_+(b)} = (p \circ i_-)_*(a) + (p \circ i_+)(b) = \id_*(a) + \id_*(b)= a+b
	\]
	Also $\ker p_* = \set[(a,-a)]{a \in H_k \enbrace[\big]{\set{\mathrm{pt}};R}}$. Also 
	\[
		\tilde{H}_k(S^0) = \begin{cases}
			0, &\text{ falls }k \not= 0\\
			R ,&\text{ falls }k=0
		\end{cases}
	\]
	\item[Induktionsschritt: $n-1  \mapsto n$:] Sei $D_n^+ = S^n \setminus \set{(-1,0,\ldots ,0)}$ und $D_n^- = S^n \setminus \set{1,0,\ldots ,0}$. Dann gilt
	$S^n = D_n^+ \cup D_n^-$ und $D_n^+, D_n^- \subseteq S^n$ offen und $\D_n^+ \simeq \set{\mathrm{pt}}$ und $D_n^- \simeq \set{\mathrm{pt}}$
	Weiter gilt 
	\[
		D_n^+ \cap D_n^- = \set[x=(x_0, \ldots , x_n) \in S^n]{x_0 \not\in \set{\pm 1}} \simeq S^{n-1} = \set{x=(0, x_1, \ldots , x_n) \in S^n}.
	\]
	Betrachte nun die Mayer-Vietoris-Folge zu $S^n = D_n^+ \cup D_n^-$: \marginnote{"$R$" aus Platzgründen weggelassen}
	\[
		\begin{tikzcd}
			\underset{=0}{\tilde{H}_k(D_n^+)} \oplus \underset{=0}{\tilde{H}_k(D_n^-)} \rar & 
			\tilde{H}_k(S^n) \rar["\partial_k", "\cong"'] &
			\tilde{H}_{k-1} \enbrace*{D_n^+ \cup D_n^-} \rar & 
			\underset{=0}{\tilde{H}_{k-1}(D_n^+)} \oplus \underset{=0}{\tilde{H}_{k-1}(D_n^-)} \\
			 & & \tilde{H}_{k-1}(S^{n-1}) \uar["\cong", "(\iota_*)"']& 
		\end{tikzcd}
	\]
	Es folgt $\tilde{H}_k(S^n;R) \cong \tilde{H}_{k-1}(S^{n-1};R) \stackrel{\text{I.A.}}{=} \begin{cases}
		R, &\text{ falls }k=1\\
		0, &\text{ sonst}
	\end{cases}$ \bewende
\end{description}
% subsection 97 (end)
% section 9 (end)
\newpage

\section{Anwendungen} % (fold)
\label{sec:10}

\subsection[Satz: Invarianz der Dimension]{Satz (Invarianz der Dimension)} % (fold)
\label{sub:101}
Sind $\mathds{R}^n$ und $\mathds{R}^m$ homöomorph, so gilt $n=m$. \todo{RevChap10}
\minisec{Beweis}
Sei $h : \mathds{R}^n \to \mathds{R}^m$ ein Homöomorphismus. Sei $x_0 \in \mathds{R}^n$ beliebig. Dann erhalten wir auch einen Homöomorphismus 
$k := h\big|_{\mathds{R}^n\setminus \set{x_0}} : \mathds{R}^n \setminus \set{x_0} \to \mathds{R}^m  \setminus \set{h(x_0)}$. Nun ist $\mathds{R}^n \setminus \set{x_0}$
homotopieäquivalent zu $S^{n-1}$ und $\mathds{R}^{m}\setminus\set{h(x_0)}$ ist homotopieäquivalent zu $S^{m-1}$. Da $k$ ein Homöomorphismus ist, ist $k_*$ ein Isomorphismus.
Wir erhalten
\begin{align*}
	H_*(S^{n-1}) \cong H_* \enbrace[\big]{\mathds{R}^n \setminus \set{x_0}}  \cong H_* \enbrace[\big]{\mathds{R}^m \setminus \set{h(x_0)}} \cong H_*(S^{m-1})
\end{align*}
Es ist aber $H_*(S^{n-1}) \cong H_*(S^{m-1})$ genau dann, wenn $n=m$. \bewende
% subsection 101 (end)

\subsection{Fixpunktsatz von Brouwer} % (fold)
\label{sub:102}
Jede stetige Abbildung $f : D^n \to D^n$ besitzt einen Fixpunkt.
\minisec{Beweis}
Durch Widerspruch: Angenommen es gibt $f : D^n \to D^n$ ohne Fixpunkt. Aus dem letzten Semester wissen wir, dass es dann eine stetige Abbildung $F : D^n \to S^{n-1}$ gibt 
mit $F|_{S^{n-1}} = \id_{S^{n-1}}$. In Homologie erhalten wir folgendes kommutatives Diagramm \marginnote{\oE: $n\ge 2$}
\[
	\begin{tikzcd}
		0 = H_{n-1}(D^n) \rar["(F)_*"] & H_{n-1}(S^{n-1}) \cong \mathds{Z} \\
		H_{n-1}(S^{n-1}) \uar["i_*"] \urar["(\id_{S^{n-1}})_* = \id_{H_{n-1}(S^{n-1})}"']
	\end{tikzcd}
\]
wobei $i : S^{n-1} \hookrightarrow D^n$ die Inklusion ist. Da $F \circ i = \id_{S^{n-1}}$ ist, gilt $F_* \circ i_* = \id_{H_{n-1}(S^{n-1})}$. \light \bewende
% subsection 102 (end)
\section*{Abbildungsgrad}

\subsection[Definition: Abbildungsgrad]{Definition} % (fold)
\label{sub:103}
Sei $f : S^n \to S^n$ eine stetige Abbildung. Da $\tilde{H}_n(S^n) \cong Z$ ist, gibt es eine ganze Zahl $d(f)$, so dass $f_*(x)= d(f) \cdot x$ für alle 
$x \in \tilde{H}_n(S^n)$. Diese Zahl heißt der \Index{Abbildungsgrad} von $f$. 
% subsection 103 (end)

\subsection[Proposition: Einfache Eigenschaften des Abbildungsgrades]{Proposition} % (fold)
\label{sub:104}
\begin{enumerate}[a)]
	\item $d(\id_{S^n}) = 1$.
	\item Sind $f$ und $g$ homotop, so gilt $d(f)=d(g)$.
	\item $d(f \circ g) = d(f) \cdot d(g)$.
\end{enumerate}
\minisec{Beweis}
\begin{enumerate}[a)]
	\item $(\id_{S^n})_*(x)= \id(x)= x$.
	\item Sind $f$ und $g$ homotop, so gilt $f_*=g_*$.
	\item Für alle $x \in \tilde{H}_n(S^n)$ gilt: $d(f \circ g) \cdot x = (f \circ g)_* (x) = f_* \circ g_*(x)= d(f) \cdot d(g) \cdot x$. Da 
	$\tilde{H}_n(S^n) \cong \mathds{Z}$ folgt $d(f \circ g) = d(f) \cdot d(g)$. \bewende
\end{enumerate}
% subsection 104 (end)

\subsection[Beispiel: Abbildungsgrad der Spiegelung]{Beispiel} % (fold)
\label{sub:105}
Sei $i \in \set{0, \ldots ,n}$ und $R_i : S^n \to S^n$ die Spiegelung $R_i(x_0, \ldots , x_n) = (x_0, \ldots , - x_i, \ldots , x_n)$. Dann ist $d(R_i) = -1$.
\minisec{Beweis}
Sei $f : S^n \to S^n$ der Homöomorphismus, der die $0$-te und die $i$-te Koordinate vertauscht: $f(x_0, \ldots , x_n) = (x_i, \ldots , x_0, \ldots , x_m)$. Dann ist
$R_i = f \circ R_0 \circ f$. Da $f$ ein Homöomorphismus ist, ist $f_*$ ein Isomorphismus und daher $d(f)$ invertierbar, also $d(f) \in \set{\pm 1}$ mit 
$d(R_i) = d(f) \cdot d(R_0) \cdot d(f)$ folgt $d(R_0)=d(R_i)$. Es genügt also $R_0$ zu betrachten.

Wir zeigen $d(R_0)= -1$ durch Induktion nach $n$.
\begin{description}
	\item[Induktionsanfang:] Sei $i_- : \set{\mathrm{pt}} \to S^0$ mit $\im i_- = \set{-1}$ und $i_+ : \set{\mathrm{pt}} \to S^0 $ mit $\im i_+ = \set{+1}$. Dann ist
	\begin{align*}
		H_0(S^0) &= \set[(i_-)_*(a)+ (i_+)_*(b)]{a,b \in H_0(\set{\mathrm{pt}})} \qquad \text{und}\\
		\tilde{H}_=(S^0) &= \set[(i_-)_*(a)- (i_+)_*(a)]{a \in H_0(\set{\mathrm{pt}})}
	\end{align*}
	Nun ist
	\begin{align*}
		(R_0)_* \enbrace*{(i_-)_*(a)- (i_+)_*(a)} = (R_0)_* (i_-)_* (a) - (R_0)_* (i_+)_* (a) &= (R_0 \circ i_-)_*(a) - (R_0 \circ i_+)_*(a)\\ 
		&= (i_+)_*(a) - (i_-)_*(a) \\ 
		&= - \enbrace[\big]{(i_-)_*(a) + (i_+)_*(a)} 
	\end{align*}
	Also folgt $d(R_0)=-1$.
	\item[Induktionsschritt:] Sei 
	\begin{align*}
		D_+^n &= \set[(x_0, \ldots , x_n) \in S^n]{x_n \not= -1} \\
		D_-^n &= \set[(x_0, \ldots , x_m) \in S^n]{x_n \not= +1}  
	\end{align*}
	Wir wissen schon, dass die Randabbildung aus der dazugehörigen Mayer-Vietoris-Folge $\tilde{H}_n(S^1) \xrightarrow{\partial_n} \tilde{H}_{n-1}(D^n_+ \cap D_-^n)$ ein 
	Isomorphismus ist. Da die Randabbildung natürlich ist, erhalten wir ein kommutatives Diagramm\todo{fertig stellen}
	\[
		\begin{tikzcd}[row sep=3em, column sep=4em]
			\tilde{H}_n(S^n) \rar["\partial_n", "\cong"'] \dar["(R_0)_*"]
			& \tilde{H}_{n}(D^n_- \cap D^n_+) \dar["\enbrace*{R_0\big|_{(D^n_- \cap D^n_+)}}_*"]  
			& \tilde{H}_{n-1}(S^{n-1}) \lar["i_*", "\cong"'] \\
			\tilde{H}_n(S^n) \rar["\partial_n", "\cong"']
			& \tilde{H}_{n+1}(D^n_- \cap D^n_+) 
			& H_{n-1}(S^{n-1})
		\end{tikzcd}
	\]
	dabei ist $i: S^{n-1} \to D^n_- \cap D_+^n$ die Homotopieäquivalenz $i(x_0, \ldots , x_{n-1}) = (x_0, \ldots , x_{n-1},0)$.
	Nach Induktionsannahme ist $\enbrace*{R_0\big|_{S^{n-1}}}_*(y)= -y$ für alle $y \in \tilde{H}_{n-1}(S^{n-1})$. Für $x \in \tilde{H}_n(S^n)$ folgt
	\begin{align*}
		(R_0)_+(x) &= (\partial_n) ^{-1} \circ  (i_*) \circ \enbrace*{R_0\big|_{S^{n-1}}}_* \enbrace[\big]{(i_* ^{-1} \circ \partial_n) \cdot (x)} \\
		&= (\partial_n) ^{-1} \circ (i_*) \enbrace*{- i_*^{-1} \enbrace*{\partial_n (x)} } \\
		&= - \partial_n ^{-1} \circ (i_*) \circ (i_*)^{-1} \circ \partial_n(x)= -x 
	\end{align*}
	Also $d(R_0)=-1$. \bewende
\end{description}
% subsection 105 (end)

\subsection[Beispiel: Grad der Punktspiegelung]{Beispiel} % (fold)
\label{sub:106}
Der Grad der Punktspiegelung $R : S^n \to S^n$ $x  \mapsto -x$ ist $(-1)^{n+1}$.
\minisec{Beweis}
$d(R)= d(R_0 \circ R_1 \circ \ldots \circ R_n) = d(R_0) \cdot \ldots \cdot d(R_n)= (-1)^{n+1}$. \bewende
% subsection 106 (end)

\subsection[Erinnerung: Vektorfelder]{Erinnerung} % (fold)
\label{sub:107}
Sei $M$ eine $C^\infty$-Mannigfaltigkeit. Ein \Index{Vektorfeld} auf $M$ ist eine stetige Abbildung $v : M \to \Tmap M$ mit $v(x) \in \Tmap_x M$ für alle $x \in M$.
Für $M=S^n$ entspricht ein Vektorfeld genau einer stetigen Abbildung $v : S^n \to \mathds{R}^{n+1}$ mit $\sprod*{v(x)}{x} =0$ für alle $x \in S^n$.
\missingfigure{$S^2$ mit Vektorfeld} 
% subsection 107 (end)

\subsection[Satz: Es gibt ein Vektorfeld ohne Nullstellen auf $S^n \iff n$ ist ungerade]{Satz} % (fold)
\label{sub:108}
Es gibt genau dann ein Vektorfeld ohne Nullstellen auf $S^n$, wenn $n$ ungerade ist.
\minisec{Beweis}
Ist $n$ ungerade, so ist $S^{n} \subseteq \mathds{R}^{n+1} = \mathds{C}^{k}$ mit $k=\frac{n+1}{2}$. Punkte in $S^n$ sind dann genau $k$-Tupel $z=(z_1, \ldots , z_k)$ mit
$\abs*{z_1}^2 + \ldots + \abs*{z_k}^2 = 1$. Ein Vektorfeld entspricht dann einer stetigen Abbildung $v : S^n \to \mathds{C}^k$, so dass für das kanonische komplexe 
Skalarprodukt $\sprod*{\cdot }{\cdot}_\mathds{C} $ auf $\mathds{C}^k$ gilt
\[
	\sprod*{z}{v(z)}_\mathds{C} \in i \mathds{R} 
\]
Definiere nun $v : S^n \to \mathds{C}^k$ durch $v(z) := i z$. Dann $i z\not= 0$ für alle $z \in S^n$ und $\sprod*{z}{i z}= -i \norm[2]{z} \in i \mathds{R}$. Sei umgekehrt
$v : S^n \to \mathds{R}^{n+1}$ ein Vektorfeld ohne Nullstellen. Zu zeigen: $d( \text{Punktspiegelung }R)=1$ (Da $d(R)=(-1)^{n+1}$ folgt dann die Behauptung).
Zeige dazu: $R$ ist homotop zu $\id$. Seien $H, K : S^n \times [0,1] \to S^n$ definiert durch 
\[
	H(x,t) := \frac{t x + (1-t)v(x)}{\norm{tx + (1-t)v(x)} } \qquad \text{bzw.}\qquad K(x,t) :=  \frac{-t x + (1-t)v(x)}{\norm{-tx + (1-t)v(x)}}
\]
Da $\sprod*{v(x)}{x}=0$ folgt für alle $x$ und $v(x)\not=0$ ist, folgt 
\begin{align*}
	\norm{t x + (1-t)v(x)}^2 &= \norm{tx}^2 + \norm{(1-t)v(x)}^2 \not=0   \\
	\norm{-t x + (1-t)v(x)}^2 &= \norm{-tx}^2 + \norm{(1-t)v(x)}^2 \not=0
\end{align*}
Also sind $H$ und $K$ wohldefiniert. $H$ ist Homotopie zwischen $\id$ und $x \mapsto \frac{v(x)}{\norm{v(x)} } $ und $K$ ist Homotopie zwischen $-\id$ und 
$x \mapsto \frac{v(x)}{\norm{v(x)} } $. \bewende
% subsection 108 (end)

\subsection{Satz (Jordanscher Kurvensatz)} % (fold)
\label{sub:109}
Sei $f \colon S^1 \to \mathds{R}^2$ eine stetige, injektive Abbildung. Dann hat $\mathds{R}^2 \setminus f(S^1)$ genau zwei Wegzusammenhangskomponenten. Eine davon ist 
beschränkt, die andere nicht.
\minisec{Beweis (mit \ref{sub:1010})}
Wir können $f$ auch als injektive Abbildung $f \colon S^1 \to \mathds{R}^2 \subseteq S^2 = \mathds{R}^1 \cup \set{\infty}$ auffassen. Nun ist 
$\tilde{H}_0 \enbrace[\big]{S^2 \setminus f(S^1)} =\mathds{Z}$. Also $H_0 \enbrace[\big]{S^2 \setminus f(S^1)}= \mathds{Z} \oplus \mathds{Z}$. Daher besteht
$S^2 \setminus f(S^1)$ aus zwei Wegzusammenhangskomponenten $U_1$ und $U_2$. Sei \oE $\infty \in U_2$. Dann sind $U_1$ und $U_2 \setminus \set{\infty}$ die
Wegzusammenhangskomponenten von $\mathds{R}^2 \setminus f(S^1)$. $U_1$ ist beschränkt und $U_2$ ist unbeschränkt. \bewende \smallskip \\
Warum ist $U_2 \setminus \set{\infty}$ noch wegzusammenhängend?
\minisec{Antwort}
Seien $x,y \in U_2$, dann gibt es eine Umgebung $V$ von $\infty \in S^2$ mit $V \setminus \set{\infty} \equiv D^2 \setminus \set{0} $,  $x,y \not\in V$ und
$f(S^1) \cap V = \emptyset$
% subsection 109 (end)

\subsection{Satz} % (fold)
\label{sub:1010}
\begin{enumerate}[a)]
	\item Sei $f \colon D^k \to S^n$ eine stetige, injektive Abbildung mit $0 \le k < n$. Dann gilt $\tilde{H}_i\enbrace[\big]{S^n \setminus f(D^k)} =0$ für alle $i$.
	\item Sei $f \colon S^k \to S^n$ eine stetige, injektive Abbildung mit $0 \le k < n$. Dann gilt 
	\[
		\tilde{H}_i \enbrace[\big]{S^n \setminus f(S^k)} = \begin{cases}
			\mathds{Z}, &\text{ falls }i=n-k-1\\
			0, &\text{ sonst}
		\end{cases}
	\]
\end{enumerate}
\minisec{Beweis (mit \ref{sub:1011})}
\begin{enumerate}[a)]
	\item Durch Induktion nach $k$. Für $k=0$ ist $S^n \setminus f(D^0)$ homöomorph zu $\mathds{R}^n$ und die Behauptung folgt.
	
	Induktionsschritt $k \mapsto k+1$: Da $D^{k+1}$ und $[0,1]^{k+1}$ homöomorph sind, können wir $D^{k+1}$ durch $[0,1]^{k+1}$ ersetzen. Sei 
	$x \in \tilde{H}_i \enbrace*{S^n \setminus f \enbrace*{[0,1]^{k+1}}}$. Es ist
	\[
		S^n \setminus f \enbrace*{\set{\nicefrac{1}{2}} \times [0,1]^{k}} = S^n \setminus f \enbrace*{\benbrace*{0,\nicefrac{1}{2}}  \times [0,1]^k} \cup S^n \setminus f \enbrace*{\benbrace*{\nicefrac{1}{2},1} \times [0,1]^k}   
	\]
	eine Vereinigung von offenen Mengen. Weiter ist 
	\begin{align*}
		\enbrace*{S^n \setminus f \enbrace*{\benbrace*{0,\nicefrac{1}{2} } \times [0,1]^k} } \cap \enbrace*{S^n \setminus f \enbrace*{\benbrace*{\nicefrac{1}{2} ,1} \times [0,1]^k}} =   S^n \setminus f \enbrace*{[0,1]^{k+1}} 
	\end{align*}
	Nun benutzten wir die zugehörige Mayer-Vietoris-Folge: \todo{schöner machen}
	\[
		\begin{tikzcd}[column sep=1.3em]
			\underset{=0 \text{ nach I.A.}}{\tilde{H}_{i+1} \enbrace*{S^n \setminus f \enbrace*{\set{\nicefrac{1}{2}}  \times [0,1]^k}}} \rar &
			\tilde{H}_i \enbrace*{S^n \setminus f \enbrace*{[0,1]^{k+1}} }  \dar["\cong"] & \\
			& \bigoplus\limits_{\textstyle\tilde{H}_i \enbrace*{S^n \setminus f \enbrace[\big]{[\nicefrac{1}{2},1] \times [0,1]^k} } }
			^{\textstyle\tilde{H}_i \enbrace*{S^n \setminus f \enbrace[\big]{[0,\nicefrac{1}{2}] \times [0,1]^k} } } \rar &
			\underset{=0 \text{ nach I.A.}}{\tilde{H}_i \enbrace*{S^n \setminus f \enbrace*{\set{\nicefrac{1}{2}} \times [0,1]^k } } }
		\end{tikzcd}
	\]
	Angenommen $x \not=0$. Dann ist das Bild von $x$ in $\tilde{H}_i \enbrace*{S^n \setminus f \enbrace*{[0,\nicefrac{1}{2}] \times [0,1]^k} }$ ungleich Null oder in
	$\tilde{H}_i \enbrace*{S^n \setminus f \enbrace*{[\nicefrac{1}{2},1] \times [0,1]^k}}$ ungleich Null. Indem wir dieses Argument iterieren erhalten wir eine Folge
	von Intervallen
	\[
		[0,1] \supseteq I_1 \supseteq I_2 \supseteq \ldots \text{ mit } \bigcap_n I_n = \set{t} 
	\]
	und das Bild von $x \in \tilde{H}_i \enbrace*{S^n \setminus f \enbrace*{[0,1]^{k+1}}}$ unter der von der Inklusion 
	$i_a \colon S^n \setminus f \enbrace*{[0,1]^{k+1}} \to S^n \setminus f \enbrace*{I_a \times [0,1]^k}$ induzierten Abbildung 
	$(i_a)_* \colon \tilde{H}_i \enbrace*{S^n \setminus f \enbrace*{[0,1]^{k+1}}} \to \tilde{H}_i \enbrace*{S^n \setminus f \enbrace*{I_a \times [0,1]^k} } $ ungleich Null 
	ist. Nun ist aber 
	\[
		\bigcup_a \enbrace*{S^n \setminus f \enbrace*{I_a \times [0,1]^k}} = S^n \setminus f \enbrace*{\set{t} \times [0,1]^k }  
	\]
	und wieder nach Induktionsannahme das Bild von $x$ in $\tilde{H}_i \enbrace*{S^n \setminus f \enbrace*{\set{t} \times [0,1]^k}} =0$ Null. Nach Proposition 
	\ref{sub:1011}(2) muss $x$ aber schon in einer $\tilde{H}_i \enbrace*{S^n \setminus f \enbrace*{I_a \times [0,1]^k} }$ trivial sein. \light
	\item Durch Induktion nach $k$. Für $k=0$ ist $S^n \setminus f(S^0) \cong \mathds{R}^n \setminus \set{0}$. Da $\mathds{R}^n \setminus \set{0} \simeq S^{n-1}$
	folgt für $k=0$ die Behauptung.
	
	Induktionsschritt $k-1 \mapsto k$: Seien $D^k_+ := \set[(x_0,\ldots ,x_k)]{x_k \ge 0}$ und $D^k_- := \set[(x_0, \ldots , x_k)]{x_k \le 0}$. Dann ist 
	$D^k_+ \cap D^k_- = \set[(x_0, \ldots ,x_k)]{x_k=0} \cong S^{k-1}$. Wieder ist 
	\[
		S^n \setminus f \enbrace*{D^k_+ \cap D^k_-} \enbrace*{S^n \setminus f(D^k_+)} \cup \enbrace*{S^n \setminus f(D^k_-)}  
	\]
	eine offene Vereinigung mit $\enbrace*{S^n \setminus f(D^k_+)} \cap (S^n \setminus f(D^k_-)) = S^n \setminus f(S^k)$. Wegen a) ist die Randabbildung in der zugehörigen
	Mayer-Vietoris-Folge
	\[
		\tilde{H}_i \enbrace*{S^n \setminus f \enbrace*{D^k_+ \cap D^k_-} } \xrightarrow{\partial_i \cong} \tilde{H}_{i-1} \enbrace*{S^n \setminus f(S^k)}   
	\]
	ein Isomorphismus. Die Behauptung folgt per Induktion. \bewende
\end{enumerate}
% subsection 1010 (end)

\subsection{Proposition} % (fold)
\label{sub:1011}
Seien $U_1 \subseteq U_2 \subseteq \ldots  \subseteq X$ offen mit $X = \bigcup_{i=1}^\infty U_i$. Seien $i_{a,b} \colon U_a \hookrightarrow U_b$ für $a < b$ und 
$i_a \colon U_a \hookrightarrow X$ die Inklusionen.
\begin{enumerate}[(1)]
	\item Für jedes $x \in H_k(X;R)$ gibt es $a >0$ und $x_a \in H_k(U_a;R)$ mit $(i_a)_*(x_a)=x$.
	\item Ist $x_a \in H_k(U_a;R)$ mit $(i_a)_*(x)=0$, so gibt es $b >a$ mit $(i_{a,b})_*(x_a)=0$.
\end{enumerate}
\minisec{Bemerkung}
Diese Proposition gilt genauso für reduzierte Homologie.
\minisec{Beweis}
Sei $\sigma \colon \abs*{\Delta^n} \to X$ ein singulärer Simplex in $X$. Dann ist $\sigma ^{-1}(U_a)$, $a \in \mathds{N}$ eine offene Überdeckung von $\abs*{\Delta^n}$. Da 
$\abs*{\Delta^n}$ kompakt ist, gibt es $a$ mit $\abs*{\Delta^n} \subseteq \sigma ^{-1}(U_a)$, also $\sigma \enbrace*{\abs*{\Delta^n}} \subseteq U_a$. Da Elemente  von
$C_n(X;R)$ endliche $R$-Linearkombinationen von singulären Simplizes sind, folgt 
\[
	C_n(X;R) = \bigcup_a C_n(U_a;R).
\]
Damit ergeben sich leicht (1) und (2). \bewende
% subsection 1011 (end)
% section 10 (end)
\newpage

\section{\CW-Komplexe} % (fold)
\label{sec:11}

\subsection[Definition: Pushout]{Definition} % (fold)
\label{sub:111}
Ein kommutatives Diagramm von topologischen Räumen der Form \todo{RevChap11}
\begin{equation*}
	\begin{tikzcd}
		A \rar["f"] \dar["g"] & X \dar["\overline{g} "] \\
		Y \rar["\overline{f}"] & Z
	\end{tikzcd} \label{eq:111} \tag{\#}
\end{equation*}
\noindent\begin{minipage}{0.65\textwidth}
	heißt ein \Index{Pushout}, falls es folgende universelle Eigenschaft hat: \smallskip\\
	\centering{$\forall \hat{f} \colon X \to \hat{Z}$, $\hat{g}: Y \to \hat{Z}$ mit 
	$\hat{f} \circ f = \hat{g} \circ g$ gilt: 
	\[
		\exists ! \varphi \colon Z \to \hat{Z} \text{ mit } \hat{f}= \varphi \circ \overline{g} , \hat{g} = \varphi \circ \overline{f}.
	\]}
\end{minipage}
\hfill \begin{minipage}{0.3\textwidth}
	\(
		\begin{tikzcd}
			A \rar["f"] \dar["g"] & X \dar["\overline{g}"] \arrow[ddr, "\hat{f}", bend left, Firebrick3] \\
			Y \rar["\overline{f}"] \arrow[drr, "\hat{g}"',bend right, Firebrick3] & Z  \drar["\varphi", DodgerBlue3] \\
			& & \hat{Z}
		\end{tikzcd}
	\)
\end{minipage}\\
Wir sagen dann auch: \eqref{eq:111} ist der Pushout von 
\(
	\begin{tikzcd}[sep=small]
		A \rar["f"] \dar["g"']& X \\
		Y
	\end{tikzcd}
\)
\minisec{Bemerkung}
Jedes Diagramm $\begin{tikzcd}[sep=small]
		A \rar["f"] \dar["g"']& X \\
		Y
	\end{tikzcd}$ lässt sich zu einem Pushout vervollständigen:
Betrachte dazu den Raum $X \cup_A Y := \nicefrac{X \coprod Y}{f(a) \sim g(a) \forall a \in A}$ mit der Quotiententopologie. Sind $\overline{f} : Y \to X \cup_A Y $  und 
$\overline{g} : X \to X \cup_A Y$ die von den Inklusionen $X \hookrightarrow X \coprod Y$ und $Y \hookrightarrow X \coprod Y$ induzierten Abbildungen. So ist 
\[
	\begin{tikzcd}
		A \rar["f"] \dar["g"'] & X \dar["\overline{g}"] \\
		Y \rar["\overline{f}"] & X \cup_A Y
	\end{tikzcd}
\]
Existenz und Eindeutigkeit
aus der universellen Eigenschaft können benutzt werden um zu zeigen, dass der Pushout eindeutig bis auf kanonischen Homöomorphismus ist.
% subsection 111 (end)

\subsection[Definition: \CW-Komplex]{Definition} % (fold)
\label{sub:112}
Ein \bet{\CW-Komplex}\index{CW-Komplex@\CW-Komplex} ist ein topologischer Raum $X$ zusammen mit einer Filtrierung der Unterräume von $X$: 
\[
	\emptyset = X^{(-1)} \subseteq X^{(0)} \subseteq X^{(1)} \subseteq X^{(2)} \subseteq  \ldots \subseteq X^{(n)} \subseteq X^{(n+1)} \subseteq \ldots \subseteq X
\]
sodass die folgenden zwei Eigenschaften erfüllt sind:
\begin{enumerate}[(i)]
	\item \Index{Zellstruktur}: Zu jedem $n \in N$ gibt einen Pushout von topologischen Räumen 
	\[
		\begin{tikzcd}[row sep = large, column sep = 5.4em]
			\smashoperator{\coprod\limits_{i \in I^{(n)}}} S^{n-1} \rar["\coprod_{i \in I^{(n)}} q_i^{(n)}"] \dar["\coprod_{i \in I^{(n)}} j_i"', hook] 
			& X^{(n-1)} \dar[hook, "k^{(n)}"] \\
			\smashoperator{\coprod\limits_{i \in I^{(n)}}} D^n \rar["\coprod_{i \in I^{(n)}} Q_i^{(n)}"] & X^{(n)}
		\end{tikzcd}
	\]
	wobei $j_i \colon S^{n-1} \hookrightarrow D^n$ und $k_n \colon X^{(n-1)} \to X^{(n)}$ die Inklusionen sind.
	\item \Index{Schwache Topologie}: Es ist $X= \bigcup_{n \in \mathds{N}} X^{(n)}$ und $U \subseteq X$ ist genau dann offen, wenn $U \cap X^{(n)} \subseteq X^{(n)}$ für 
	alle $n$ offen ist.
\end{enumerate}
Eine Abbildung $f \colon X \to Y$ zwischen \CW-Komplexen heißt \Index{zellulär}, falls $f \enbrace*{X^{(n)}} \to Y^{(n)}$ für alle $n \in \mathds{N}$. 
% subsection 112 (end)

\subsection[Bemerkungen zur Definition von \CW-Komplexen]{Bemerkung} % (fold)
\label{sub:113}
\begin{enumerate}[(i)]
	\item Die Abbildungen $q_i^{(n)}$ und $Q_i^{(n)}$ sind \emph{nicht} Teil der Struktur eines \CW-Komplex. Nur die Existenz von $q_i^{(n)}$ und $Q_i^{(n)}$ wird 
	gefordert --- sie ist nicht notwendig eindeutig. Hat man $q_i^{(n)}$ und $Q_i^{(n)}$ \emph{gewählt}, so heißen $q_i^{(n)}$ die \Index{anklebende Abbildung} der $i$-ten
	$n$-Zelle und $Q_i^{(n)}$ die \Index{charakteristische Abbildung} der $i$-ten $n$-Zelle.
	\item Sei $\mathring{D}^n = D^n \setminus S^{n-1}$ das Innere von $D^n$. Die Abbildungen $Q_i^{(n)}$ schränken sich zu einem Homöomorphismus 
	$\coprod_{i \in I^{(n)}} \mathring{D}^n \to X^{(n)} \setminus X^{(n-1)}$ ein. Insbesondere lässt ich $I^{(n)}$ mit der Menge der Wegzusammenhangskomponenten von 
	$X^{(n)} \setminus X^{(n-1)}$ identifizieren. Die Wegzusammenhangskomponenten heißen die \bet{offenen Zellen}\index{offene Zelle} von $X$. Damit ist jeder \CW-Komplex 
	die disjunkte Vereinigung seiner offenen Zellen.
	
	\ifxetexorluatex \faWarningSign \, \fi Achtung: Jede offene $n$-Zelle ist offen in $X^{(n)}$, aber nicht notwendig in $X$.
	\item Der Abschluss einer offenen Zelle $Q_i^{(n)}(\mathring{D}^n)$ ist $Q_i^{(n)}(D^n)$ und insbesondere kompakt. Die $Q_i^{(n)}(D^n)$ heißen die
	\bet{abgeschlossenen Zellen}\index{abgeschlossene Zelle} und sind als Abschluss von offenen Zellen unabhängig von der Wahl der $q_i^{(n)}$, $Q_i^{(n)}$.
	
	\ifxetexorluatex \faWarningSign \, \fi Achtung: Abgeschlossene Zellen sind Bilder von $D^n$ unter stetigen Abbildungen, aber nicht notwendig homöomorph zu $D^n$.
	\item Ein \CW-Komplex $X$ heißt \bet{endlich}\index{CW-Komplex@\CW-Komplex!endlicher}, wenn er nur aus endlich vielen Zellen besteht, also wenn 
	$I = \bigcup_{n \in \mathds{N} } I^{(n)}$ endlich ist. Insbesondere gibt es dann ein $n$ mit $X=X^{(n)}$.
	
	Die Dimension von $X$ ist die maximale Dimension von Zellen von $X$: $\dim X = N$, falls $X^{(N)}=X$, $X^{(N-1)} \subsetneq X$. Gibt es kein solches $N$ so
	$\dim X = \infty$.
\end{enumerate}
% subsection 113 (end)

\subsection[Beispiele für \CW-Komplexe]{Beispiele} % (fold)
\label{sub:114}
\begin{enumerate}[(i)]
	\item $S^n$ mit 
	\[
		(S^n)^{(k)} = \begin{cases}
			\emptyset, &\text{ falls }k=-1\\
			\set{(1,0,\ldots ,0)}, &\text{ falls } k=0, \ldots n-1 \\
			S^n , &\text{ falls } k \ge n 
		\end{cases}
	\]
	ist ein \CW-Komplex via $S^n \cong \nicefrac{D^n}{S^{n-1}}$
	\[
		\begin{tikzcd}
			S^{n-1} \rar \dar & \enbrace*{S^n}^{(n-1)}=\set{(1,0,\ldots ,0)}  \dar \\
			D^n \rar & S^n  
		\end{tikzcd}
	\]
	\item $S^n$ mit $(S^n)^{(k)} = \set{(x_0, \ldots ,x_k,0, \ldots ,0) \in S^n} $ ist ebenfalls eine \CW-Struktur auf $S^n$.
	\item Sei $K=(V,\Sigma)$ ein simplizialer Komplex. Es ist $\abs*{\Delta^n} \cong D^n$ und $\abs*{\partial \Delta^n} \cong S^{n-1}$. Sei 
	$\Sigma^{(n)} := \Sigma_0 \cup \Sigma_1 \cup \ldots \cup \Sigma_n$. Dann wir $\abs*{K}$ durch $\abs*{K}^{(n)} := \abs*{\enbrace*{V,\Sigma^{(n)}}}$ für $n \ge 0$ zu 
	einem \CW-Komplex. Ist $f \colon K \to L$ eine simpliziale Abbildung, so ist $\abs*{f} \colon \abs*{K} \to \abs*{L}$ eine zelluläre Abbildung.
	\item Der $n$-dimensionale reelle projektive Raum $\mathds{R}P^n$ kann definiert werden durch 
	\[
		\mathds{R}P^n := \nicefrac{\mathds{R}^{n+1} \setminus \set{0}}{x \sim \lambda x} = \nicefrac{S^n}{x \sim -x}
	\]
	Punkte in $\mathds{R}P^n$ sind Äquivalenzklassen von $(n+1)$-Tupeln reeller Zahlen und werden als die sogenannten homogenen Koordinaten $[x_0 : \ldots : x_n]$ 
	geschrieben. Es ist 
	\[
		\mathds{R}P^n = \set[{[x_0 : \ldots : x_n]}]{(x_0, \ldots ,x_n) \in S^n \rule{0cm}{1em}} 
		= \set[{[x_0 : \ldots : x_n]}]{(x_0, \ldots ,x_n) \in S^n, x_n \ge 0 \rule{0cm}{1em}}
	\]
	Wir erhalten einen Homöomorphismus $f^{(n)} \colon \nicefrac{D^n}{x \sim -x, x \in S^{n-1}} \to \mathds{R}P^n$ mit
	\[
		f \enbrace*{x_1, \ldots ,x_n} = \benbrace*{x_1 : \ldots : x_n : \sqrt{1- \sum\nolimits_{i=1}^{n} x_i^2}} 
	\]
	Ist $Q^{(n)} : D^n \to \mathds{R}P^n$ die Komposition von $f^{(n)}$ mit der Projektion $D^n \twoheadrightarrow \nicefrac{D^n}{x \sim -x, x \in S^{n-1}}$ und 
	$q^{(n)} \colon S^{n-1} \to \mathds{R}P^{n-1}$ die Projektion, so erhalten wir einen Pushout
	\[
		\begin{tikzcd}
			S^{n-1} \rar["q^{(n)}"] \dar[hook]& \mathds{R}P^{n-1} \dar[hook] \\
			D^n \rar["Q^{(n)}"] & \mathds{R}P^n
		\end{tikzcd}
	\]
	Induktiv erhalten wir einen \CW-Struktur auf $\mathds{R}P^n$ mit genau einer $k$-Zelle für $k=0,\ldots ,n$ und 
	\[
		\enbrace*{\mathds{R}P^n}^{(k)} \cong \mathds{R}P^k 
	\]
	\item Der $n$-dimensionale komplexe projektive Raum $\mathds{C}P^n$ kann definiert werden durch 
	\[
		\mathds{C}P^n := \nicefrac{\mathds{C}^{n+1} \setminus \set{0}}{z \sim \lambda z} = \nicefrac{S^{2n+1}}{z \sim \lambda z, \lambda \in S^1}
	\]
	Punkte in $\mathds{C}P^n$ sind Äquivalenzklassen von $(n+1)$-Tupeln komplexer Zahlen und werden als die sogenannten homogenen Koordinaten $[x_0 : \ldots : x_n]$ 
	geschrieben. Es ist 
	\[
		\mathds{C}P^n = \set[{[z_0 : \ldots : z_n]}]{(z_0, \ldots ,z_n) \in S^{2n+1} \rule{0cm}{1em}} 
		= \set[{[z_0 : \ldots : z_n]}]{(z_0, \ldots ,z_n) \in S^{2n+1}, z_n \ge 0 \rule{0cm}{1em}}
	\]
	Wir erhalten einen Homöomorphismus $f^{(n)} \colon \nicefrac{D^{2n+1}}{z \sim \lambda z, \lambda \in S^{2n}} \to \mathds{C}P^n$ mit
	\[
		f \enbrace*{z_1, \ldots ,z_n} = \benbrace*{z_1 : \ldots : z_n : \sqrt{1- \sum\nolimits_{i=1}^{n} \abs{z_i}^2}} 
	\]
	Ist $Q^{(2n)} : D^{2n} \to \mathds{C}P^n$ die Komposition von $f^{(n)}$ mit der Projektion $D^{2n} \twoheadrightarrow \nicefrac{D^{2n}}{z \sim \lambda z, z \in S^{2n}}$ und 
	$q^{(n)} \colon S^{2n+1} \to \mathds{C}P^{n-1}$ die Projektion, so erhalten wir einen Pushout
	\[
		\begin{tikzcd}
			S^{2n-1} \rar["q^{(n)}"] \dar[hook]& \mathds{C}P^{n-1} \dar[hook] \\
			D^{2n} \rar["Q^{(n)}"] & \mathds{C}P^n
		\end{tikzcd}
	\]
	Induktiv erhalten wir einen \CW-Struktur auf $\mathds{C}P^n$ mit genau einer $k$-Zelle für $k=0,2,\ldots ,2n$ und 
	\[
		\enbrace*{\mathds{C}P^n}^{(k)} \cong \mathds{C}P^{\floor*{\frac{k}{2}}} 
	\]
\end{enumerate}
% subsection 114 (end)

\subsection[Lemma: Kompaktheit von Teilmengen eines \CW-Komplexes]{Lemma} % (fold)
\label{sub:115}
Sei $X$ ein \CW-Komplex und $K \subseteq X$. Dann ist $K$ genau dann kompakt, wenn $K$ abgeschlossen ist und $K$ nur endlich viele offene Zellen von $X$ schneidet.
\minisec{Beweis}
Zunächst: Ist $S \subseteq X$ ein Unterraum der jede offene Zelle höchstens in einem Punkt schneidet, so ist $S$ diskret. Für $S_0 \subseteq S$ ist 
$S_0 \cap (X^{(n)} \setminus X^{(n-1)})$ abgeschlossen in $X^{(n)}$, da $S_0$ jede offene $n$-Zelle von $X$ höchstens in einem Punkt schneidet.  Daher ist 
$S_0 \cap X^{(n)} \subseteq X^{(n)}$ abgeschlossen für alle $n$. Damit ist auch $S_0 \subseteq X$ abgeschlossen und $S$ diskret.

Sei nun $K \subseteq X$ kompakt. Dann ist $K$ sicher abgeschlossen. Sei $S \subseteq K$ ein Teilraum der aus jeder offenen Zelle von $X$, die $K$ schneidet, genau einen 
Punkt enthält. Wegen der Vorüberlegung ist $S$ diskret. Da $S \subseteq K$ abgeschlossen und $K$ kompakt ist, ist $S$ auch kompakt. Damit ist $S$ endlich und $K$
schneidet nur endlich viele offene Zellen.

Ist nun umgekehrt $K$ abgeschlossen und schneidet nur endlich viele offene Zellen, so ist $K$ enthalten in einer endlichen Vereinigung von abgeschlossen Zellen von $X$.
Da diese abgeschlossen Zellen kompakt sind, ist auch diese Vereinigung kompakt, K ist also eine abgeschlossener Teilraum eines kompakten Teilraums und damit selbst kompakt.
\bewende
% subsection 115 (end)

\subsection[Korollar: Ein \CW-Komplex ist genau dann kompakt, wenn er endlich ist.]{Korollar} % (fold)
\label{sub:116}
Ein \CW-Komplex ist genau dann kompakt, wenn er endlich ist.
% subsection 116 (end)
% section 11 (end)
\newpage

\section{Zelluläre Homologie} % (fold)
\label{sec:12}

\subsection{Definition} % (fold)
\label{sub:121}
Sei $X$ ein \CW-Komplex. Dann heißt 
\[
	C_n^{\cell}(X;R) := H_n \enbrace*{X^{(n)}, X^{(n-1)};R} 
\]
der \bet{$\mathbf{n}$-te zelluläre Kettenmodul}\index{n-ter zellulärer Kettenmodul@$n$-ter zelluläre Kettenmodul} von X mit Koeffizienten in $R$. Die $n$-te Randabbildung
$\partial^{\cell}_n : C_n^{\cell}(X;R) \to C_{n-1}^{\cell}(X;R)$ wird durch die Komposition 
\[
	\begin{tikzcd}[column sep=huge, row sep=large]
		H_n \enbrace*{X^{(n)}, X^{(n-1)};R} \rar["\partial_n^{\enbrace{X^{(n)}, X^{(n-1)}}}"] & H_{n-1} \enbrace*{X^{(n-1)};R} \dar["\enbrace*{j^{(n-1)}}_*"] \\
		& H_{n-1} \enbrace*{X^{(n-1)}, X^{(n-2)};R} 
	\end{tikzcd}
\]
definiert. Dabei ist $\partial_n^{\enbrace{X^{(n)}, X^{(n-1)}}}$ die Randabbildung aus der Paarfolge für $\enbrace*{X^{(n)},X^{(n-1)}}$ und $j^{(n-1)}$ die Inklusion
$\enbrace*{X^{(n-1)},\emptyset} \hookrightarrow \enbrace*{X^{(n-1)}, X^{(n-2)}}$.
% subsection 121 (end)

\subsection{Lemma} % (fold)
\label{sub:122}
\[
	\partial_{n-1}^{\cell} \circ \partial_n^{\cell} =0
\]
\minisec{Beweis}
\[
	\begin{tikzcd}
		H_n \enbrace*{X^{(n)}, X^{(n-1)}} \rar["\partial_n"] & H_{n-1}(X^{(n-1)}) \dar[equal] \rar & \cdots\\
		& H_{n-1}\enbrace*{X^{n-1}} \rar["\enbrace*{j^{(n-1)}}_*"] & H_{n-1} \enbrace*{X^{(n-1)}, X^{(n-2)}} \rar["\partial_{n-1}"] & H_{n-2} \enbrace*{X^{(n-2)}} 
		\rar & \cdots \\ 
		& & & H_{n-2} \enbrace*{X^{(n-2)}}  \rar["\enbrace*{j^{(n-2)}}_*"] & H_{n-2} \enbrace*{X^{(n-2)}, X^{(n-3)}} 
	\end{tikzcd}
\]
Also $\partial^{\cell}_{n-1} \circ \partial^{\cell}_n =0$, da schon $\partial_{n-2} \circ j^{(n-1)}_* =0$. \bewende
% subsection 122 (end)

\subsection{Definition} % (fold)
\label{sub:123}
$\enbrace*{C_*^{\cell}(X;R),\partial_*^{\cell}}$ heißt der \Index{zelluläre Kettenkomplex} von $X$ über $R$. Seine Homologie $H_*^{\cell}(X;R)$ heißt die \Index{zelluläre 
Homologie} von $X$.
% subsection 123 (end)

\subsection{Bemerkung} % (fold)
\label{sub:124}
Der zelluläre Kettenkomplex und die zelluläre Homologie sind Funktoren auf der Kategorie der \CW-Komplexe und zellulären Abbildungen.
% subsection 124 (end)

\subsection{Lemma} % (fold)
\label{sub:125}
Sei $X$ ein \CW-Komplex. Wähle $q_i^{(n)}$ und $Q_i^{(n)}$ für die Zellstruktur auf $X$.
\begin{enumerate}[a)]
	\item 
	\[
		\enbrace*{\coprod_{i \in I^{(n)}} Q_i^{(n)}, \coprod_{i \in I^{(n)}} q_i^{(n)}}_* \colon \enbrace*{\coprod_i D^n, \coprod_i S^{n-1};R}  \longrightarrow 
		\enbrace*{X^{(n)}, X^{(n-1)};R} 
	\]
	induziert einen Isomorphismus in $H_*$.
	\item \[
		H_k \enbrace*{\coprod_i D^n, \coprod_i S^{n-1};R} = \begin{cases}
			R[I^{(n)}], &\text{ falls }k=n\\
			0, & \text{ sonst}
		\end{cases}
	\] 
\end{enumerate}
\minisec{Beweis}
\begin{enumerate}[a)]
	\item Sei $D^n_0 = D^n \setminus \set{0}$, $\frac{1}{2} D^n = \set[\frac{x}{2} ]{x \in D^n}$ und $\frac{1}{2} D^n_0 = \frac{1}{2}D^n \setminus \set{0}$. Wir definieren
	\[
		X^{(n-1)}_{++} := X^{(n)} \setminus \bigcup_{i \in I^{(n)}} Q_i^{(n)}(\set{0}) \enspace\text{ und } \enspace
		X^{(n-1)}_+ := X^{(n)} \setminus \bigcup_{i \in I^{(n)}} Q_i^{(n)}\enbrace*{\frac{1}{2} D^n}
	\]
	Dann sind 
	$X^{(n-1)} \hookrightarrow X_+^{(n-1)} \hookrightarrow X_{++}^{(n-1)}$ Homotopieäquivalenzen. Betrachte:
	\[
		\begin{tikzcd}[sep=large]
			\enbrace*{\coprod D^n, \coprod S^{n-1}} \dar["\enbrace*{\coprod Q_i^{(n)}, \coprod q_i^{(n)}}"] \rar["(1)"] & \enbrace*{\coprod D^n, \coprod D^n_0} \dar &
			\enbrace*{\coprod \frac{1}{2} D^n, \coprod_{i \in I^{(n)}}   \frac{1}{2} D_0^n} \lar["(2)"] \dar["(3)"] \\
			\enbrace*{X^{(n)}, X^{(n-1)}} \rar["(5)"] & \enbrace*{X^{(n)}, X^{(n-1)}_{++}}    
			& \enbrace*{X^{(n)} \setminus X^{(n-1)}_+ , X^{(n-1)}_{++} \setminus X^{(n-1)}_+} \lar["(4)"] 
		\end{tikzcd}
	\]
	$(1), (2), (5)$ sind Homotopieäquivalenzen und induzieren Isomorphismen in Homologie. $(3)$ ist ein Homöomorphismus und induziert einen Isomorphismus in Homologie.
	$(4)$ erfüllt die Vorraussetzungen für den Ausschneideisomorphismus und induziert daher auch einen Isomorphismus in Homologie. Damit folgt a).
	\item Es gilt
	\begin{align*}
		H_k \enbrace*{\coprod_i D^n, \coprod_i S^{n-1};R} \xleftarrow{\cong} \bigoplus_i H_k (D^n,S^{n-1};R) &\xrightarrow{\cong} \bigoplus_i \tilde{H}_{k-1}(S^{n-1};R)\\
		&\cong \bigoplus_i \begin{cases}
			R, &\text{ falls }k=n\\
			0, &\text{ sonst}
		\end{cases} \\
		&\cong \begin{cases}
			R[I^{(n)}], &\text{ falls }k=n\\
			0, &\text{ sonst}
		\end{cases} \bewende
	\end{align*}
\end{enumerate}
% subsection 125 (end)

\subsection{Bemerkung} % (fold)
\label{sub:126}
\missingfigure{Diagramm mit 3 Zeilen}
Die Randabbildung $\partial_n^{\cell}$ wird unter diesen Isomorphismen zu einer $I^{(n-1)} \times I^{(n)}$-Matrix. Für $j \in I^{(n-1)}$ und $k \in I^{(n)}$ ist der 
$(j,k)$-Eintrag dieser Matrix genau der Abbildungsgrad einer Selbstabbildung der $(n-1)$-Sphäre. Diese $\enbrace*{Q_i^{(n-1)}, q_i^{(n-1)}}$ induzieren einen Homöomorphismus
$\nicefrac{\coprod_i D^{n-1}}{\coprod_i S^{n-2}} \xleftarrow{\cong} \nicefrac{X^{(n)}}{X^{(n-1)}}$. Durch Komposition mit der Projektion auf die $j$-te Zelle: 
$\nicefrac{\coprod_i D^{n-1}}{\coprod_i S^{n-2}} \xrightarrow{p_j^{(n-1)}}  \nicefrac{D^{n-1}}{S^{n-2}} \cong S^{n-1}$. Der $(j,k)$-te Eintrag ist nun der 
Abbildungsgrad von 
\[
	\begin{tikzcd}
		S^{n-1} \rar["q^{(n-2)}_k"] & X^{(n-1)} \rar & \nicefrac{X^{(n-1)}}{X^{(n-2)}} \rar["\cong"] & \nicefrac{\coprod_i D^{n-1}}{\coprod_i S^{n-2}} \rar["p_j^{(n-1)}"]
		& S^{n-1}
	\end{tikzcd}
\]
% subsection 126 (end)

\subsection[Definition: Natürlicher Isomorphismus]{Definition} % (fold)
\label{sub:127}
Eine natürliche Transformation $\tau$ zwischen Funktoren $F,G\colon \mathcal{C} \to \mathcal{D}$ heißt ein \Index{natürlicher Isomorphismus}, wenn 
$\tau_C \colon F(C) \to G(C)$ für alle Objekte $C$ von $\mathcal{C}$ ein Isomorphismus in $\mathcal{D}$ ist.
% subsection 127 (end)

\subsection{Satz} % (fold)
\label{sub:128}
Für \CW-Komplexe gibt es einen natürlichen Isomorphismus $\tau_X \colon H_*^{\cell}(X;R) \xrightarrow{\cong} H_*(X;R)$.
\minisec{Beweis (mit Lemma \ref{sub:129})}
Betrachte folgendes Diagramm \marginnote{"$R$" aus Platzgründen weggelassen}
\[
	\begin{tikzcd}[font=\small, row sep=large]
		0=H_n \enbrace*{X^{(n-1)}} \arrow[dr] & & H_n \enbrace*{X^{(n+1)}} \rar["\cong", "\text{Lemma b)}"'] & H_n(X)\\  
		& h_n \enbrace*{X^{(n)}} \arrow[ur, two heads] \\
		H_{n+1} \enbrace*{X^{(n+1)},X^{(n)}} \arrow[rr,"\partial_{n+1}^{\cell}"] & & H_n \enbrace*{X^{(n)},X^{(n-1)}} \arrow[rr,"\partial_n^{\cell}"] & &
		H_{n-1}  \enbrace*{X^{(n-1)},X^{(n-2)}}  
	\end{tikzcd}
\]
\todo{Diagramm vervollständigen}
Es folgt, dass $j_*^{(n)}$ ein Isomorphismus ist, $H_n(X^{(n)};R) \xrightarrow{\cong} \ker \partial_n^{\cell}$ und $\im \partial_{n+1}^{X^{(n+1)}, X^{(n)}} \xrightarrow{\cong} \im \partial_{n+1}^{\cell}$. Weiter erhalten wir einen Isomorphismus
\[
	\begin{tikzcd}
		H_n(X;R) \cong H_n \enbrace*{X^{(n+1)};R} & \nicefrac{H_n(X^{(n)};R)}{ \im \partial_{n+1}^{X^{(n+1)}, X^{(n)}}} \lar["\cong"] 
	\end{tikzcd}
\]
Damit folgt die Behauptung. \bewende
% subsection 128 (end)

\subsection{Lemma} % (fold)
\label{sub:129}
Sei $X$ ein \CW-Komplex.
\begin{enumerate}[a)]
	\item Für $k >n$ ist $H_k(X^{(n)};R)=0$.
	\item Für $k <n$ induziert die Inklusion $l_n : X^{(n)} \hookrightarrow X$. einen Isomorphismus $(l_n)_* : H_k(X^{(n)};R) \to H_k(X;R)$.
\end{enumerate}
\minisec{Beweis}
\begin{enumerate}[a)]
	\item Ist $k>n$, so erhalten wir aus der Paarfolge einen Isomorphismus 
	\[
		\begin{tikzcd}[column sep=small]
			\underset{=0}{H_{k+1}\enbrace*{X^{(n)}, X^{(n-1)};R}} \rar &  H_k(X^{(n-1)};R) \rar["\cong"] & H_k(X^{(n)};R) \rar 
			& H_k \underset{=0}{\enbrace*{X^{(n)}, X^{(n-1)};R} }
		\end{tikzcd}
	\]
	Da $H_k(X^{(0)};R)=0$ für $k>0$ folgt a) nun durch endliche Induktion für $n=0,1, \ldots , k-1$.
	\item Ist $n >k$, so erhalten wir 
	\[
		\begin{tikzcd}[column sep=small]
			\underset{=0}{H_{k+1} \enbrace*{X^{(n+1)},X^{(n)};R}} \rar & H_k \enbrace*{X^{(n)};R} \rar["\cong"] & H_k \enbrace*{X^{(n+1)};R} \rar 
			& \underset{=0}{H_k \enbrace*{X^{(n+1)},X^{(n)};R}   }
		\end{tikzcd}
	\]
	Für $m \ge n >k$ induziert die Inklusion $X^{(n)} \hookrightarrow X^{(m)}$ daher einen Isomorphismus $H_k(X^{(n)};R) \to H_k(X^{(m)};R)$. Ist $\dim X$ endlich, so folgt
	die Behauptung. 
	
	Für den allgemeinen Fall beobachten wir zunächst, dass es zu jeder stetigen Abbildung $\sigma\colon\abs*{\Delta^j}\to X$ ein $n$ gibt mit $\im \sigma \subseteq X^{(n)}$,
	da $\im \sigma \subseteq X$ kompakt ist und daher nur endlich viele offene Zellen von $X$ trifft (\ref{sub:115}).
	\[
		\Longrightarrow C_*(X;R) = \bigcup_n C_*(X^{(n)};R)
	\]
	Damit folgt \begin{enumerate}[(1)]
		\item $\forall x \in H_k(X;R) : \exists n : x \in \im \enbrace*{H_k(X^{(n)};R) \to H_k(X;R)} $
		\item $\forall x\in H_k(X^{(n)};R)$ sodass das Bild von $x$ in $H_k(X;R)$ trivial ist, so $\exists n \ge n$, sodass das Bild von $x$ in $H_k(X^{(m)};R)$ trivial ist.
	\end{enumerate}
	Da für $m \ge n >k$ $H_k(X^{(n)};R) \to H_k(X^{(m)};R)$ ein Isomorphismus ist, ergibt sich damit die Behauptung. \bewende
\end{enumerate}
% subsection 129 (end)

\subsection{Korollar} % (fold)
\label{sub:1210}
Die Eulercharakteristik eines endlichen simplizialen Komplexes ist eine topologische Invariante.
\minisec{Beweis}
Sei $a_n$ die Anzahl der $n$-Simplizes von $K$. Dann gibt es eine \CW-Struktur auf $\abs*{K}$, die für jedes $n$ genau $a_n$ Zellen hat. 

\[
	\chi(K) = \sum_{i=0}^{\infty} (-1)^i a_i = \sum_{i=0}^{\infty} (-1)^i \Rg C_i^{\cell} \enbrace[\big]{\abs*{K}} \stackrel{\text{Kapitel \ref{sec:4}}}{=}
	\sum_{i=0}^{\infty} (-1)^i \Rg H_i^{\cell} \enbrace[\big]{\abs*{K}} = \sum_{i=0}^{\infty} (-1)^i \Rg H_i \enbrace[\big]{\abs*{K_i} }    
\]
Da die $H_i \enbrace[\big]{\abs*{K} } $ eine topologische Invariante sind, ist auch $\chi(K)$ eine topologische Invariante. \bewende
% subsection 1210 (end)
% section 12 (end)
\cleardoubleoddemptypage
\pagenumbering{Alph}
\setcounter{page}{1}
\printindex
\listoffigures
\todototoc
\listoftodos[Todo's und andere Baustellen]
\end{document}